\documentclass[11pt,oneside,openany]{book}

\usepackage[T1]{fontenc}
\usepackage{lmodern}
\usepackage[letterpaper,margin=1.10in,includeheadfoot]{geometry}
\usepackage{amsmath,amssymb,amsthm,mathtools}
\usepackage{microtype}
\usepackage{xcolor}
\usepackage{titlesec}
\usepackage{fancyhdr}
\usepackage{needspace}
\usepackage[titles]{tocloft}
\usepackage{hyperref}

\ifdefined\pdftrailerid
  \pdftrailerid{MathGenome Advanced Homological Algebra}
\fi

\definecolor{ink}{HTML}{172923}
\definecolor{deepgreen}{HTML}{2B5242}
\definecolor{rulegreen}{HTML}{8A9F92}
\hypersetup{
  colorlinks=true,linkcolor=deepgreen,urlcolor=deepgreen,citecolor=deepgreen,
  pdftitle={Advanced Homological Algebra},
  pdfsubject={Spectral Sequences, Derived Categories, and Cohomology},
  pdfauthor={MathGenome},bookmarksnumbered=true
}
\setcounter{tocdepth}{1}
\setcounter{secnumdepth}{2}
\setlength{\cftbeforechapskip}{0.50em}
\setlength{\headheight}{15pt}
\setlength{\emergencystretch}{2em}

\titleformat{\chapter}[display]
  {\normalfont\color{ink}}{\color{deepgreen}\large\scshape Chapter \thechapter}{9pt}
  {\huge\bfseries}
\titlespacing*{\chapter}{0pt}{4pt}{25pt}
\titleformat{\section}{\normalfont\Large\bfseries\color{ink}}{\thesection}{0.7em}{}
\titleformat{\subsection}{\normalfont\large\bfseries\color{ink}}{\thesubsection}{0.7em}{}

\pagestyle{fancy}
\fancyhf{}
\fancyhead[L]{\small\sffamily\color{deepgreen}MathGenome}
\fancyhead[R]{\small\nouppercase{\leftmark}}
\fancyfoot[C]{\small\thepage}
\renewcommand{\headrulewidth}{0.35pt}
\renewcommand{\chaptermark}[1]{\markboth{\thechapter.\ #1}{}}
\fancypagestyle{plain}{%
  \fancyhf{}\fancyfoot[C]{\small\thepage}\renewcommand{\headrulewidth}{0pt}}

\newtheorem{theorem}{Theorem}[chapter]
\newtheorem{lemma}[theorem]{Lemma}
\newtheorem{proposition}[theorem]{Proposition}
\newtheorem{corollary}[theorem]{Corollary}
\theoremstyle{definition}
\newtheorem{definition}[theorem]{Definition}
\newtheorem{example}[theorem]{Example}
\theoremstyle{remark}
\newtheorem{remark}[theorem]{Remark}
\allowdisplaybreaks[1]

\begin{document}
\frontmatter
\hypersetup{pageanchor=false}
\begin{titlepage}
\thispagestyle{empty}
\thispdfpagelabel{Cover}
\begin{flushleft}
\vspace*{0.65in}
{\large\scshape\color{deepgreen}MathGenome\par}
\vspace{0.35in}
{\color{rulegreen}\rule{0.94\textwidth}{0.65pt}\par}
\vspace{0.43in}
{\color{ink}\fontsize{36}{42}\selectfont\bfseries
Advanced\\[4pt] Homological Algebra\par}
\vspace{0.32in}
{\color{deepgreen}\Large Spectral Sequences, Derived Categories,\\[3pt] and Cohomology\par}
\vfill
{\large Mathematical lecture notes\par}
\vspace{0.10in}
{\normalsize MathGenome v0.21.0 dependency baseline\par}
\vspace*{0.45in}
\end{flushleft}
\end{titlepage}
\hypersetup{pageanchor=true}

\begingroup
\small
\tableofcontents
\endgroup

\chapter*{Scope and conventions}
\addcontentsline{toc}{chapter}{Scope and conventions}

Filtered complexes lead to spectral sequences whose pages can be computed
before their abutments are identified. We begin by constructing those pages
and proving the comparison with the filtered cohomology. The two filtrations
of a total complex give a first test of the comparison. The middle chapters
develop localization at quasi-isomorphisms, the standard truncations, and
bounded total derived functors. They also identify higher extension
splicing with shifted composition. These constructions give two spectral
sequences for hypercohomology and a comparison for the derived functor of a
composition. The closing chapters derive group homology and cohomology,
Lie algebra cohomology, and sheaf cohomology results from those tools and
record the hypotheses each application uses.

Unless otherwise stated, complexes are cochain complexes and shifts are
\(K[n]^i=K^{i+n}\), with differential multiplied by \((-1)^n\).
A decreasing filtration is written \(F^pK\), and a cohomological spectral
sequence has differential of bidegree \((r,1-r)\). For commuting
bicomplex differentials the total differential used here is
\(d_h+(-1)^p d_v\) on bidegree \((p,q)\). Homological constructions specify
their own indices and left and right scalar structures. A finite
filtration of an object has both a lowest and a highest effective index,
which may depend on the object or on the cochain degree.

We use ordinary homology and cohomology sequences, comparison of
object resolutions, and the module versions of Ext, Tor, universal
coefficients, and K\"unneth established in the algebra core
\cite[Chapters 10--11]{p103}. We specify where an abstract construction uses
an abelian category, where a concrete calculation uses modules, and where
an argument uses exactness of unions or a bound on the complexes. The
bounded constructions here are prerequisites for later treatments of
unbounded totalizations, derived tensor products and internal Hom, and dg
enhancements. The separate result map records the mathematical sources,
proof dependencies, and candidate Lean declarations at the fixed dependency
baseline.

\mainmatter
\chapter[Filtered complexes and convergence]
  {Filtered Complexes, Exact Couples,\\and Convergence}
\label{ch:filtrations}

A filtered complex gives successive approximations to its cohomology.
The pages and differentials can be constructed in any abelian category.
Identifying their limiting pieces with cohomology requires an additional
argument about the filtration. We give both the finite argument and a
module version allowing infinitely many filtration steps. The indexing
throughout is cohomological, with decreasing filtrations. The two
constructions below follow the same indexing as
\cite[Sections 12.21 and 12.24]{stacks-homology}. The filtered-complex
and exact-couple arguments can also be found in
\cite[Sections 5.4, 5.5, and 5.9]{weibel}.

\section{The pages of a filtered complex}

\begin{definition}[Filtered complexes and associated graded objects]
\label{res:F01}
Let $\mathcal A$ be an abelian category. A decreasingly filtered cochain
complex $(K,F)$ consists of subobjects $F^pK^n\subset K^n$, for integers
$p,n$, satisfying
\[
 F^{p+1}K^n\subset F^pK^n,\qquad
 d(F^pK^n)\subset F^pK^{n+1}.
\]
A filtered map is a cochain map carrying each $F^p$ into the corresponding
$F^p$ of its target. Each $F^pK$ is itself a cochain complex, and
$\operatorname{gr}_F^pK=F^pK/F^{p+1}K$ is a quotient complex. Its differential
is induced by the restriction of $d$ to $F^pK$.
\end{definition}

The restriction squares to zero because the inclusion of the subcomplex
is monic. The quotient differential exists because $d$ preserves the
next filtration step. Its square is zero because the quotient projection
is epic. Both constructions respect identity maps and composition.

\begin{lemma}[Exact sequences between filtration steps]
\label{res:F02}
The filtration steps give two natural short exact sequences of complexes.
\begin{align*}
 0&\longrightarrow F^{p+1}K\longrightarrow F^pK
    \longrightarrow \operatorname{gr}_F^pK\longrightarrow 0,\\
 0&\longrightarrow \operatorname{gr}_F^{p+1}K
    \longrightarrow F^pK/F^{p+2}K
    \longrightarrow \operatorname{gr}_F^pK\longrightarrow 0 .
\end{align*}
In the second sequence the connecting morphism sends a cycle in the
quotient, represented by $x$ in $F^pK^n$, to the class represented by
$dx$ in $F^{p+1}K^{n+1}/F^{p+2}K^{n+1}$.
\end{lemma}

\begin{proof}
The first sequence is exact by the definition of the quotient. In the
second, the kernel of $F^pK/F^{p+2}K\to F^pK/F^{p+1}K$ is
$F^{p+1}K/F^{p+2}K$. These statements hold in each degree, and the
restriction and quotient differentials commute with the displayed maps
by Definition~\ref{res:F01}. To compute the connecting morphism, take the
differential of a lift of a quotient cycle. It lies in the kernel
subcomplex and is a cycle there. Changing the lift by a kernel element
changes the resulting cycle by a kernel boundary. Changing the quotient
cycle by a boundary allows one to subtract the differential of a lift,
which does not alter the image under the connecting morphism. Every
operation occurs in the natural diagram of short exact sequences, so
this description also proves naturality. In an abelian category,
the lifts in this argument are taken after epic pullback, as described
below.
\end{proof}

\paragraph{Working with subobjects.}
We use intersections, inverse images, images, and finite sums of
subobjects. A formula involving a representative means a morphism
from a test object. An epic map need not lift that morphism on the
original test object, but does after pulling back the epimorphism.
Epimorphisms in an abelian category are stable under pullback. Two
successive lifts can therefore be made over a common epic cover.
Containment in a subobject descends along that cover: if $T'\to T$
is epic and the composite $T'\to T\to X$ factors through $Y\subset X$,
then $T\to X\to X/Y$ vanishes and $T\to X$ factors through $Y$.
Representatives of a sum of subobjects can be decomposed after an
epic pullback along the map from their direct sum onto their sum.
These rules justify the representative arguments below without
countable sums or projective generators in $\mathcal A$.

\begin{proposition}[Cycles, boundaries, and page differentials]
\label{res:F03}
For $r\geq 0$ and $n=p+q$, put
\[
 A_r^{p,n}=F^pK^n\cap d^{-1}(F^{p+r}K^{n+1}).
\]
On the graded piece $\operatorname{gr}_F^pK^n$, put
\begin{align*}
 Z_r^{p,q}
   &=\bigl(A_r^{p,n}+F^{p+1}K^n\bigr)/F^{p+1}K^n,\\
 B_r^{p,q}
   &=\bigl((F^pK^n\cap d(F^{p-r+1}K^{n-1}))+F^{p+1}K^n\bigr)
      /F^{p+1}K^n.
\end{align*}
Then $B_r^{p,q}\subset Z_r^{p,q}$. Put $E_r^{p,q}=Z_r^{p,q}/B_r^{p,q}$.
There is a well-defined map
\[
 d_r^{p,q}\colon E_r^{p,q}\longrightarrow E_r^{p+r,q-r+1},
 \qquad [x]\longmapsto [dx],
\]
and $d_r^2=0$.
\end{proposition}

\begin{proof}
Every boundary in $F^pK^n$ lies in $A_r^{p,n}$ since its differential
vanishes. For $r=0$, $A_0^{p,n}=F^pK^n$ and $B_0=0$. Thus $E_0$ is
the graded piece and $d_0$ is its differential. For $r\geq1$ the
subobject identity
\[
 F^pK^n\cap d(F^{p-r+1}K^{n-1})=d A_{r-1}^{p-r+1,n-1}
\]
holds by pulling back an epic map to the image of $d$.
Also $A_r^{p,n}\cap F^{p+1}K^n=A_{r-1}^{p+1,n}$. Since the boundary
subobject is contained in $A_r$, the modular law for subobjects
gives the quotient presentation
\begin{equation}
 \label{eq:page-quotient}
 E_r^{p,q}=
 \frac{A_r^{p,n}}
 {A_{r-1}^{p+1,n}+d A_{r-1}^{p-r+1,n-1}}\qquad(r\geq1).
\end{equation}
For $x\in A_r^{p,n}$, its differential lies in $F^{p+r}K^{n+1}$
and has zero differential, so it lies in the numerator at the
target. Adding a term from $A_{r-1}^{p+1,n}$ changes $dx$ by an
element of $dA_{r-1}^{p+1,n}$ in the target denominator. Adding a
boundary does not change $dx$. Thus the formula descends. Applying
it twice gives zero.
\end{proof}

\begin{theorem}[The first pages and the passage to the next page]
\label{res:F04}
The pages in Proposition~\ref{res:F03} form a cohomological
spectral sequence. More explicitly,
\[
 E_0^{p,q}=\operatorname{gr}_F^pK^{p+q},\quad
 E_1^{p,q}=H^{p+q}(\operatorname{gr}_F^pK),\quad
 E_{r+1}^{p,q}\simeq H^{p,q}(E_r,d_r),
\]
where the last object is the kernel of the outgoing differential
modulo the image of the incoming differential. Under the $E_1$
identification, $d_1$ is the connecting morphism of the second exact
sequence in Lemma~\ref{res:F02}.
\end{theorem}

\begin{proof}
The assertion for $E_0$ follows from Proposition~\ref{res:F03}. Its
cycles are exactly the classes whose differential lands in $F^{p+1}$,
and its boundaries come from $d(F^pK^{p+q-1})$, which gives $E_1$.
The connecting morphism in Lemma~\ref{res:F02} had the formula
$[x]\mapsto[dx]$, giving the assertion about $d_1$ with the chosen signs.

For $r\geq1$, let $n=p+q$ and let $x$ represent a cycle of $(E_r,d_r)$.
By \eqref{eq:page-quotient} at the target, write
\[
 dx=u+dv,\qquad
 u\in A_{r-1}^{p+r+1,n+1},\quad v\in A_{r-1}^{p+1,n}.
\]
Then $x-v$ represents the same class, and $d(x-v)=u\in F^{p+r+1}$.
Thus the $d_r$-cycles are represented by $A_{r+1}^{p,n}$. The converse
holds because the differential of an element of $A_{r+1}^{p,n}$ is
closed and lies in $F^{p+r+1}$, hence belongs to the first summand of
the target denominator.

The incoming $d_r$ contributes $d A_r^{p-r,n-1}$. This also contains
$d A_{r-1}^{p-r+1,n-1}$, since the latter antecedents satisfy
$d y\in F^p$. The only remaining quotient relation from
\eqref{eq:page-quotient} comes from
$A_{r+1}^{p,n}\cap A_{r-1}^{p+1,n}=A_r^{p+1,n}$. Consequently the
cohomology of $(E_r,d_r)$ is canonically the quotient of
$A_{r+1}^{p,n}$ by
$A_r^{p+1,n}+d A_r^{p-r,n-1}$, which is $E_{r+1}^{p,q}$.
The epic pullback argument above supplies the indicated decompositions
and proves the equalities of subobjects for arbitrary $\mathcal A$.
\end{proof}

An application starting with the second page can use the subsystem
$r\geq2$ with the same bidegrees and page transitions. Its first
specified page is then $E_2$ in our notation. The identification
of either earlier page still requires the filtered complex.

\section{Exact couples and functoriality}

\begin{theorem}[The exact couple of a filtration]
\label{res:F05}
Put $D_1^{p,n}=H^n(F^pK)$ and regard $E_1^{p,n-p}$ as indexed by
$(p,n)$. The first short exact sequence of Lemma~\ref{res:F02}
gives an exact couple, with maps
\[
\begin{split}
 i&\colon D_1^{p+1,n}\longrightarrow D_1^{p,n},\\
 j&\colon D_1^{p,n}\longrightarrow E_1^{p,n-p},\\
 k&\colon E_1^{p,n-p}\longrightarrow D_1^{p+1,n+1}.
\end{split}
\]
The derived exact couples give the same $E_r$ and $d_r$ as
Proposition~\ref{res:F03} for all $r\geq1$.
\end{theorem}

\begin{proof}
The long exact cohomology sequence supplies exactness at each of the
three types of terms. To derive an exact couple, temporarily suppress
the indices. Set $D'=\operatorname{im}i$ and
$E'=\ker(jk)/\operatorname{im}(jk)$. The differential $jk$ squares to
zero because $kj=0$. The new maps are
\[
 i'(a')=i(a'),\qquad j'(ia)=[j(a)],\qquad k'([e])=k(e).
\]
Changing $a$ by $\ker i=\operatorname{im}k$ changes $j(a)$ by a
boundary. If $jk(e)=0$, then $k(e)\in\ker j=\operatorname{im}i$,
and $k$ kills boundaries because $kj=0$. Thus the three maps exist.
For exactness at the target of $i'$, if $j'(ia)=0$ then
$j(a)=jk(e)$ for a representative $e$. Hence $a-k(e)\in\ker j
=\operatorname{im}i$, so $ia\in\operatorname{im}i'$. At its source,
if $i(a')=0$ then $a'=k(e)$, and $a'\in\operatorname{im}i$ implies
$jk(e)=0$. Exactness at $E'$ follows because $k(e)=0$ implies
$e=j(a)$ and $[e]=j'(ia)$. The other inclusions follow by
composition. The same arguments apply componentwise to a bigraded
couple and with epic lifts in $\mathcal A$.

After $r-1$ derivations the terms and the indices of the maps are
\[
\begin{split}
 D_r^{p,n}
  &=\operatorname{im}\bigl(H^n(F^{p+r-1}K)\longrightarrow
        H^n(F^pK)\bigr),\\
 j_r&\colon D_r^{p,n}\longrightarrow E_r^{p+r-1,n-p-r+1},\\
 k_r&\colon E_r^{p,n-p}\longrightarrow D_r^{p+1,n+1}.
\end{split}
\]
To compare with \eqref{eq:page-quotient}, $j_r$ takes a closed
representative in $F^{p+r-1}$ of an element of $D_r^{p,n}$ and uses
that representative in $E_r$. Two such representatives differ by
a boundary from $F^p$ lying in $F^{p+r-1}$, hence give the same
page class. The map $k_r$ sends $[x]$ to the cohomology class of
$dx$ in $F^{p+1}$, which lies in the image from $F^{p+r}$.
The denominator in \eqref{eq:page-quotient} changes $dx$ by a
boundary in $F^{p+1}$. For $r=1$ these are exactly the maps of the
long exact sequence of the first short exact sequence of
Lemma~\ref{res:F02}. On derivation the new $j$ uses a lift from
one deeper filtration step. The new $k$ uses a representative
$x\in A_{r+1}^{p,n}$, which exists by Theorem~\ref{res:F04}.
These are the stated formulas for $r+1$. Finally $j_rk_r$
sends $[x]$ to $[dx]$ in bidegree $(p+r,q-r+1)$, proving
agreement of the constructions.
\end{proof}

\begin{proposition}[Functoriality and comparison of pages]
\label{res:F06}
A filtered cochain map induces bidegree-zero maps on all pages,
commuting with the page differentials and the isomorphisms
$E_{r+1}\simeq H(E_r,d_r)$. These maps preserve identities and
composition. An isomorphism in every bidegree of one page induces
isomorphisms on all later pages.
\end{proposition}

\begin{proof}
A filtered cochain map carries $A_r$ into $A_r$ because it commutes
with $d$. The same holds for the two summands in the denominator of
\eqref{eq:page-quotient}. The induced quotient map sends $[x]$ to
the class of its image. Consequently it commutes with $[x]\mapsto[dx]$.
The next-page identification in Theorem~\ref{res:F04} was induced
by inclusion of $A_{r+1}$, so it is natural as well. Maps on quotients
inherit identities and composition from the maps on complexes.
An isomorphism of the bigraded objects commuting with the differentials
is an isomorphism of their complexes and hence induces an isomorphism
on the next page. Induction gives the assertion for every later page.
\end{proof}

\section{The limiting filtration}

\begin{theorem}[Convergence under degreewise finite filtrations]
\label{res:S01}
Suppose that for each $n$ there are integers $a_n,b_n$ such that
$F^pK^n=K^n$ for $p\leq a_n$ and $F^pK^n=0$ for $p\geq b_n$.
The page terms are eventually stable in every bidegree and, for
each fixed total degree, the stabilization is uniform over its
bidegrees. Define
\[
 F^pH^n(K)=\operatorname{im}\bigl(H^n(F^pK)\longrightarrow H^n(K)\bigr).
\]
This is a finite decreasing filtration. The stable page is
naturally identified with its associated graded by
\[
 E_\infty^{p,q}\simeq
 F^pH^{p+q}(K)/F^{p+1}H^{p+q}(K).
\]
\end{theorem}

\begin{proof}
Write $Z^n=\ker(d\colon K^n\to K^{n+1})$ and
$B^n=\operatorname{im}(d\colon K^{n-1}\to K^n)$.
Choose $r\geq1$ large enough that $F^{p+r}K^{n+1}=0$ and
$F^{p-r+1}K^{n-1}=K^{n-1}$. Then the numerator of
\eqref{eq:page-quotient} is $F^pK^n\cap Z^n$, its first denominator
summand is $F^{p+1}K^n\cap Z^n$, and the second is
$F^pK^n\cap B^n$. Increasing $r$ leaves this presentation
unchanged. There are only finitely many possibly nonzero terms
of $E_0$ on each total degree, so a common bound for these
bidegrees suffices. All other terms remain zero.

The image defining the cohomology filtration is the image of
$F^pK^n\cap Z^n$ in $Z^n/B^n$. This proves that the filtration
has the same two bounds in degree $n$ as the original
filtration in that degree. The kernel of
$F^pK^n\cap Z^n\longrightarrow F^pH^n/F^{p+1}H^n$ is
$(F^{p+1}K^n\cap Z^n)+(F^pK^n\cap B^n)$. Indeed, a
representative mapping into $F^{p+1}H^n$ agrees in cohomology
with a cycle in $F^{p+1}K^n$, so their difference is a boundary
already in $F^pK^n$. Conversely both summands map to zero.
This identifies the quotient with the stable page. The construction
uses only natural subobject operations, so the identification
commutes with filtered cochain maps.
\end{proof}

Thus in the finite situation the spectral sequence determines the
successive quotients of an actual finite filtration of the target.
The filtration is exhaustive, separated, and complete. Different
conventions for convergence beyond this situation appear in
\cite[Chapter 3]{mccleary} and \cite[Section 5.2]{weibel}.
We will name the filtration properties that a limit claim provides.

\begin{theorem}[Bounded below and exhaustive chain filtrations]
\label{res:S02}
Let $C_\bullet$ be a chain complex of modules with an increasing
filtration preserved by the differential. Assume that in each degree
$n$, $F_aC_n=0$ for $a$ sufficiently small and
$\bigcup_a F_aC_n=C_n$. Define the homological page indices so that
$E^r_{a,b}$ has total chain degree $a+b$. For each bidegree there is
an eventual numerator of cycles in the associated graded, while
the boundaries form an increasing family. The quotient of
that numerator by the union of those boundaries is
$E^\infty_{a,b}\simeq F_aH_{a+b}(C)/F_{a-1}H_{a+b}(C)$.
This identification is natural. The target filtration is bounded
below in each degree and exhaustive.
\end{theorem}

\begin{proof}
Apply Proposition~\ref{res:F03} to $K^n=C_{-n}$ and
$F^pK^n=F_{-p}C_{-n}$. Then $(p,q)=(-a,-b)$, and the
differential has homological bidegree $(-r,r-1)$.
In the decreasing cochain indexing $F^{p+r}K^{n+1}=0$ for large
$r$ by the lower bound on the increasing chain filtration.
The first formula for $Z_r$ in Proposition~\ref{res:F03}
therefore becomes
\[
 \bigl((F^pK^n\cap Z^n)+F^{p+1}K^n\bigr)/F^{p+1}K^n
\]
for all sufficiently large $r$. On the boundary side,
exhaustivity gives
$\bigcup_r (F^pK^n\cap d(F^{p-r+1}K^{n-1}))=F^pK^n\cap B^n$.
Indeed a boundary has a preimage in a finite filtration step,
and every element in the union is already a boundary.
Forming the quotient after adding $F^{p+1}K^n$ commutes with
this directed union, since membership means that a representative
is the sum of an element in some step and an element in
$F^{p+1}K^n$. Quotienting the eventual cycle subobject of the
graded piece by this union gives the quotient
$(F^pK^n\cap Z^n)/((F^{p+1}K^n\cap Z^n)+(F^pK^n\cap B^n))$
by the modular law. It is the graded piece of $H^n(K)$ by the
image and kernel computation in Theorem~\ref{res:S01}.
Every cohomology class has a cycle in some step by exhaustivity,
and the zero bounds descend to the target. All constructions
preserve maps of filtered complexes. Reversing the indices gives
the assertion for $C$.
\end{proof}

The boundary family in Theorem~\ref{res:S02} need not stabilize
at a finite page. For example, in the decreasing cochain
convention over a field, put $K^0=\bigoplus_{i\geq1}k e_i$ and
$K^1=\bigoplus_{i\geq1}k f_i$, with $d e_i=f_i$. Give $e_i$
filtration degree $-i$ and every $f_i$ degree $0$, where
filtration degree $s$ means inclusion in $F^p$ exactly for
$p\leq s$. In bidegree $(0,1)$, $E_r$ for $r\geq1$ is
spanned by the classes of $f_i$ with $i\geq r$.
Its eventual quotient by the union of boundaries is zero,
as is $H^1(K)$, although no single page has this value.

\begin{corollary}[Comparison with the target]
\label{res:S03}
Let $u\colon(K,F)\to(L,F)$ be a filtered cochain map of two
complexes satisfying the finiteness hypotheses of
Theorem~\ref{res:S01}. If $u$ gives an isomorphism in every
bidegree on some page, then $H^n(u)$ is an isomorphism on every
target filtration step, and its inverse preserves the filtrations.
\end{corollary}

\begin{proof}
By Proposition~\ref{res:F06}, the maps are isomorphisms on every
later page. Theorem~\ref{res:S01} identifies the resulting maps
with isomorphisms on all associated graded pieces of $H^n(u)$.
Fix $n$ and take a common $b_n$ for the two finite filtrations.
At $p=b_n$ the map is an isomorphism between zero objects.
The short exact sequence for a step and its successor, on the
source and target, gives a commutative diagram. If the map at
$p+1$ is an isomorphism, the left map and the map on the
$p$th graded quotient are isomorphisms. The snake lemma gives
zero kernel and cokernel for the middle map at $p$. Downward
induction to a common lower bound proves the statement at every
step. These stepwise inverses are compatible with inclusions,
which proves the filtration assertion about the inverse.
\end{proof}

\section{Collapse, edge maps, and products}

\begin{proposition}[Collapse and extensions]
\label{res:S04}
Under the hypotheses of Theorem~\ref{res:S01}, if a page $E_s$ is
concentrated in one row $q=q_0$ with $s\geq2$, or in one column
$p=p_0$ with $s\geq1$, it is stable. Each target cohomology object
is naturally isomorphic to its unique possibly nonzero graded
piece, via the relevant quotient and inclusion. If several pieces
survive on a diagonal, they determine a succession of extension
problems.
\end{proposition}

\begin{proof}
Every differential on or after the given page changes the row
by $1-r\ne0$ in the first case, or the column by $r\ne0$ in the
second. The exactness of the passage to the next page therefore
gives $E_s=E_\infty$. On each total degree the finite filtration
has at most one nonzero quotient. The steps before that quotient
equal the whole object and the steps after it are zero. The
displayed identification follows from Theorem~\ref{res:S01}.
If there are several surviving pieces, each next step is
specified by an exact sequence
$0\to F^{p+1}H^n\to F^pH^n\to E_\infty^{p,n-p}\to0$.
The pieces alone need not specify its extension class.
For example, two-step filtrations of $\mathbb Z/4$ by
$2(\mathbb Z/4)$ and of $\mathbb Z/2\oplus\mathbb Z/2$ by its
first factor have the same graded pieces. The first extension does
not split, since every element outside the subgroup has order four.
The second is split by its second factor.
\end{proof}

\begin{theorem}[First-quadrant edge maps and five terms]
\label{res:S05}
Suppose that Theorem~\ref{res:S01} applies and that $E_2$ vanishes
outside $p,q\geq0$. For $n\geq0$ there are natural edge maps
$E_2^{n,0}\to H^n(K)$ and $H^n(K)\to E_2^{0,n}$.
In low degrees they fit into an exact sequence
\[
 0\longrightarrow E_2^{1,0}\longrightarrow H^1(K)
 \longrightarrow E_2^{0,1}\xrightarrow{d_2}
 E_2^{2,0}\longrightarrow H^2(K).
\]
\end{theorem}

\begin{proof}
For $r\geq2$, no differential leaves row zero and no differential
enters column zero. Thus the stable terms in row zero are quotients
of the corresponding $E_2$ terms, and those in column zero are
subobjects. In total degree $n$, the only possibly nonzero
stable terms have $0\leq p\leq n$, so $F^0H^n=H^n$ and
$F^{n+1}H^n=0$. Composing the quotient onto $E_\infty^{n,0}$
with the identification with $F^nH^n$ and its inclusion defines
the row edge. Composing projection from $H^n$ onto $F^0H^n/F^1H^n$
with the identification with $E_\infty^{0,n}$ and its inclusion
defines the column edge.

An inspection of possible source and target bidegrees gives
\[
 E_\infty^{1,0}=E_2^{1,0},\qquad
 E_\infty^{0,1}=\ker(d_2^{0,1}),\qquad
 E_\infty^{2,0}=\operatorname{coker}(d_2^{0,1}).
\]
In fact the only outgoing differential from $(0,1)$ for $r\geq2$
is at $r=2$, and the only possible incoming differential to
$(2,0)$ is the same map. The filtration of $H^1$ gives the
injection of $F^1H^1$ and the quotient by it. That quotient
equals the kernel of $d_2$. The last edge map factors as the
quotient by the image of $d_2$ followed by the inclusion of
$F^2H^2$, proving exactness at $E_2^{2,0}$.

For a nonzero $d_2$ over a field $k$, take
$K^1=k a$, $K^2=k b$, and $d a=b$, with filtration degrees
$0$ for $a$ and $2$ for $b$. On $E_0$ the only nonzero positions are
$(0,1)$ and $(2,0)$. The maps $d_0$ and $d_1$ vanish for degree
reasons, so the same classes span $E_2$. The formula of
Proposition~\ref{res:F03} gives $d_2[a]=[b]$, an isomorphism.
The page $E_3$ is zero, in agreement with the acyclicity of $K$.
\end{proof}

\begin{proposition}[Multiplicative filtrations]
\label{res:S06}
Suppose $K$ is a cochain differential graded algebra over a commutative
ring, with $F^pK\cdot F^sK\subset F^{p+s}K$.
Every page carries a product adding bidegrees. The sign in the
derivation identity
\[
 d_r([x][y])=d_r([x])[y]+(-1)^{|x|}[x]d_r([y])
\]
depends on the total degree of $x$.
The product on $E_{r+1}$ agrees with the cohomology product for
$d_r$. Under the hypotheses of Theorem~\ref{res:S01}, its stable
product is the associated graded product of the cohomology algebra.
If the algebra is unital and $1\in F^0K^0$, the class of $1$
gives the unit on every page.
\end{proposition}

\begin{proof}
The product on $E_0$ is well defined by the multiplicativity of
the filtration. For $r\geq1$, let $x\in A_r^{p,n}$ and
$y\in A_r^{s,m}$ be homogeneous and put $P=p+s$. The identity
$d(xy)=(dx)y+(-1)^n x\,dy$ places the differential in
$F^{P+r}$, so $xy\in A_r^{P,n+m}$. If $x$ is changed by an
element $v\in A_{r-1}^{p+1,n}$, then $v y\in F^{P+1}$ and
$d(vy)\in F^{P+r}$, giving the first denominator summand of
\eqref{eq:page-quotient}. If $x$ is changed by $du$ with
$u\in A_{r-1}^{p-r+1,n-1}$, write
\[
 (du)y=d(uy)-(-1)^{n-1}u(dy).
\]
Here $uy$ lies in $F^{P-r+1}$ and $d(uy)$ lies in $F^P$,
so $d(uy)$ belongs to the second denominator summand.
Also $u(dy)$ lies in $F^{P+1}$ and its differential lies in
$F^{P+r}$, so it belongs to the first. If the change is made in
$y$, the same containment for $xv$ is immediate. For a change
$du$ in $y$, use $x(du)=(-1)^n(d(xu)-(dx)u)$, placing the
first term in the second denominator summand and the second
term in the first. This verifies well-definedness without any
commutativity assumption on $K$.

Associativity, the unit assertion, and the displayed derivation
identity descend from their representatives. The next-page
isomorphism was induced by the inclusion of $A_{r+1}$, and
that inclusion preserves products. Finally, the image filtration
on $H^*(K)$ is multiplicative because products of cycles in
$F^p$ and $F^s$ are cycles in $F^{p+s}$. The stable
identification in Theorem~\ref{res:S01} sends a closed
representative to its cohomology class modulo the next step,
so it preserves the graded product.
\end{proof}

\begin{example}[Stable pages and a nonseparated target]
\label{res:S07}
Let $k$ be a field. In degrees zero and one take
$K^0=\bigoplus_{i\geq0} k e_i$ and $K^1=\bigoplus_{i\geq0} k f_i$.
Define $d e_i=f_i-f_{i+1}$. In each degree, filter by the tails
$i\geq\max(p,0)$. This decreasing filtration is exhaustive and
separated on the complex. The associated graded differential is
an isomorphism at every nonzero graded step, so $E_1=0$ and
every later page is zero. Nevertheless $H^1(K)\simeq k$ and
$F^pH^1(K)=H^1(K)$ for every $p$. In particular the stable
page is compatible with the zero associated graded of the target,
but the target filtration fails separation.
\end{example}

To see the cohomology assertion, let $\sigma\colon K^1\to k$
sum the coefficients. It kills $d(K^0)$. For a finitely supported
sequence $(b_i)$ whose total sum is zero, put
$a_i=\sum_{j=0}^i b_j$. These $a_i$ are finitely supported and
$d(\sum_i a_i e_i)=\sum_i b_i f_i$, so $\ker\sigma=d(K^0)$.
The same recurrence for a zero differential shows $H^0(K)=0$.
The class of $f_i$ generates $H^1(K)$ for every $i$, so every
tail maps onto it. The completeness comparison has zero codomain
because $\varprojlim_p H^1(K)/F^pH^1(K)=0$. On replacing each direct sum
by its tail completion, which is the corresponding product, the
same recurrence solves $d a=b$ for arbitrary $(b_i)$.
The completed complex is acyclic and has the same $E_0$ page.
This verifies directly why the finite hypotheses in
Theorem~\ref{res:S01} matter and why the pages alone cannot detect
whether passing to the completion changes cohomology. General
criteria using complete filtrations require separate hypotheses
and are treated in \cite[Section 5.5]{weibel}.

\chapter{Double Complexes and Tensor Calculations}
\label{ch:double-complexes}

A double complex has two filtrations on its total complex. The two spectral
sequences compute the same total cohomology under a common finiteness
hypothesis, but the two filtrations on that cohomology need not coincide.
We use cohomological indices for the general constructions and switch to
homological indices for the tensor calculations. The sign and indexing
conventions below specify how to pass between the two.

\section{Totalization and its filtrations}

\begin{proposition}[Totalization and interchange of the axes]
\label{res:B01}
Let \(C^{p,q}\), \(p,q\in\mathbb Z\), be objects of an additive category
with differentials \(d_h\colon C^{p,q}\to C^{p+1,q}\) and
\(d_v\colon C^{p,q}\to C^{p,q+1}\). Suppose that
\(d_h^2=d_v^2=0\) and that the horizontal and vertical squares commute.
On \(C^{p,q}\) define
\[
    D=d_h+(-1)^p d_v .
\]
This gives a differential on the direct-sum totalization when the requisite
diagonal coproducts exist, and on the product totalization when the
diagonal products exist. Multiplying \(d_v\) by \((-1)^p\) changes to the
anticommuting convention and leaves \(D\) unchanged. The transposed double
complex \(C^{\mathrm t}\) has \((C^{\mathrm t})^{q,p}=C^{p,q}\) with its
horizontal and vertical differentials interchanged. Multiplication by
\((-1)^{pq}\) from \(C^{p,q}\) to \((C^{\mathrm t})^{q,p}\) induces a
natural isomorphism of total complexes in either totalization.
\end{proposition}

\begin{proof}
Write \(a_{p,q}\) for the injection of a summand and \(b_{p,q}\) for the
projection to a factor, with the total degree specified by \(p+q\).
For the direct sum, the outgoing differential of each summand is the
finite sum
\[
  D a_{p,q}=a_{p+1,q}d_h+(-1)^p a_{p,q+1}d_v.
\]
The coproduct property gives a unique map with these formulas. For
the product, the incoming map to each output factor is
\[
  b_{p,q}D=d_h b_{p-1,q}+(-1)^p d_v b_{p,q-1}.
\]
Only two input coordinates occur, so the product property again gives
a map. Composing twice and collecting the components gives the two
square terms \(d_h^2,d_v^2\) and the mixed terms
\[
     (-1)^{p+1}d_vd_h+(-1)^p d_hd_v
\]
on input bidegree \((p,q)\). They all vanish. The same finite component
calculation applies to the product by projecting to an output factor.
With \(\widetilde d_v=(-1)^p d_v\), the mixed terms express
\(d_h\widetilde d_v+\widetilde d_vd_h=0\). Rescaling by \((-1)^p\)
again recovers the commuting convention.

On the transposed totalization the differential, in the original
coordinates, is \(d_v+(-1)^q d_h\). If \(T\) multiplies the input
\(C^{p,q}\) by \((-1)^{pq}\), the coefficients of \(d_h\) and \(d_v\)
in \(TD\) are respectively
\((-1)^{(p+1)q}\) and \((-1)^p(-1)^{p(q+1)}\).
These agree with the coefficients of \(d_h\) and \(d_v\) in
\(D^{\mathrm t}T\). The formula defines an inverse by the same
multipliers and reindexing in the opposite direction. These formulas
commute with maps preserving both differential directions.
\end{proof}

The signs specialize to the tensor total complex: when both differentials
lower degree, the same formula \(D=d_h+(-1)^p d_v\) works, since replacing
\(p+1\) by \(p-1\) does not change the parity calculation.
See \cite[Sections 5.1 and 5.6]{weibel} and
\cite[Chapter 12, Sections 18 and 25]{stacks-homology}.

\begin{lemma}[Finite diagonals and the column and row filtrations]
\label{res:B02}
Let \(C\) be a double cochain complex in an abelian category and suppose
that for each \(n\) the set
\[
    I_n=\{(p,q)\colon p+q=n,\ C^{p,q}\ne 0\}
\]
is finite. The canonical map from the diagonal direct sum to the diagonal
product is an isomorphism of complexes. Write \(K=\operatorname{Tot}C\)
for their common value. There are two natural decreasing filtrations by
subcomplexes
\begin{align*}
    F_{\mathrm c}^{a} K^n
       &=\bigoplus_{\substack{p+q=n\\p\ge a}} C^{p,q},\\
    F_{\mathrm r}^{a} K^n
       &=\bigoplus_{\substack{p+q=n\\q\ge a}} C^{p,q}.
\end{align*}
Each filtration is finite in every degree. In particular, these statements
apply when \(C^{p,q}=0\) unless \(p\ge p_0\) and \(q\ge q_0\) for fixed
integers \(p_0,q_0\).
\end{lemma}

\begin{proof}
The nonzero factors are indexed by \(I_n\), and the remaining zero factors
may be removed by the universal properties of products and coproducts.
The diagonal comparison is the biproduct isomorphism whose \((i,j)\)-entry
is the identity for \(i=j\) and zero otherwise. The component formulas in
Proposition~\ref{res:B01} show that these comparisons intertwine the
differentials. In the column filtration, \(d_h\) increases the column
and \(d_v\) preserves it, so both carry the indicated summands into the
corresponding subobject in degree \(n+1\). In the row filtration, \(d_v\)
increases the row and \(d_h\) preserves it. Each resulting inclusion is
the split inclusion of a finite sub-biproduct in a given degree; the
splittings need not be maps of complexes.

For a nonzero diagonal, a threshold at most the minimum relevant index
gives the whole degree, and a threshold greater than the maximum gives
zero. The zero diagonal satisfies both bounds. For the translated
first quadrant, a nonzero summand in degree \(n\) must have
\(p_0\le p\le n-q_0\). This gives the required finiteness and allows
the bounds to depend on \(n\).
\end{proof}

Outside the finite-diagonal setting the two totalizations can have
different cohomology. For example, take \(C^{p,-p}=k\) for every integer
\(p\), with \(k\) a nonzero field, and all other terms and differentials
zero. The diagonal comparison on \(H^0\) is the inclusion of finite
support sequences in all sequences. The convergence statements below
use exactly the finite-diagonal hypothesis unless another bound is
explicitly given.

\section{The two spectral sequences}

Use \(H_v^q(C^{p,*})\) for the cohomology of the vertical complex in
column \(p\). The horizontal maps are maps of these unsigned vertical
complexes and induce a differential on \(p\mapsto H_v^q(C^{p,*})\).
The notation \(H_h^p H_v^q(C)\) refers to its degree \(p\) cohomology.
The notation for the opposite order is defined by transposing the axes.

\begin{proposition}[The first pages and their signs]
\label{res:B03}
Suppose \(C\) satisfies Lemma~\ref{res:B02}. For the column filtration
on \(\operatorname{Tot}C\) the first pages, in cohomological convention,
are
\begin{align*}
    E_{0,\mathrm c}^{p,q}&=C^{p,q},
        &d_{0,\mathrm c}^{p,q}&=(-1)^p d_v,\\
    E_{1,\mathrm c}^{p,q}&=H_v^q(C^{p,*}),
        &d_{1,\mathrm c}^{p,q}&=H_v^q(d_h),\\
    E_{2,\mathrm c}^{p,q}&=H_h^p H_v^q(C).
\end{align*}
For the row filtration use \(p\) for filtration degree, so that it is
the original row index. Through the signed interchange of
Proposition~\ref{res:B01} the row pages identify with the column pages
of \(C^{\mathrm t}\). In particular
\[
    E_{1,\mathrm r}^{p,q}\simeq H_h^q(C^{*,p}),
    \qquad
    E_{2,\mathrm r}^{p,q}\simeq H_v^p H_h^q(C).
\]
These identifications are natural for maps of double complexes.
\end{proposition}

\begin{proof}
In total degree \(p+q\) the quotient of the column terms at thresholds
\(p,p+1\) is the biproduct factor \(C^{p,q}\). The horizontal term of \(D\)
lands in the subobject being quotiented out, so the associated graded
differential is \((-1)^p d_v\). Multiplication of a differential by
\(-1\) leaves its kernel and image unchanged, so there is a canonical
identification of the first page with the unsigned vertical cohomology.
For a representative in column \(p\) that is vertically closed, its
total differential is \(d_h\) of the representative, now in column
\(p+1\). The commutation of the unsigned differentials implies that
\(d_h\) takes vertical cycles to cycles and vertical boundaries to
boundaries.

The same argument works in an abelian category without element notation.
The inclusion of the kernel of \(d_v\) in \(C^{p,q}\) has composite
with \(D\) factoring through the column \(p+1\) summand, and the
induced map to the quotient of cycles by boundaries is the map induced
by \(d_h\). This is the first filtered differential of
Theorem~\ref{res:F04}. Taking its cohomology gives the second page.

The interchange isomorphism takes \(F_{\mathrm r}\) for \(C\) to
\(F_{\mathrm c}\) for \(C^{\mathrm t}\). Equivalently, in the original
coordinates \(C^{q,p}\) on the raw row page has \(d_0=d_h\), and the raw
\(d_1\) is induced by \((-1)^q d_v\). The signed identification on
\(E_1^{p,q}\) multiplies by \((-1)^{pq}\). Increasing \(p\) by one changes
this multiplier by \((-1)^q\), which accounts for the raw sign.
Thus the two descriptions give isomorphic spectral sequences with the
same convention on later pages. The constructions use natural
sub-biproduct inclusions, kernels, and images.
\end{proof}

\begin{theorem}[Convergence, comparison, and assembly]
\label{res:B04}
Suppose \(C\) has finite diagonals. Both spectral sequences of
Proposition~\ref{res:B03} converge to \(H^*(\operatorname{Tot}C)\).
More precisely, each gives, in degree \(n\), the finite decreasing
filtration by images of
\[
  H^n(F^a\operatorname{Tot}C)\longrightarrow H^n(\operatorname{Tot}C),
\]
and \(E_\infty^{p,n-p}\) is naturally isomorphic to its \(p\)-th graded
piece. If a map \(C\to C'\) between two such double complexes is a
quasi-isomorphism on each column, or on each row, then the map on
totalizations is a quasi-isomorphism. In particular, exactness of every
column, or of every row, implies acyclicity of the total complex.

Suppose moreover \(C^{p,q}=0\) for \(q<0\), and a map of horizontal
complexes \(L^*\to C^{*,0}\) is annihilated by \(d_v\) and induces
\(L^p\simeq H_v^0(C^{p,*})\) for every \(p\). If the other vertical
cohomology vanishes, the induced map \(L^*\to\operatorname{Tot}C\) is
a quasi-isomorphism. The transposed assertion holds for an augmentation
into \(C^{0,*}\) when \(C^{p,q}=0\) for \(p<0\).
\end{theorem}

\begin{proof}
By Lemma~\ref{res:B02}, each filtration satisfies the degreewise finite
hypothesis of Theorem~\ref{res:S01}. Its conclusion supplies the
stabilized page and graded isomorphisms, including their compatibility
with the induced map on cohomology. A columnwise quasi-isomorphism is an
isomorphism on the \(E_{1,\mathrm c}\) page by
Proposition~\ref{res:B03}. Hence it is an isomorphism on the stabilized
page. In each total degree both induced filtrations are finite.
For a morphism between such filtered objects, descend through
\[
    0\longrightarrow F^{a+1}H^n
      \longrightarrow F^aH^n
      \longrightarrow \operatorname{gr}_F^a H^n
      \longrightarrow 0 .
\]
Starting at a large common zero threshold, an isomorphism on the first
and third terms gives one on the middle term by the short five lemma.
At a common lower threshold this gives the map on the whole cohomology
object. This is also the finite comparison of
Corollary~\ref{res:S03}. The row result follows by signed interchange,
and comparison with the zero double complex gives the acyclicity claim.

For the augmentation, view \(L^p\) as the double complex concentrated
in bidegrees \((p,0)\). The hypothesis about \(d_v\) makes the
augmentation a map of double complexes. It is a quasi-isomorphism on
each column: in degree zero it induces the given isomorphism, and in
the other degrees both cohomology objects vanish. Its totalization is
the stated map from \(L^*\). For the row case transpose and use
Proposition~\ref{res:B01}; on the axis the interchange multiplier is
\((-1)^0=1\).
\end{proof}

The same argument applies to chain complexes with all indices
negated. Suppose a first-quadrant chain double complex has a vertical
augmentation commuting with the horizontal maps, exact in positive
vertical degrees and inducing an isomorphism on vertical \(H_0\).
The augmentation annihilates boundaries from vertical degree one
and defines a map to a double complex concentrated in vertical
degree zero. It is a quasi-isomorphism on each column, so the induced
map from the total complex is a quasi-isomorphism as well.
An analogous statement holds
for horizontal resolutions. The module assembly and augmented
comparisons in \cite[Chapter 10]{p103} specialize this argument. Compare
\cite[Lemmas 12.25.1, 12.25.3, and 12.25.4]{stacks-homology} and
\cite[Exercise 5.6.4]{weibel}.

\section{Tensor complexes and change of rings}

For a chain double complex \(C_{p,q}\), put \(C^{i,j}=C_{-i,-j}\).
Descending column threshold \(i\ge -a\) then means original chain
column \(p\le a\), which defines the usual increasing filtration
\(F_a\) of the sum totalization. Replacing \((i,j)\) by \((-p,-q)\)
sends the cohomological page differential of bidegree
\((r,1-r)\) to a chain page differential of bidegree \((-r,r-1)\).
Thus the same proofs compute successive quotients
\(F_pH_{p+q}/F_{p-1}H_{p+q}\). For the row filtration we continue to put
its filtration index first, so that the displayed first index is the
original vertical degree in the chain double complex.

In this section rings have identity, and module actions preserve it.
We write the right \(R\)-module first in \(\otimes_R\) and in
\(\operatorname{Tor}^R_n\), and the left \(R\)-module second. Unless an
additional module action is specified, these constructions take values
in abelian groups. We use homological grading, with resolutions in
degrees \(0,1,\ldots\). The differential on the tensor double complex
of right and left complexes \(P,Q\) is
\[
   \partial(x\otimes y)=\partial_P x\otimes y
            +(-1)^i x\otimes\partial_Q y,\qquad x\in P_i.
\]
For a flat module, tensoring in the other variable is exact. In
particular, it carries the cycle and boundary short exact sequences to
short exact sequences and gives the canonical maps identifying
homology after tensoring with the tensor products of the homology
groups. This also gives the comparisons on the first pages used below.

\begin{theorem}[Balance of Tor]
\label{res:B05}
For a right \(R\)-module \(A\) and a left \(R\)-module \(M\), let
\(P_\bullet\to A\) and \(Q_\bullet\to M\) be projective resolutions in
their respective module categories. There is a natural comparison,
independent of the choices of projective resolutions,
\[
        H_n(P_\bullet\otimes_R M)
        \;\simeq\; H_n(A\otimes_R Q_\bullet).
\]
Both groups compute \(\operatorname{Tor}_n^R(A,M)\).
\end{theorem}

\begin{proof}
The double complex has terms \(P_i\otimes_R Q_j\) with the two commuting
differentials. Each \(P_i\) is a flat right module, so the augmented
column is exact after tensoring and computes \(P_i\otimes_R M\) in
degree zero. Thus the map from the total complex induced by
\(Q_\bullet\to M[0]\) is a quasi-isomorphism by
Theorem~\ref{res:B04}. Likewise each \(Q_j\) is a flat left module, so
augmentation by \(P_\bullet\to A[0]\) is a quasi-isomorphism from the
same total complex to \(A\otimes_R Q_\bullet\). The map from the first
homology group to the second is the inverse of the first total
comparison followed by the second. Both augmentations are chain maps:
they vanish away from the corresponding degree-zero row or column,
commute along that row or column, and annihilate boundaries entering
it. In the horizontal degree-zero quotient the tensor differential in
the remaining direction has sign \((-1)^0=1\).

For a morphism of module pairs, lift the two coefficient maps to maps
of projective resolutions. The comparison theorem for projective
resolutions gives lifts unique up to chain homotopy
\cite[Chapter 10]{p103}. If two lifts in the right variable differ
by \(\partial h+h\partial\), their maps on total complexes differ by
the homotopy \(x\otimes y\mapsto h(x)\otimes y\). Indeed the vertical
terms in its homotopy identity cancel since \(h\) increases \(i\) by one.
If two lifts in the left variable differ by
\(\partial h+h\partial\), use
\(x\otimes y\mapsto(-1)^i x\otimes h(y)\);
now the horizontal terms cancel. This proves independence on
homology. Both augmentations commute with the lifted pair of maps.
Applying the argument to lifts of identity maps compares two choices
of resolutions by inverse maps on homology. These identities prove
naturality and independence for the displayed comparison. The two
augmentations determine the comparison on homology. See
\cite[Chapter 11]{p103} for the module resolution models and
\cite[Application 5.6.3]{weibel} for the spectral-sequence argument.
\end{proof}

\begin{theorem}[The K\"unneth spectral sequence]
\label{res:B06}
Let \(P_\bullet\) be a chain complex of flat right \(R\)-modules with
\(P_i=0\) for \(i<L\) for some integer \(L\). For a left \(R\)-module \(M\)
there is a natural homological spectral sequence
\[
    E^2_{p,q}=\operatorname{Tor}^R_p(H_q(P),M)
       \quad\Longrightarrow\quad H_{p+q}(P_\bullet\otimes_R M).
\]
Its support lies in \(p\ge0,\ q\ge L\), with a finite target filtration
in each degree and differential \(d_r\) of bidegree \((-r,r-1)\).
\end{theorem}

\begin{proof}
Choose a projective left resolution \(Q_\bullet\to M\).
Put \(D_{i,j}=P_i\otimes_R Q_j\). In total degree \(n\), possible
nonzero terms have \(L\le i\le n,\ j=n-i\ge0\), so the
finite-diagonal comparison applies even when \(L<0\).
Flatness of \(P_i\) makes the vertical augmentation a
quasi-isomorphism on each column, giving the natural target
identification
\[
     H_n(\operatorname{Tot} D)\simeq H_n(P_\bullet\otimes_R M).
\]
For the other filtration, put \(p=j\) and \(q=i\).
Flatness of \(Q_p\) gives horizontal homology
\[
      E^1_{p,q}\simeq H_q(P)\otimes_R Q_p .
\]
Under the signed interchange that orders the \(Q\) index first, the
induced differential is \(1\otimes\partial_Q\). Equivalently, in the
untransposed coordinates it has the constant factor \((-1)^q\) in
each fixed \(q\). The cohomological sign calculation of
Proposition~\ref{res:B03} with the indices negated gives the isomorphism
of the two descriptions. Taking homology in \(p\) and using
Theorem~\ref{res:B05} gives the displayed \(E^2\). The finite-diagonal
bound and Theorem~\ref{res:B04} supply both convergence and the finite
filtration on the target.

A chain map of right complexes gives a map of the double complexes.
For a coefficient map in \(M\), use a lifted map of the projective
left resolutions. Two lifts differ by a homotopy which, after applying
\(H_q(P)\otimes_R -\), is an ordinary homotopy of the first-page
complexes in the \(p\)-direction. Hence they induce the same map on
\(E^2\), and therefore on all later pages since those maps are induced
on homology of pages. The augmentation comparisons commute with the
maps on the two targets. Taking lifts between two choices for the same
\(M\) yields natural comparisons from \(E^2\) onward and on the
filtered target. The resulting spectral sequence is the one of
\cite[Theorem 5.6.4]{weibel}.
\end{proof}

\begin{corollary}[The two-column case]
\label{res:B08}
In Theorem~\ref{res:B06}, suppose every \(H_q(P)\) admits a flat
resolution of length at most one as a right \(R\)-module. Then for
every \(n\) the target filtration gives a natural short exact sequence
\[
\begin{split}
 0 \longrightarrow H_n(P)\otimes_R M
    \longrightarrow H_n(P_\bullet\otimes_R M)
    \longrightarrow \operatorname{Tor}_1^R(H_{n-1}(P),M)
    \longrightarrow 0 .
\end{split}
\]
\end{corollary}

\begin{proof}
To see the vanishing for \(p>1\), use the length-one flat resolution
of \(H_q(P)\) in the right variable and the chosen projective resolution
of \(M\) in the left variable. The two augmentation comparisons of
Theorem~\ref{res:B05} work for this pair because the terms on each
side are flat. The computation by the right resolution has terms
only in degrees zero and one. Thus all differentials from the second
page onward vanish.
Indeed each differential moves at least two columns.
Only columns \(0,1\) can contain nonzero terms. The finite filtration
on \(H_n\) therefore has \(F_{-1}=0\), \(F_0\) equal to the first
nonzero graded piece, and \(F_1=H_n\). Their identifications are the
terms in the sequence above. The left map can be read at \(F_0\).
A cycle in \(P_n\) tensored with a lift of \(m\in M\) to \(Q_0\)
represents the corresponding class in the filtered total complex,
whose augmentation is the cycle tensored with \(m\). Changing the
lift changes this representative by a vertical boundary, with the
appropriate tensor sign. This gives the usual cycle-tensor edge map.

For example, the length bound holds if every boundary module
\(B_q(P)\) is flat. Exactness of
\(0\to Z_q(P)\to P_q\to B_{q-1}(P)\to0\) and the long exact sequence
of Tor with any left module show \(Z_q(P)\) is flat.
The sequence \(0\to B_q(P)\to Z_q(P)\to H_q(P)\to0\) is then the
required flat resolution. This is the hypothesis and calculation in
\cite[Formula 5.6.5]{weibel}. Additional splitting conclusions use
the hypotheses of the corresponding coefficient theorem, for
example those on the complexes and ring in
\cite[Chapter 11]{p103}.
\end{proof}

\begin{theorem}[Change of rings for Tor]
\label{res:B07}
Let \(f\colon R\to S\) be a homomorphism of rings, \(A\) a right
\(R\)-module and \(M\) a left \(S\)-module. Regard \(M\) as a left
\(R\)-module by restriction. There is a natural first-quadrant
homological spectral sequence
\[
 E^2_{p,q}=\operatorname{Tor}^S_p\bigl(\operatorname{Tor}^R_q(A,S),M\bigr)
       \quad\Longrightarrow\quad \operatorname{Tor}^R_{p+q}(A,M).
\]
In the inner Tor, \(S\) is the \(R\)-\(S\) bimodule obtained from \(f\)
and right multiplication. Its right \(S\)-module structure supplies
the first input of the outer Tor. The target filtration is finite
in each degree.
\end{theorem}

\begin{proof}
Choose a projective right \(R\)-resolution \(P_\bullet\to A\)
and a projective left \(S\)-resolution \(Q_\bullet\to M\).
Use the double complex
\[
    D_{i,j}=P_i\otimes_R Q_j
          \simeq (P_i\otimes_R S)\otimes_S Q_j .
\]
The comparison sends \((x\otimes s)\otimes y\) to \(x\otimes sy\),
with inverse \(x\otimes y\mapsto(x\otimes1)\otimes y\). The forward
map respects the \(R\)-balance relation because
\((xr)\otimes sy=x\otimes f(r)(sy)\), and the \(S\)-balance relation
because \(s(ty)=(st)y\). For the inverse, the corresponding relation is
\[
   ((xr)\otimes1)\otimes y=(x\otimes f(r))\otimes y
                         =(x\otimes 1)\otimes f(r)y.
\]
The two formulas are mutually inverse by \(S\)-balancing and commute
with both differentials.

The augmented complex \(Q_\bullet\to M\) remains exact in left
\(R\)-modules, since restriction of scalars is exact.
For each \(i\) the right module \(P_i\) is \(R\)-flat. Thus the
vertical augmentations \(P_i\otimes_R Q_\bullet\to P_i\otimes_R M\)
are quasi-isomorphisms. Theorem~\ref{res:B04} gives
\[
  H_n(\operatorname{Tot}D)\simeq H_n(P_\bullet\otimes_R M)
          =\operatorname{Tor}^R_n(A,M).
\]
This argument uses the flatness of \(P_i\) over \(R\) and the
exactness of restriction for \(Q_\bullet\). The flatness used for
the other filtration is over \(S\).

Specifically, each \(Q_j\) is \(S\)-flat. With the row index first,
horizontal homology gives
\[
  E^1_{p,q}\simeq H_q(P_\bullet\otimes_R S)\otimes_S Q_p
      =\operatorname{Tor}^R_q(A,S)\otimes_S Q_p .
\]
Here \(P_\bullet\otimes_R S\) is a complex of right \(S\)-modules
whose maps come from those of \(P_\bullet\). Its homology gives
the right action used in the statement. The sign comparison in
Proposition~\ref{res:B03} identifies the \(p\)-direction with the
resolution differential on \(Q\); its homology gives
\(E^2_{p,q}\) by Theorem~\ref{res:B05} over \(S\).
The two indices are nonnegative, so the total degree bounds and the
conclusion of Theorem~\ref{res:B04} apply. For maps in \(A\) and \(M\)
with \(f\) fixed, lifts to their respective resolutions induce maps
of the double complexes. Homotopic choices in \(P\) give the same maps
on the horizontal homology page; homotopic choices in \(Q\) give the
same maps on the homology of that page. Augmentation and the common
total homotopy identities in Theorem~\ref{res:B05} verify naturality
on the target. This proves the change-of-rings comparison of
\cite[Theorem 5.6.6]{weibel}.
\end{proof}

\chapter{Homotopy and Derived Categories}
\label{ch:derived-categories}

The category of complexes remembers all the differentials. Its homotopy
category identifies maps differing by a cochain homotopy. Its derived
category also makes every quasi-isomorphism invertible. We first keep track
of the maps and signs in the homotopy category, so that the long exact
sequence obtained after localization has the ordinary connecting maps.

\section{Shifts, cones, and triangles}

\begin{proposition}[Shifts and the homotopy category]\label{res:K01}
Let \(\mathcal A\) be an additive category. For a complex \(K\) and
\(m\in\mathbb Z\), put
\[
 (K[m])^n=K^{n+m},\qquad d_{K[m]}^n=(-1)^m d_K^{\,n+m}.
\]
The shifts act on the homotopy category \(\mathsf K(\mathcal A)\) and
compose by addition. If \(\mathcal A\) is abelian, the cohomology functors
on complexes factor through \(\mathsf K(\mathcal A)\), and the
identifications \(H^n(K[m])=H^{n+m}(K)\) are natural.
\end{proposition}
\begin{proof}
The shifted differential squares to zero. If \(f-g=d_L s+s d_K\), where
\(s^n:K^n\to L^{n-1}\), the components \((-1)^m s^{n+m}\) give the
corresponding homotopy on the shifts. Morphisms homotopic to zero form an
additive ideal because precomposition or postcomposition of a
null-homotopy by a cochain map gives a null-homotopy for the composite. Quotienting each
Hom group by that ideal defines \(\mathsf K(\mathcal A)\). On cycles,
the term \(s d_K\) is zero and \(d_L s\) factors through boundaries, so
homotopic maps induce the same cohomology map. Multiplying a differential
by \(-1\) leaves its kernel and image subobjects unchanged. This gives
the cohomology identification and its naturality. See
\cite[Tags 013H and 05RS]{stacks-derived} and
\cite[Section 4.1]{yekutieli}.
\end{proof}

\begin{proposition}[The mapping cone and its maps]\label{res:K02}
For a cochain map \(f:X\to Y\) in an additive category, define
\[
 C(f)^n=Y^n\oplus X^{n+1},\qquad
 d_{C(f)}(y,x)=(d_Y y+f x,-d_X x).
\]
There are cochain maps \(i_f:Y\to C(f)\) and \(p_f:C(f)\to X[1]\),
the inclusion and the projection. Given maps \(a:X\to X'\), \(b:Y\to
Y'\), and a homotopy \(bf-f'a=d_{Y'}s+s d_X\), the formula
\[
 (y,x)\longmapsto (b y+s x,a x)
\]
defines a cone map \(C(f)\to C(f')\). It intertwines the inclusions and
projections with \(b\) and \(a[1]\), respectively.
\end{proposition}
\begin{proof}
In the square of the cone differential the only potentially nonzero
off-diagonal entry is \(d_Y f-f d_X\), which vanishes. For the claimed
cone map, equality of the two off-diagonal entries of the chain map
equation is \(d_{Y'}s+f'a=bf-sd_X\). This is precisely the homotopy
equation. The other entries are the equations for \(a\) and \(b\).
The intertwining identities are immediate from the inclusion and
projection matrices. In an exactly commuting square, take \(s=0\).
For exactly commuting squares this construction preserves composition.
For a square specified only up to homotopy, the formula uses the
chosen homotopy. See \cite[Tags 014E, 014F, and 08RI]{stacks-derived}.
\end{proof}

\begin{proposition}[Contractibility and the cone criterion]\label{res:K03}
The cone of the identity map is contractible. In an abelian category,
\(f\) is a quasi-isomorphism if and only if \(C(f)\) is acyclic.
There exist quasi-isomorphisms of module complexes which are not
homotopy equivalences.
\end{proposition}
\begin{proof}
On \(C(1_X)\) the homotopy \(h^n(y,x)=(0,y)\) satisfies
\(d h+h d=1\). For any \(f\) there is a degreewise split short exact
sequence
\[
 0\longrightarrow Y\xrightarrow{i_f} C(f)
       \xrightarrow{p_f} X[1]\longrightarrow 0.
\]
The connecting map \(H^n(X[1])\to H^{n+1}(Y)\) is \(H^{n+1}(f)\).
Indeed, lifting a cycle of \(X[1]\) by the section of the projection
leaves \(f x\) in the first component after applying the differential. The
long exact sequence proves the criterion in both directions.

For the last assertion, put \(\mathbb Z\xrightarrow{2}\mathbb Z\)
in degrees \(-1,0\). The augmentation to \((\mathbb Z/2)[0]\) is a
quasi-isomorphism. Every map back is zero because every homomorphism
\(\mathbb Z/2\to\mathbb Z\) is zero. Its composite with the augmentation
cannot be homotopic to the identity of the stalk since they have
different \(H^0\). This is the module example from
\cite[Chapter 10]{p103}. The cone criterion is also recorded in
\cite[Tag 05RT]{stacks-derived}.
\end{proof}

\begin{proposition}[Degreewise split sequences]\label{res:K04}
Suppose \(0\to A\xrightarrow{u}B\xrightarrow{v}C\to0\) is degreewise
split. Choose graded identifications \(B^n=A^n\oplus C^n\) compatible
with \(u,v\), and write
\[
 d_B(a,c)=(d_A a+\alpha c,d_C c).
\]
Then \(\alpha:C\to A[1]\) is a cochain map. Its homotopy class is
independent of the splitting and natural in maps of short exact sequences.
The cone row \((u,i_u,-p_u)\) is canonically isomorphic in
\(\mathsf K(\mathcal A)\) to \((u,v,\alpha)\).
\end{proposition}
\begin{proof}
The off-diagonal entry of \(d_B^2=0\) gives
\(d_A\alpha+\alpha d_C=0\), with the sign required for \(A[1]\).
If one graded section differs from another by \(u t\), their
off-diagonal terms differ by \(d_A t-t d_C\), which is a
null-homotopy for maps into \(A[1]\). For a map of short exact sequences
\((a,b,c)\), the difference of the two graded lifts \(b s-s'c\) of \(c\)
is \(u' t\). The equation
\(a\alpha-\alpha'c=d_{A'}t-t d_C\) proves naturality in the homotopy
category.

The map \(q:C(u)\to C\) is \(v\) on \(B\) and zero on \(A[1]\).
In the graded coordinates \((a,c,a')\) of \(C(u)\), a chain section
is \(c\mapsto(0,c,-\alpha c)\). The kernel is \(C(1_A)\), which is
contractible by Proposition~\ref{res:K03}, hence \(q\) is a homotopy
equivalence. Under the chain section, \(-p_u\) is precisely \(\alpha\).
The ensuing sign convention is that the third arrow of a cone triangle
is \(-p_f\). It makes the boundary from a split sequence agree with
the usual cocycle-lift boundary. See
\cite[Tags 05SS and 014L]{stacks-derived}.
\end{proof}

\paragraph{Comparison of five exact terms.}
In a commuting diagram of two exact rows \(A\to B\to C\to D\to E\),
the middle vertical map is an isomorphism when the four others
are. This comparison works in an abelian category using the epic
pullbacks described in Chapter~\ref{ch:filtrations}. For injectivity,
let an element of \(C\) map to zero in the bottom row. Its image in
\(D\) is zero by the comparison at \(D\), so it lifts locally to
\(B\). The image of this lift in the bottom row comes from the
bottom \(A\) by exactness. Lift it to the top \(A\); injectivity
at \(B\) shows that the original lift already comes from \(A\),
so its image in \(C\) is zero. For surjectivity, take an element
of the bottom \(C\) and lift its image in the bottom \(D\) to the
top \(D\). That lift dies in the top \(E\) by injectivity at \(E\).
It therefore lifts locally to the top \(C\). The difference from
the desired bottom element comes from the bottom \(B\), and can
be corrected by lifting through the isomorphism at \(B\).
The zero and image conditions descend along the epic covers.

\begin{theorem}[The triangulated structure on the homotopy category]
\label{res:K05}
The shift \([1]\) and the triangles isomorphic to
\[
 X\xrightarrow{f}Y\xrightarrow{i_f}C(f)
     \xrightarrow{-p_f}X[1]
\]
make \(\mathsf K(\mathcal A)\) a triangulated category.
\end{theorem}
\begin{proof}
Every map is the first arrow of its cone triangle. The identity has a
contractible cone, so its row is isomorphic to \(X\to X\to0\to X[1]\).
The class of triangles is defined up to isomorphism.

\emph{Rotation.} The cone of \(i_f\) has coordinates
\((y,x,y')\in Y^n\oplus X^{n+1}\oplus Y^{n+1}\) and differential
\[
 (d_Yy+fx+y',-d_Xx,-d_Yy').
\]
The map \((y,x,y')\mapsto-x\) to \(X[1]\) has chain section
\(x\mapsto(0,-x,fx)\) and contractible kernel \(C(1_Y)\). After this
identification, the last arrow of the cone triangle for \(i_f\) is
\(-f[1]\). This proves rotation with the stated sign. For inverse
rotation, take the cone of \(p_f[-1]:C(f)[-1]\to X\). Its coordinates
are \((x,y,x')\in X^n\oplus Y^n\oplus X^{n+1}\) with differential
\((d_Xx+x',d_Yy+fx',-d_Xx')\). The map to \(Y\) is
\((x,y,x')\mapsto fx-y\), with chain section \(y\mapsto(0,-y,0)\)
and kernel \(C(1_X)\). The inclusion of \(X\) corresponds to \(f\),
and the last map on the section corresponds to \(i_f\). Thus inverse
rotation also preserves the class.

\emph{Morphisms of triangles.} A commutative square in the homotopy
category can be represented by a square of cochain maps together with
a homotopy between the two composites. Proposition~\ref{res:K02}
supplies the third map of triangles.

If the first two components of a morphism of cone triangles are
isomorphisms in \(\mathsf K(\mathcal A)\), its third component is an
isomorphism as well. To see this without using an octahedron, fix a
complex \(T\). If a cochain map \(u:T\to Y\) has \(i_fu\) null-homotopic,
write the homotopy into \(C(f)^{n-1}=Y^{n-1}\oplus X^n\) as \((s,t)\).
Its second component says that \(t:T\to X\) is a cochain map. Its
first component gives \(u-ft=d_Ys+s d_T\). Conversely, \(i_ff\)
is null-homotopic by \(x\mapsto(0,x)\). This proves exactness at
\(\operatorname{Hom}(T,Y)\). Applying the same calculation to the
rotations gives exact rows of five abelian groups with middle group
\(\operatorname{Hom}(T,C(f))\) and adjacent groups from \(X,Y\) and
their shifts. The morphism of triangles induces a map of these rows,
and the four outer vertical maps are isomorphisms. The preceding
five-term comparison makes the middle map an isomorphism for each \(T\). If \(c:C\to C'\)
is the third component, surjectivity for \(T=C'\) gives a right
inverse to \(c\), and injectivity for \(T=C\) makes it a left
inverse. In particular the morphism axiom, applied to identities
on the first two terms, compares any two distinguished triangles
on a fixed first arrow by an isomorphism fixing those terms.

\emph{The octahedron.} For \(X\xrightarrow{f}Y\xrightarrow{g}Z\)
the two maps
\[
 a:C(f)\to C(gf),\quad (y,x)\mapsto(gy,x),\qquad
 b:C(gf)\to C(g),\quad (z,x)\mapsto(z,fx)
\]
extend the squares on the first two terms. They give the first two
arrows of the fourth triangle. In fact, in the coordinates
\((z,x,y',x')\) of \(C(a)\), its differential is
\[
 (d_Zz+gfx+gy',\,-d_Xx+x',\,-d_Yy'-fx',\,d_Xx').
\]
The chain map \(C(a)\to C(g)\), \((z,x,y',x')\mapsto(z,fx+y')\),
has chain section \((z,y')\mapsto(z,0,y',0)\). Its kernel is
\(C(1_{X[1]})\), with coordinates \((x,x')\) and differential
\((-d_Xx+x',d_Xx')\). Consequently the cone triangle for \(a\) is
isomorphic to
\[
 C(f)\xrightarrow{a}C(gf)\xrightarrow{b}C(g)
   \xrightarrow{i_f[1](-p_g)}C(f)[1].
\]
The outer compatibilities are \(a i_f=i_{gf}g\),
\(b i_{gf}=i_g\), \(p_{gf}a=p_f\), and \(p_g b=f[1]p_{gf}\),
all read directly from the coordinate maps. The closing arrow of the
fourth triangle is \((z,y')\mapsto(-y',0)\), so it agrees with the
composite required by the same choice of cone sign. This gives the
octahedron for composable representatives. For prescribed triangles
on \(f,g,gf\), choose the three comparison isomorphisms just constructed
and conjugate the maps between their third terms by those
isomorphisms. These comparison maps fix \(X,Y,Z\) and their shifts,
so all commuting faces remain commuting. The fourth triangle remains
distinguished because it is isomorphic to the one for the explicit
cone maps. Compare \cite[Tags 014C, 014Q, and 014S]{stacks-derived}.
\end{proof}

\begin{proposition}[Exact sequences from triangles]\label{res:K06}
For a distinguished triangle, cohomology gives a natural long exact
sequence if \(\mathcal A\) is abelian. Each of
\(\operatorname{Hom}_{\mathsf K(\mathcal A)}(T,-)\) and
\(\operatorname{Hom}_{\mathsf K(\mathcal A)}(-,T)\) gives an exact
sequence extended by shifts, with arrows reversed in the second case.
\end{proposition}
\begin{proof}
By isomorphism and rotation it suffices to check exactness at the
middle term of a cone row. Cohomology has the asserted exactness by
the split sequence in Proposition~\ref{res:K03}. Rotation appends the
shifted terms. Use the arrows of the triangle, including its
closing sign. They give natural maps under morphisms of triangles.

For the covariant Hom sequence, suppose \(u:T\to Y\) has \(i_fu\)
null-homotopic, with homotopy \((s,t)\) into
\(Y^{n-1}\oplus X^n=C(f)^{n-1}\). Its second component equation says
that \(t:T\to X\) is a cochain map. The first says
\(u-ft=d_Ys+s d_T\), so \([u]\) is in the image of \([f]\).
Conversely \(i_ff\) is null-homotopic by the homotopy \(x\mapsto(0,x)\).
For the contravariant sequence let \(v:Y\to T\) satisfy
\(vf=d_Ts+s d_X\). The map \(C(f)\to T\),
\((y,x)\mapsto vy+sx\), is a cochain map and restricts to \(v\).
The converse again follows from \(i_ff\simeq0\). Naturality follows
from composition. The shifted long sequences follow by applying
these computations to the rotated cone rows.
Compare \cite[Tags 0149 and 05RS]{stacks-derived} and
\cite[Propositions 5.2.6--5.2.7]{yekutieli}.
\end{proof}

\section{Quasi-isomorphisms and localization}

For the remainder of the chapter \(\mathcal A\) is abelian. The
localization is taken in a universe in which the roofs under
consideration form sets. Bounds on complexes, when imposed, are bounds
on their terms unless cohomology is explicitly specified.

\begin{proposition}[The acyclic subcategory]\label{res:D01}
The acyclic complexes form a thick triangulated subcategory of
\(\mathsf K(\mathcal A)\). Its cone class, the maps whose cones belong
to this subcategory, is the class of quasi-isomorphisms.
\end{proposition}
\begin{proof}
Acyclicity means \(H^n=0\) for all \(n\). It is preserved under
isomorphism in the homotopy category and under shifts by
Proposition~\ref{res:K01}, under cones between acyclic objects by
Proposition~\ref{res:K06}, and under retracts because the cohomology of
a retract is a retract of the cohomology. These are the required
closure properties. The assertion about cones is
Proposition~\ref{res:K03}. In particular quasi-isomorphisms have
two-out-of-three because their cohomology maps do.
\end{proof}

\begin{theorem}[Roofs and their refinements]\label{res:D02}
The quasi-isomorphisms in \(\mathsf K(\mathcal A)\) admit calculi of
right and left fractions. In particular a morphism in the localization
can be written
\[
 X\xleftarrow{s}T\xrightarrow{f}Y,\qquad s\text{ a quasi-isomorphism}.
\]
Two roofs are equal precisely when they can be refined to a common
roof, with both arrows agreeing in \(\mathsf K(\mathcal A)\). There is
also the dual presentation \(X\to T\xleftarrow{s}Y\).
\end{theorem}
\begin{proof}
Here are the two fraction conditions for roofs of the displayed
orientation. Given \(f:X\to Z\) and a quasi-isomorphism \(s:Y\to Z\),
set \(W=C(f,-s)[-1]\), where \((f,-s):X\oplus Y\to Z\).
The cone projection, shifted back, gives maps \(W\to X\) and \(W\to Y\).
Their composites with \(f\) and \(s\) agree in \(\mathsf K(\mathcal A)\)
by the cone homotopy. The projection to \(X\) is a quasi-isomorphism
because its kernel is \(C(-s)[-1]\) and the sequence is degreewise split.
This gives the lifting square for a common denominator.

For cancellation, suppose \(s:Y\to Z\) is a quasi-isomorphism and
\(sf=sg\) in \(\mathsf K(\mathcal A)\) for \(f,g:X\to Y\).
Exactness of Hom applied to the cone triangle for \(s\) factors \(f-g\)
through a map \(u:X\to C(s)[-1]\). The projection
\(C(u)[-1]\to X\) is a quasi-isomorphism since \(C(s)[-1]\)
is acyclic, and precomposing with it kills \(f-g\). This is the
cancellation condition. The left calculus follows from the same
argument with cones in place of the shifted cones, equivalently by
applying the right calculus to the opposite category with cochain
degrees reversed.

For clarity, these conditions give the roof description as follows.
Identities and composites of denominators are denominators. For two
successive roofs, apply the lifting square to the middle numerator and
the next denominator. Precompose the first denominator and the last
numerator with the two maps of that square. Any two choices of square
can first be brought to a common denominator by a further lifting
square, then made equal where necessary by cancellation. The same
construction compares a composite after either input is refined.
For transitivity of common refinement, the maps to the middle roof
are quasi-isomorphisms by two-out-of-three. Refine those two maps by a
lifting square. To check associativity, write the roofs as
\(X_{i-1}\xleftarrow{s_i}T_i\xrightarrow{f_i}X_i\) for \(i=1,2,3\).
Choose \(W\to T_1,T_2\) and \(P\to T_2,T_3\) to lift \(f_1\)
against \(s_2\) and \(f_2\) against \(s_3\), respectively.
The maps \(W\to T_1\) and \(P\to T_2\) are denominators.
Next lift \(W\to T_2\) against \(P\to T_2\), obtaining a
denominator \(Q\to W\) and a map \(Q\to P\).
Then \(Q\to W\to T_1\to X_0\) and \(Q\to P\to T_3\to X_3\)
give the same composite roof for both bracketings, by the three
commuting lifting squares. The image of a
roof under a functor \(F\) inverting quasi-isomorphisms has to be
\(F(f)F(s)^{-1}\). That formula is unchanged under refinement and
respects composition. This proves the localization property and
the claimed equality criterion. See also
\cite[Tags 05R2, 05R6, and 05RU]{stacks-derived} and
\cite[Sections 6.1--6.2]{yekutieli}.
\end{proof}

\begin{theorem}[The derived category]\label{res:D03}
Let \(\mathsf D(\mathcal A)\) be \(\mathsf K(\mathcal A)\) localized at
quasi-isomorphisms. It is the Verdier quotient by the acyclic
subcategory. The quotient \(Q\) preserves distinguished triangles,
and the \(H^n\) factor through \(Q\). A morphism in \(\mathsf D(\mathcal A)\)
is invertible exactly when all its cohomology maps are invertible.
\end{theorem}
\begin{proof}
Write \(W\) for the quasi-isomorphisms. They are the maps whose cones
are acyclic, by Proposition~\ref{res:D01}. Both fraction calculi
hold by Theorem~\ref{res:D02}. This also gives an additive
localization: for two roofs with the same ends take a common
denominator and add the numerators. Any two such choices have a
common further refinement, which proves independence and the
abelian group identities. When precomposing two roofs written
over the same denominator, use the same lifting square for both.
When postcomposing, lift the two numerators against the next
denominator over a common refinement. The sum of their lifts
is a lift of the sum. This proves bilinearity and makes \(Q\)
additive. The zero object stays zero: in a roof out of zero the
denominator has acyclic source and is refined by the map from zero,
and in a roof to zero the numerator is already zero. The biproduct
equations therefore pass to the localization.

For each integer \(m\), the shift \([m]\) preserves \(W\) by
Proposition~\ref{res:K01}. Shifting both arrows of a roof thus
gives the descended shifts, with the same addition law on
shift indices and with the canonical commutation isomorphisms
for \(Q\). Consider now a morphism between distinguished triangles
in \(\mathsf K(\mathcal A)\) whose first two components belong to
\(W\). In each degree the five consecutive groups or objects
\[
 H^n(X)\longrightarrow H^n(Y)\longrightarrow H^n(Z)
       \longrightarrow H^{n+1}(X)\longrightarrow H^{n+1}(Y)
\]
form an exact row by Proposition~\ref{res:K06}, where \(X,Y,Z\)
are the first three objects of the triangle. Naturality gives a
map of exact rows. Both maps on the left and both maps on the
right are isomorphisms, so the five-term comparison just proved
makes the middle map an isomorphism. The third component therefore belongs to \(W\).
In particular a commuting square on the first arrows with vertical
maps in \(W\) extends by Theorem~\ref{res:K05} to a morphism of
triangles all of whose components lie in \(W\).

These give \texttt{IsCompatibleWithTriangulation} for \(W\) in the
fixed Mathlib localization theorem. The required left calculus is
Theorem~\ref{res:D02}. The zero objects and preadditive structures
are those described above. Shifts are additive before localization
by the termwise addition of maps, and afterward because shifting
the sum of roofs over a common denominator gives the sum of their
shifts over the shifted denominator. The commutation of \(Q\) with
these shifts was constructed above. These are the inputs to
\cite[\texttt{CategoryTheory.Triangulated.Localization}]{mathlib-triangulated}.
Its declarations \texttt{pretriangulated},
\texttt{isTriangulated\_functor}, and \texttt{isTriangulated}
give the pretriangulated structure with distinguished triangles
isomorphic to images of distinguished triangles, exactness of
\(Q\), and the octahedral axiom, respectively. The starting
category satisfies the triangulated axioms by Theorem~\ref{res:K05}.

To verify the Verdier universal property, an exact functor
inverting \(W\) kills each acyclic object, since \(0\to C\) lies
in \(W\) for acyclic \(C\). Conversely let an exact functor \(F\)
kill all acyclic objects. For \(s\in W\) it sends the cone
triangle of \(s\) to a distinguished triangle whose third
object is zero. The first arrow of such a triangle is an
isomorphism. Indeed, representable Hom is exact by rotation and
the morphism axiom. For the latter, if \(A\to B\to C\) is the start
of a distinguished triangle and a map \(T\to B\) composes to zero
in \(C\), extend the square from \(T\to0\) to \(B\to C\) to a
map from \(T\to0\to T[1]\xrightarrow{-1}T[1]\) to the rotated
triangle. The commuting last square, shifted back, gives the
factorization through \(A\). Successive arrow composites vanish
by comparing an identity triangle to the given triangle and then
rotating. With third term zero, the resulting Hom sequence makes
the first arrow an isomorphism on every representable Hom, hence
an isomorphism. Consequently \(F\) inverts \(W\).

By Theorem~\ref{res:D02}, \(F\) factors uniquely through \(Q\), by
sending a roof \((s,f)\) to \(F(f)F(s)^{-1}\). The descended
functor is additive by the common-denominator sum. The natural
commutation with shifts of \(F\) is natural on roofs as well:
the squares for \(F(s)\) may be inverted, then composed with the
squares for \(F(f)\). It takes a triangle isomorphic to an image
under \(Q\) to a triangle isomorphic to the corresponding image
under \(F\). Thus the factorization is exact. Its shift data and
the factorizations of natural transformations are determined by
their restrictions to \(Q\): every object is represented in
\(\mathsf K(\mathcal A)\) and naturality on inverses of denominators
follows from naturality on denominators. This proves the exact
universal property for the acyclic subcategory.

Each \(H^n\) inverts quasi-isomorphisms by definition. The localization
property in Theorem~\ref{res:D02} makes its factorization unique.
If a roof represents a map \(w\) with all \(H^n(w)\) invertible, then
the numerator has all \(H^n\) invertible because the denominator does.
Both arrows of the roof therefore become isomorphisms, so \(w\) is an
isomorphism. The converse is functorial. In particular an object of
the derived category with all cohomology zero is zero.
The long cohomology sequence from Proposition~\ref{res:K06} descends
to derived triangles by their description as localized cone rows.
\end{proof}

\begin{theorem}[Triangles from short exact sequences]\label{res:D04}
Every short exact sequence \(0\to A\xrightarrow{u}B\xrightarrow{v}C
\to0\) of complexes induces a distinguished triangle
\[
 A\xrightarrow{u}B\xrightarrow{v}C
   \xrightarrow{\delta}A[1]\quad\text{in }\mathsf D(\mathcal A).
\]
The construction is natural in maps of short exact sequences and its
cohomology sequence has the usual connecting homomorphisms.
\end{theorem}
\begin{proof}
The map \(q:C(u)\to C\), \(q(b,a)=v(b)\), is a chain epimorphism
whose kernel is \(C(1_A)\). Thus \(Q(q)\) is invertible. Transporting
the cone triangle defines
\(\delta=Q(-p_u)Q(q)^{-1}\). A map of short exact sequences induces
the strict cone map from Proposition~\ref{res:K02} with \(s=0\).
It intertwines \(q\) and \(-p_u\). Inverting \(Q(q)\) proves
naturality without any choice of a splitting. A cone cocycle
\((b,a)\) projecting to the given cohomology class satisfies
\(d_Bb=-u(a)\). The ordinary boundary is the class of \(-a\),
which is induced by \(-p_u\). In an abelian category this is the
same diagram chase in the pullback of \(Z^n(C)\) along \(B^n\to
C^n\). The resulting boundary kills boundaries and is given by
\(-p_u\) after the isomorphism on cohomology.
See \cite[Tag 0152]{stacks-derived}.
\end{proof}

\section{Good truncation and bounded models}

\begin{proposition}[Good truncation]\label{res:D05}
For any cochain complex \(K\), the good truncations are
\[
\begin{aligned}
 (\tau_{\le n}K)^i&=
  \begin{cases}
   K^i & i<n,\\
   \ker(d_K^n) & i=n,\\
   0 & i>n,
  \end{cases}
&
 (\tau_{\ge n}K)^i&=
  \begin{cases}
   0 & i<n,\\
   \operatorname{coker}(d_K^{n-1}) & i=n,\\
   K^i & i>n.
  \end{cases}
\end{aligned}
\]
They preserve cohomology on the indicated side and give zero on
the other side. They induce functors on the derived category and a
natural triangle
\[
 \tau_{\le n}K\longrightarrow K\longrightarrow\tau_{\ge n+1}K
       \longrightarrow(\tau_{\le n}K)[1].
\]
Shifts commute with these constructions, with the truncation indices
translated by the shift.
\end{proposition}
\begin{proof}
The incoming differential at the kernel factors through it by
\(d^2=0\). The outgoing differential at the cokernel descends for
the same reason. The edge cohomology groups are both the relevant
cycle object modulo the preceding boundaries. The other cohomology
degrees follow from the displayed complexes. Chain maps between
the truncations inherit the induced kernel or cokernel map. For the
upper truncation a homotopy restricts at \(n\) to the source cycles,
where the discarded term \(s^{n+1}d^n\) vanishes. For the lower
truncation the degree \(n+1\) homotopy component is followed by the
projection onto the target cokernel in degree \(n\). This projection
kills the discarded term \(d^{n-1}s^n\) in degree \(n\). These observations make truncation
functorial on the homotopy category. The cohomology computation
shows it preserves quasi-isomorphisms, so it descends to
\(\mathsf D(\mathcal A)\).

The quotient \(K/\tau_{\le n}K\) maps naturally to
\(\tau_{\ge n+1}K\). Its kernel is the two-term identity complex on
\(\operatorname{im}d_K^n\) in degrees \(n,n+1\), since
\(K^n/\ker(d_K^n)=\operatorname{im}d_K^n\). Thus the quotient map
is a quasi-isomorphism, and Theorem~\ref{res:D04} supplies the
triangle. Naturality on complexes passes to the localization because
the triangle maps commute with quasi-isomorphisms and therefore
with their inverses. For example
\((\tau_{\le n}K)[m]\cong\tau_{\le n-m}(K[m])\), since multiplying
the differential by a sign does not change its kernel. The same
calculation applies to the lower truncation.
See \cite[Tags 05RV and 08J5]{stacks-derived} and
\cite[Definition 7.3.6]{yekutieli}.
\end{proof}

\begin{proposition}[Strict models and bounded localizations]\label{res:D06}
Write \(\mathsf K^+,\mathsf K^-,\mathsf K^b\) for the homotopy
categories with the indicated bounds on terms, and
\(\mathsf D^+,\mathsf D^-,\mathsf D^b\) for the full subcategories
with the corresponding bounds on cohomology. Localization of each
\(\mathsf K^\ast\) at its quasi-isomorphisms gives an equivalence
with \(\mathsf D^\ast\), for \(\ast\in\{+,-,b\}\). These are
triangulated subcategories of the corresponding unbounded
categories.
\end{proposition}
\begin{proof}
Upper and lower cohomological bounds are represented by strict ones
by Proposition~\ref{res:D05}. If both bounds are imposed, the two
truncations can be applied consecutively. Their bounds are preserved
under shifts and any cone of a map between objects with the chosen
type of bounds again has those bounds (for cohomological bounds,
use the long exact sequence). Hence the full subcategories are
triangulated. The fraction constructions of Theorem~\ref{res:D02}
use only shifts, cones and finite sums, and thus give fraction
calculi on the corresponding bounded homotopy categories.

For fullness, take a right roof between strict bounded representatives
\(X,Y\), with middle complex \(T\). If \(X,Y\) are zero below \(a\),
then \(H^i(T)=0\) for \(i<a\) because the denominator is a
quasi-isomorphism. Both maps from \(T\) factor through
\(T\to\tau_{\ge a}T\), since their degree \(a\) maps kill
\(\operatorname{im}d_T^{a-1}\). That quotient is a quasi-isomorphism.
If \(X,Y\) are zero above \(b\), precompose both maps by
\(\tau_{\le b}T\to T\), a quasi-isomorphism by the upper bound on
\(H^i(T)\). With both bounds apply both operations. For faithfulness,
an equality of two right roofs has a common refinement in
\(\mathsf K\). The strict bounded complexes that occur in the two
roofs admit a common bound of the relevant type. Truncating the
refining complex exactly as above keeps the comparison maps and
the quasi-isomorphisms. Equality of composites in the homotopy
category is preserved for upper truncation by precomposition.
For lower truncation a homotopy into a complex zero below \(a\)
descends, because its degree \(a\) component has zero target and
its degree \(a+1\) component has unchanged source. Thus the
truncated refinement proves equality in the bounded localization.
See \cite[Tags 05RQ, 05RV, and 05RW]{stacks-derived}.
\end{proof}

\begin{theorem}[The standard heart]\label{res:D07}
The cohomological subcategories
\[
 \mathsf D^{\le0}=\{K:H^i(K)=0\text{ for }i>0\},\qquad
 \mathsf D^{\ge0}=\{K:H^i(K)=0\text{ for }i<0\}
\]
form the standard \(t\)-structure. Its heart, their intersection, is
equivalent to \(\mathcal A\) by the stalk embedding \(A\mapsto A[0]\)
and the inverse \(H^0\).
\end{theorem}
\begin{proof}
Proposition~\ref{res:K01} gives
\(\mathsf D^{\le0}[1]\subseteq\mathsf D^{\le0}\) and
\(\mathsf D^{\ge0}[-1]\subseteq\mathsf D^{\ge0}\). For
orthogonality take \(X\in\mathsf D^{\le0}\), \(Y\in\mathsf D^{\ge1}\).
Represent \(X,Y\) by complexes supported in degrees at most \(0\)
and at least \(1\), respectively. In a right roof \(X\leftarrow T
\to Y\), replace \(T\) by \(\tau_{\le0}T\) without changing the
roof, since \(X\leftarrow T\) is a quasi-isomorphism. The map from
this new middle term to \(Y\) is zero degree by degree. This proves
\(\operatorname{Hom}(X,Y)=0\). The truncation triangle for \(n=0\)
provides the remaining \(t\)-structure axiom.

If an object has cohomology only in degree zero, apply the two good
truncations at zero to obtain its stalk \(H^0(K)[0]\) with canonical
quasi-isomorphisms in the truncation zigzag. Between two stalks
the roof computation of Proposition~\ref{res:D06} can be made
with both bounds zero. Its denominator is then an ordinary
isomorphism in \(\mathcal A\). Thus the stalk embedding is fully
faithful and \(H^0\) is its inverse on the heart. In particular for
stalks \(A,B\) and \(m<0\) we have
\(\operatorname{Hom}_{\mathsf D}(A[0],B[m])=0\) by orthogonality.
See \cite[Tags 05RV and 06XS]{stacks-derived} and
\cite[Proposition 7.5.1]{yekutieli}.
\end{proof}

\begin{proposition}[The bounded Postnikov tower and splitting]
\label{res:D08}
An object \(K\) with bounded cohomology is obtained from its cohomology
stalks by the finite tower of triangles
\[
 \tau_{\le i-1}K\longrightarrow\tau_{\le i}K
 \longrightarrow H^i(K)[-i]\xrightarrow{\epsilon_i}
            (\tau_{\le i-1}K)[1].
\]
If
\(\operatorname{Ext}^{\,i-j+1}_{\mathcal A}(H^i(K),H^j(K))=0\)
for all \(j<i\), then \(K\) is isomorphic to
\(\bigoplus_i H^i(K)[-i]\). Here Ext denotes the shifted Hom groups
between stalks. The comparison with resolutions is given in
Theorem~\ref{res:R11}. The splitting obtained from these
vanishing conditions need not be canonical.
\end{proposition}
\begin{proof}
Apply Proposition~\ref{res:D05} to \(\tau_{\le i}K\) and cut at
\(i-1\). The upper and lower pieces are the stated truncation and
the degree \(i\) cohomology stalk. The tower begins at zero and
ends at \(K\) because of the cohomological bounds. The maps
\(\epsilon_i\) record the attachments.

Inductively suppose \(\tau_{\le i-1}K\cong
\bigoplus_{j<i} H^j(K)[-j]\). The components of \(\epsilon_i\) then
belong to the groups
\[
 \operatorname{Hom}_{\mathsf D(\mathcal A)}
       \bigl(H^i(K)[0],H^j(K)[\,i-j+1\,]\bigr).
\]
The direct sum is finite and all these groups vanish by assumption.
Thus \(\epsilon_i=0\). Applying
\(\operatorname{Hom}_{\mathsf D}(H^i(K)[-i],-)\) to the triangle
gives a section of the map to \(H^i(K)[-i]\). Indeed the Hom
exactness in Proposition~\ref{res:K06} holds in any triangulated
category by the triangle axioms. If \(g u=0\) for \(g\) the
second arrow, the square from the zero triangle
\(T\to 0\to T[1]\xrightarrow{-1}T[1]\) to the rotated target
triangle lifts by the morphism axiom, and unshifting the lift
expresses \(u\) through the first arrow. The contravariant case
is dual. The section and the inclusion from the preceding
truncation now give a map from their direct sum to
\(\tau_{\le i}K\). It is an isomorphism on cohomology by the long
exact sequence, so it is an isomorphism by
Theorem~\ref{res:D03}. The section is a choice, which proves the
claim without assigning a preferred splitting.
Compare \cite[Tags 08J5, 0GM4, and 0EWX]{stacks-derived}.
\end{proof}

\chapter{Bounded Total Derived Functors}
\label{ch:derived-functors}

Throughout this chapter, categories are abelian, and the functors used to
construct total derived functors are additive. An assertion that a category
has enough injectives or projectives applies to the source of the corresponding
construction. We use \(\mathrm K(\mathcal A)\) for its homotopy category,
\(D(\mathcal A)\) for its localization, and \(Q_{\mathcal A}\) for the
localization functor from \(\mathrm K(\mathcal A)\). The superscripts \(+\)
and \(-\) mean bounded below and bounded above, respectively. On the complex
side the bounds are on terms, and on the derived side they are on cohomology.
The comparison between these descriptions is Proposition~\ref{res:D06}.

\section{Replacements and computation of morphisms}

We first construct replacements of strictly bounded complexes. Here is the
projective construction in an abelian category with enough projectives.
For a complex \(X\) zero in degrees \(>N\), start with \(P^i=0\) for \(i>N\).
Suppose the part of \(P\to X\) above degree \(n\) has been constructed. Set
\[
  W^n=X^n\times_{Z^{n+1}X}Z^{n+1}P,
\]
using \(d_X^n\) and the already defined comparison on cycles. Choose an
epimorphism \(P^n\twoheadrightarrow W^n\) from a projective object. Its first
projection defines \(p^n\), and its second projection, followed by the cycle
inclusion, defines \(d_P^n\). The pullback identities give \(p^{n+1}d_P^n=d_X^np^n\)
and \(d_P^{n+1}d_P^n=0\).

Pulling back \(P^n\twoheadrightarrow W^n\) along \(Z^nX\to W^n\),
\(z\mapsto(z,0)\), gives \(Z^nP\twoheadrightarrow Z^nX\). Thus
\(W^n\to X^n\) is an epimorphism by the same assertion in degree \(n+1\),
so \(p^n\) is an epimorphism. Finally, the pullback used at degree \(n-1\)
shows that \(d_P^{n-1}\) surjects onto the inverse image of \(B^nX\) in
\(Z^nP\). Indeed \(X^{n-1}\twoheadrightarrow B^nX\) is an epimorphism,
and so is its pullback to that inverse image. Consequently the induced
map \(Z^nP/B^nP\to Z^nX/B^nX\) is both monic and epic. This proves it is an
isomorphism for every \(n\).

\begin{theorem}[Bounded-below injective replacement]\label{res:R01}
If \(\mathcal A\) has enough injectives, a bounded-below complex \(X\)
admits a termwise monic quasi-isomorphism \(X\to I\), with \(I\) termwise
injective and zero below the same lower bound. Every object of
\(D^+(\mathcal A)\) admits such an injective representative.
\end{theorem}
\begin{proof}
Apply the preceding construction in \(\mathcal A^{\mathrm{op}}\) to
the bounded-above complex with degree \(i\) term \((X^{-i})^{\mathrm{op}}\)
and degree \(i\) differential \((d_X^{-i-1})^{\mathrm{op}}\). This reversal
interchanges kernels and cokernels, so the comparison returns a
termwise monic quasi-isomorphism with the stated lower bound. Use
Proposition~\ref{res:D06} first if the input is an object of the derived
category. See \cite[Tag 013K]{stacks-derived}.
\end{proof}

\begin{theorem}[Bounded-above projective replacement]\label{res:R02}
If \(\mathcal A\) has enough projectives, a bounded-above complex \(X\)
admits a termwise epic quasi-isomorphism \(P\to X\), with \(P\) termwise
projective and zero above the same upper bound. Every object of
\(D^-(\mathcal A)\) admits such a projective representative.
\end{theorem}
\begin{proof}
The downward induction above constructs the epic comparison, and
Proposition~\ref{res:D06} supplies a bounded-above complex for a derived
object. The construction replaces a complex rather than only an object
in degree zero. Compare \cite[Tag 0646]{stacks-derived} and the ordinary
module resolutions of \cite[Chapter 10]{p103}.
\end{proof}

\begin{theorem}[Orthogonality of the bounded models]\label{res:R03}
A complex \(I\) is K-injective if every map to \(I\) from an acyclic
complex is null-homotopic. Dually \(P\) is K-projective if every map
from \(P\) to an acyclic complex is null-homotopic. A bounded-below
complex of injectives is K-injective, and a bounded-above complex of
projectives is K-projective.
\end{theorem}
\begin{proof}
For \(u:C\to I\) with \(C\) acyclic, construct
\(s^n:C^n\to I^{n-1}\) starting at the lower bound of \(I\). If the
homotopy identity has been established through degree \(n-1\), then
\(u^n-d_I^{n-1}s^n\) kills \(B^nC=Z^nC\) by that identity and the chain
map property of \(u\). It therefore factors through
\(\operatorname{im}d_C^n\subseteq C^{n+1}\). Extend this factorization
to \(s^{n+1}:C^{n+1}\to I^n\) by injectivity. This gives
\(u=d_Is+sd_C\) in every degree. Passing to the opposite category
and reversing cochain degrees gives the projective statement.
See \cite[Tags 013R, 0647]{stacks-derived}.
\end{proof}

\begin{theorem}[Morphisms computed on K-models]\label{res:R04}
For K-injective \(I\) and K-projective \(P\), the maps induced by
localization
\[
 \operatorname{Hom}_{\mathrm K}(X,I)\longrightarrow
     \operatorname{Hom}_{D}(QX,QI),\qquad
 \operatorname{Hom}_{\mathrm K}(P,X)\longrightarrow
     \operatorname{Hom}_{D}(QP,QX)
\]
are bijective and natural. A quasi-isomorphism between two K-injectives,
or between two K-projectives, is a homotopy equivalence.
\end{theorem}
\begin{proof}
For a quasi-isomorphism \(s:U\to X\), its cone is acyclic by
Proposition~\ref{res:K03}. Applying \(\operatorname{Hom}_{\mathrm K}(-,I)\)
to the cone triangle and using Proposition~\ref{res:K06} shows that
precomposition by \(s\) is a bijection. Thus a roof
\(X\xleftarrow{s}U\xrightarrow{u} I\) has a unique representative
\(v:X\to I\) in the homotopy category with \(vs=u\).
Refining two equal roofs leaves that representative unchanged, by the
same bijection. The right roof argument and
\(\operatorname{Hom}_{\mathrm K}(P,-)\) prove the dual assertion.
In each case an inverse to the image of a quasi-isomorphism is between
objects on which \(Q\) is fully faithful. It lifts to a two-sided inverse
in \(\mathrm K\), which is precisely the homotopy equivalence assertion.
This also supplies uniqueness up to homotopy of each resolution lift.
See \cite[Tags 05TG, 064B]{stacks-derived}.
\end{proof}

\section{Total functors and acyclic models}

Let \(F_K\) denote termwise application of \(F:\mathcal A\to\mathcal B\)
on either bounded homotopy category. It sends homotopies to homotopies.
Write \(U=Q_{\mathcal B}F_K\), so that \(U\) already takes values in the
target derived category. We use ordinary categorical composition order in
this chapter.

\begin{definition}[Total derived functor and its comparison]
\label{res:R05}
A total right derived functor comprises a functor \(RF\) on the
localized source and a natural transformation \(\eta:U\to RF Q_{\mathcal A}\).
For every functor \(H\) on the localized source, composition with \(\eta\)
identifies natural transformations \(RF\to H\) with transformations
\(U\to H Q_{\mathcal A}\). On the left, a comparison \(\epsilon:LF Q_{\mathcal A}\to U\)
identifies natural transformations \(H\to LF\) with transformations
\(H Q_{\mathcal A}\to U\).
\end{definition}

The comparisons characterize the functors up to unique compatible
isomorphism. In the right case, factor the second comparison through the
first and the first through the second. Both compositions induce the
original comparison, so both are identities by uniqueness. The left case
reverses the arrows. A natural transformation between the original
functors induces one between their right derived functors by the same
factorization. The left version uses the dual factorization.
See \cite[Tags 05SE, 05SV]{stacks-derived} and
\cite[Chapter 8]{yekutieli}.

\begin{theorem}[Total right derived functor on \(D^+\)]\label{res:R06}
With enough injectives in \(\mathcal A\), an additive \(F\) has an
exact total right derived functor \(RF:D^+(\mathcal A)\to D^+(\mathcal B)\).
It is computed by \(F(I)\) after replacement \(X\to I\) from
Theorem~\ref{res:R01}. The comparison on a model \(I\) is an isomorphism,
and natural transformations of additive functors are derived using
these same models.
\end{theorem}
\begin{proof}
The full homotopy subcategory of bounded-below complexes of injectives
is closed under shifts and cones, since their terms are finite biproducts
of injectives. Its composite with \(Q_{\mathcal A}\) is an exact
equivalence onto \(D^+(\mathcal A)\) by Theorems~\ref{res:R01}
and~\ref{res:R04}. Transport \(Q_{\mathcal B}F_K\) restricted to that
subcategory along an inverse equivalence. An additive functor
preserves the finite biproducts in the cone matrix and its signs, so the
transported functor is exact. Degreewise \(F(X\to I)\) gives \(\eta\),
which is natural because replacement lifts are unique up to homotopy
and \(F_K\) preserves homotopies.

To verify the universal property, a transformation \(U\to H Q_{\mathcal A}\)
can be restricted to the model subcategory. There \(\eta\) is an isomorphism,
so it gives the required map \(RF\to H\) via the equivalence. Naturality
on \(X\to I\) forces the original transformation on every complex,
which proves both the factorization and its uniqueness.
See \cite[Tags 05TI, 0159]{stacks-derived}.
\end{proof}

\begin{theorem}[Total left derived functor on \(D^-\)]\label{res:R07}
With enough projectives in \(\mathcal A\), an additive \(F\) has an
exact total left derived functor \(LF:D^-(\mathcal A)\to D^-(\mathcal B)\).
It is computed by \(F(P)\) after replacement \(P\to X\) from
Theorem~\ref{res:R02}.
\end{theorem}
\begin{proof}
Use the full homotopy subcategory of bounded-above complexes of
projectives. Shifts and cones preserve it. Theorems~\ref{res:R02}
and~\ref{res:R04} make its localization an exact equivalence onto
\(D^-(\mathcal A)\). Restrict \(Q_{\mathcal B}F_K\) and transport
across the equivalence. The comparison is induced by \(F(P\to X)\).
Given a transformation \(H Q_{\mathcal A}\to U\), restrict to projective
models, where this comparison is invertible, and transport back.
Naturality on \(P\to X\) shows it agrees with the given map for every \(X\).
This construction uses the same cochain shifts as the right one.
If expressed by duality, the reversal in Theorem~\ref{res:R01}
changes \([n]\) into \([-n]\) and interchanges the two bounds.
See \cite[Theorem 8.4.23]{yekutieli}.
\end{proof}

The unit and counit in these constructions commute with shifts:
the shift of a model is a model for the shifted complex, and
the resolution comparison comes from its shifted map. Uniqueness
of the comparison in the homotopy category makes this compatible
with using any choice of models.

\begin{theorem}[Computing with acyclic terms]\label{res:R08}
Call an object \(A\) right \(F\)-acyclic if the comparison
\(F(A)[0]\to RF(A[0])\) is an isomorphism, and define left
\(F\)-acyclicity using the reverse comparison to \(F(A)[0]\).
For a bounded-below complex \(C\) of right \(F\)-acyclic objects,
\(F(C)\to RF(C)\) is an isomorphism. For a bounded-above complex of left
\(F\)-acyclic objects, \(LF(C)\to F(C)\) is an isomorphism.
If \(F\) is left exact, right acyclicity reduces to the vanishing of the
positive right derived functors. For a right exact \(F\), the corresponding
left condition is the vanishing of the positive left derived functors.
\end{theorem}
\begin{proof}
In the right case begin with a finite complex. The degreewise split
sequence peeling off its bottom term gives a triangle by
Theorem~\ref{res:D04}. Application of \(F\) preserves this split sequence
degreewise. Naturality of the comparison gives a map of triangles,
and the long exact sequences show inductively that it is an isomorphism
on cohomology, starting with a shifted stalk.
Now fix a degree \(i\) and replace \(C\) by its brutal quotient truncation
in degrees \(\le i+1\). Its kernel has no terms below \(i+2\).
Both that kernel after \(F\) and its right derived image have
vanishing cohomology in degrees \(i,i+1\). For the derived image this
uses the bound-preserving replacement of Theorem~\ref{res:R01}.
The long exact sequences thus reduce comparison in degree \(i\)
to the already proved finite case.

For the left case, reverse the degrees and use the opposite categories.
The top-term sequence now peels off a top stalk, and the discarded
degrees are below those being tested. For a left exact \(F\), applying it
to the first two steps of an injective resolution identifies \(F(A)\)
with the zeroth cohomology of \(F(I)\). The replacement starts in degree
zero, so all negative cohomology vanishes. The right exact case uses the
first two projectives and cokernels. See
\cite[Tags 015E, 05TA]{stacks-derived}.
\end{proof}

\begin{theorem}[Classical derived functors and the delta-functor property]
\label{res:R09}
For a stalk \(A[0]\), the groups or objects
\(R^iF(A)=H^i(RF(A[0]))\) and \(L_iF(A)=H^{-i}(LF(A[0]))\)
are computed by ordinary injective and projective resolutions.
They vanish for \(i>0\) on injective and projective objects,
respectively. For left exact \(F\), \(R^0F=F\), and the right derived
functors form a universal cohomological delta-functor. For right exact
\(F\), \(L_0F=F\) and the left derived functors have the homological
analogue.
\end{theorem}
\begin{proof}
The replacements of a stalk are the ordinary resolutions, with a
projective resolution indexed in nonpositive cochain degrees. An
injective or projective stalk is itself an allowed model. The degree-zero
claims are those in Theorem~\ref{res:R08}. For a short exact sequence
of stalks use Theorem~\ref{res:D04} and apply the exact total functor.
The cohomology sequences give the connecting maps.

For universality on the right, for each object \(A\) embed it in an
injective \(I\) with quotient \(B\). The connecting map
\(R^{i-1}F(B)\to R^iF(A)\) is epic for \(i>0\), and its kernel is the
image from \(R^{i-1}F(I)\). Given a degree-zero transformation to
another delta-functor, construct its higher components using this
surjection and the corresponding connecting map of the target.
The previous degree of the transformation respects the image from \(I\),
so the construction factors through the surjection. For a morphism
\(A\to A'\), extend its composite \(A\to I'\) to \(I\to I'\).
Naturality of the two long exact sequences, and the surjection just
used, prove naturality and independence of the embedding. Compatibility
with connecting maps for an arbitrary short exact sequence follows by
mapping that sequence into \(0\to A\to I\to B\to0\), extending \(A\to I\)
over its middle object. Surjectivity also proves uniqueness.
Dually, given a degree-zero map from a homological delta-functor into
\(L_0F\), choose a projective epimorphism \(P\to A\) with kernel \(J\).
The inclusion \(L_iF(A)\to L_{i-1}F(J)\) for \(i>0\) forces the
higher maps by the same argument in the opposite category. See
\cite[Tags 015A, 05TE, 015B]{stacks-derived}.
\end{proof}

\begin{proposition}[Exact functors and localization]\label{res:R10}
An exact \(F\) descends termwise to an exact functor \(D(\mathcal A)\to D(\mathcal B)\)
commuting with cohomology. On the bounded derived categories on which
they exist, \(RF\) and \(LF\) identify with this functor.
\end{proposition}
\begin{proof}
Exactness identifies the kernels, images, and quotients defining
cohomology of \(F(X)\) with the \(F\)-images of those for \(X\).
It follows that \(F_K\) preserves quasi-isomorphisms, so the localization
property in Theorem~\ref{res:D03} gives the functor. Because
termwise application preserves cone matrices and shifts, the functor is
exact. Its comparison before localization is invertible, so it satisfies
both universal properties in Definition~\ref{res:R05} whenever
restricted to the relevant bounded category. See
\cite[Tag 015F]{stacks-derived}.
\end{proof}

\section{Ext, Tor, and comparison of compositions}

\begin{theorem}[Ext as shifted Hom]\label{res:R11}
For \(A,B\in\mathcal A\) the group
\(\operatorname{Hom}_{D(\mathcal A)}(A[0],B[n])\) agrees for \(n\ge0\)
with Ext computed by either of the available bounded resolutions,
covariantly in \(B\) and contravariantly in \(A\).
In degree zero it is \(\operatorname{Hom}_{\mathcal A}(A,B)\), and
it vanishes for \(n<0\).
\end{theorem}
\begin{proof}
The fully faithful heart embedding of Theorem~\ref{res:D07} gives
the zeroth formula. Its orthogonality gives the negative vanishing,
whether or not either kind of resolution is available.
For an injective resolution \(B\to I\), Theorem~\ref{res:R04}
computes the morphism group by homotopy classes of maps \(A[0]\to I[n]\).
Those maps are morphisms \(A\to I^n\) killed by the next differential,
modulo morphisms factoring through the previous differential via
\(A\to I^{n-1}\). The signs of \(I[n]\) do not change either subgroup.
This gives \(H^n\operatorname{Hom}(A,I^\bullet)\), the stated vanishing
and the zeroth identification. A projective resolution \(P\to A\) gives
the same argument for maps \(P\to B[n]\). In the usual cochain Hom
convention its differential on degree \(j\) is
\(d(u)=d_B u-(-1)^j u d_P\), so cycles are precisely maps to \(B[j]\);
boundaries are precisely homotopies to zero. If one writes the
resolution Hom differential as unsigned precomposition, multiply
its degree-\(j\) term by \((-1)^{j(j+1)/2}\) to identify the two
conventions. The Hom bijections are natural and lifts unique up to
homotopy, giving the bifunctor comparison. This identifies, in
particular, the resolution-based module Ext of \cite[Chapter 11]{p103}
with the derived model without identifying their representing complexes.
See \cite[Tags 06XQ--06XS]{stacks-derived}.
\end{proof}

\begin{theorem}[Right derived composition and acyclicity]\label{res:R12}
For left exact \(F:\mathcal A\to\mathcal B\) and \(G:\mathcal B\to\mathcal C\),
assume \(\mathcal A,\mathcal B\) have enough injectives. The canonical map
\(R(GF)\to RG\,RF\) on \(D^+(\mathcal A)\) is invertible if and only if
\(F(I)\) is right \(G\)-acyclic for every injective \(I\) of \(\mathcal A\).
\end{theorem}
\begin{proof}
Compose the comparison from \(G\) on the termwise complex \(F_K(X)\)
with \(RG\) of the comparison from \(F\) on \(X\). This is a
transformation \(Q_{\mathcal C}(GF)_K\to RG\,RF Q_{\mathcal A}\);
Definition~\ref{res:R05} induces the displayed map. Evaluate it on a
bounded-below injective model \(I^\bullet\) of \(X\). There its first
term is \(GF(I^\bullet)\), the intermediate comparison for \(F\)
is invertible, and it is exactly the comparison
\(G(F(I^\bullet))\to RG(F(I^\bullet))\).
The latter is invertible under the acyclicity condition by
Theorem~\ref{res:R08}. Conversely, evaluate an invertible comparison
at each injective stalk \(I[0]\). Its component is then precisely
\(G(F(I))[0]\to RG(F(I)[0])\). This is right \(G\)-acyclicity by
definition. See \cite[Tags 05T2, 015M]{stacks-derived}.
\end{proof}

\begin{proposition}[Ext boundaries and extension classes]\label{res:R13}
For a short exact sequence \(0\to A\to B\to C\to0\), write its derived
boundary as \(\delta:C\to A[1]\). The triangle map in the second Ext
variable is postcomposition with \(\delta[n]\); in the first it takes
\(u\) to \(u[1]\delta\). Under the injective Hom comparison of
Theorem~\ref{res:R11}, the ordinary second-variable boundary has the
displayed sign, and the ordinary first-variable boundary takes \(u\)
to \((-1)^{n+1}u[1]\delta\).
With extension classes given by the triangle convention of
Theorem~\ref{res:D04}, the assignment of \(\delta\) to a short exact
sequence identifies Baer sum with addition in \(\operatorname{Ext}^1(C,A)\).
\end{proposition}
\begin{proof}
The two representable long exact sequences of the derived triangle
give these maps, by the argument of Proposition~\ref{res:K06}.
For the second variable, regard \(A\to B\to C\) as a
three-term acyclic complex. Apply the Cartan--Eilenberg
construction of Theorem~\ref{res:G01}. Each successive
cokernel complex is acyclic by Lemma~\ref{res:G07}, so the
injective object for horizontal cohomology is always zero.
Each horizontal row is therefore a split short exact
sequence, and the three columns resolve \(A,B,C\).
This gives the injective horseshoe in this abelian category.
If its off-diagonal differential is
\(\alpha:I_C^j\to I_A^{j+1}\), the connecting class on a cocycle \(z\)
is \(\alpha z\), by taking its lift in the \(I_C\)-summand. The
map \(\delta\) is represented by \(\alpha:I_C^\bullet\to I_A^\bullet[1]\)
by Theorem~\ref{res:D04}. Thus postcomposition with \(\delta[n]\)
has the claimed sign.

For the first variable, use an injective resolution \(I^\bullet\) of
the fixed target, represent \(u\) by a cocycle \(u:A\to I^n\), and
extend it to \(w:B\to I^n\). Write \(d_Iw=v\pi\) for the quotient
\(\pi:B\to C\), so the usual Hom connecting map is represented by \(v\).
On the two-term cone \((A\xrightarrow{i}B)\) in degrees \(-1,0\),
the map \(u[1]\delta\) has components \((-u,0)\).
Adding the homotopy with degree-zero component \(w\) gives components
\((0,(-1)^{n+1}v\pi)\): the degree-zero target differential is that
of \(I[n+1]\), and the degree \(-1\) homotopy term is \(wi=u\).
This proves the comparison, with the sign shown in the statement.

Here is the extension classification in any abelian category.
By Theorem~\ref{res:D02} a morphism \(C[0]\to A[1]\) is represented
by a roof \(C[0]\leftarrow L^\bullet\to A[1]\) whose left arrow is
a quasi-isomorphism. Replace \(L\) first by \(\tau_{\le0}L\) and then by
its good quotient truncation \(\tau_{\ge-1}\). Both replacements
are quasi-isomorphisms and the arrows restrict and descend. For the
arrow to \(A[1]\), descent in degree \(-1\) follows since it kills
the boundaries there. The roof is now given by a short exact sequence
\(0\to L^{-1}\xrightarrow{d}L^0\to C\to0\) and a map \(a:L^{-1}\to A\).
Push this sequence out along \(-a\). The resulting cone roof maps
to \(A[1]\) by the negative projection, and hence represents the
given morphism. A homotopy change of \(a\) adds a map factoring
through \(d\). Pushing out along maps differing this way gives
equivalent extensions: in the biproduct presentation of the pushout,
add the corresponding multiple of the \(L^0\)-coordinate to the
\(A\)-coordinate. If two roofs agree, pass to a common refinement
as in Theorem~\ref{res:D02} and apply the same two truncations.
The left maps agree on cycles in degree zero, and the right maps
differ at most by the homotopy term just treated. The refined
extension therefore maps to both pushouts by maps inducing the
identity on \(A\) and \(C\). Conversely an equivalence of extensions
gives equal cone roofs. This proves the bijection.

Taking direct sums of the cone roofs, pulling back by the diagonal
of \(C\), and pushing out by the codiagonal of \(A\) sends two
boundaries to their sum. This is the Baer addition operation.
The projection convention of \cite[Tags 06XP, 06XU, 06XV]{stacks-derived}
uses the positive projection for the Yoneda class; with our cone
matrix it gives the negative of the triangle boundary. The
classification and addition statements agree after this uniform
sign change. See also \cite[Chapter 11]{p103}.
\end{proof}

\begin{proposition}[Classical Tor from a total left derived functor]
\label{res:R14}
Fix a ring \(R\) and a left \(R\)-module \(M\). For a right
\(R\)-module \(A\), \(H^{-n}\bigl(L(-\otimes_R M)(A[0])\bigr)\)
is naturally \(\operatorname{Tor}_n^R(A,M)\). In degree zero this is
\(A\otimes_R M\), and its connecting maps agree with the classical
ones with the chain indices reversed from cochain indices.
\end{proposition}
\begin{proof}
Choose a projective right-module chain resolution \(P_\bullet\to A\),
place \(P_n\) in degree \(-n\) with differential \(d_n:P_n\to P_{n-1}\),
and apply Theorem~\ref{res:R07}. On this model both constructions
take the homology of \(P_\bullet\otimes_R M\). Resolution comparison
makes the identification independent of the choice and natural in both
module variables. The balanced computation via a left-module resolution
is the comparison of Theorem~\ref{res:B05}. Right exactness gives
the degree-zero statement. For a short exact sequence of right modules
use a compatible horseshoe resolution. Its degreewise split sequence
remains exact after tensoring and gives both connecting maps by the
same lift and differential. The sign of the map in the derived
triangle is the one identified in Theorem~\ref{res:D04}. See
\cite[Section 3.2]{weibel} and \cite[Chapter 11]{p103}.
\end{proof}

\Needspace{24\baselineskip}
\section{Higher Yoneda extensions and products}

For objects \(X,Y\) of the abelian category and \(n\ge1\), an
\(n\)-extension of \(X\) by \(Y\) is an exact sequence
\[
  0\longrightarrow Y\longrightarrow E_{n-1}\longrightarrow
  \cdots\longrightarrow E_0\longrightarrow X\longrightarrow 0 .
\]
Take equivalence generated by maps of such exact sequences that
are the identity on both endpoints. The Baer operation is direct
sum followed by pullback at the right endpoint along the diagonal
of \(X\) and pushout at the left endpoint along the codiagonal of
\(Y\). For \(n=0\) use \(\operatorname{Hom}(X,Y)\). Splicing an
\(m\)-extension of \(Y\) by \(Z\) on the left of an \(n\)-extension
of \(X\) by \(Y\) uses the composite through \(Y\) at their junction.

\begin{theorem}[The higher Yoneda product]\label{res:R15}
In any abelian category the equivalence classes of \(n\)-extensions
of \(X\) by \(Y\), with the Baer operation, are naturally isomorphic
to \(\operatorname{Ext}^n(X,Y)\) as defined in
\ref{res:R11}. In degree one the isomorphism uses the positive
projection convention of \ref{res:R13}. Under these isomorphisms,
splicing a class \(\alpha\) of degree \(m\) on the left of a class
\(\beta\) of degree \(n\) is
\[
       \alpha[n]\circ\beta:X\longrightarrow Z[n+m],
       \qquad
       \beta:X\longrightarrow Y[n],\quad
       \alpha:Y\longrightarrow Z[m].
\]
Degree-zero classes act by pullback or pushout on extensions and by
ordinary composition in degree zero.
\end{theorem}

\begin{proof}
Put the terms from \(Y\) through \(E_0\) in degrees \(-n,\ldots,0\)
with the displayed differentials, forming a cochain complex \(C_E\).
The augmentation \(s_E:C_E\to X[0]\) is a quasi-isomorphism.
Projection \(p_E:C_E\to Y[n]\) is a chain map, equal to the identity
on \(Y\) in degree \(-n\). Write \(u_E=p_Es_E^{-1}\) for the
unscaled roof and put
\[
             y_E=(-1)^{n(n-1)/2}u_E .
\]
These expressions are invariant under maps of extensions with
identity endpoints, since the induced chain maps commute with
augmentation and projection and are quasi-isomorphisms.

We verify that the unscaled roofs already give a bijection of sets.
Represent a morphism \(X[0]\to Y[n]\) by a roof with middle complex
\(L\), using \ref{res:D02}. Apply the good truncations of
\ref{res:D05}, first \(\tau_{\le0}\), then \(\tau_{\ge -n}\).
They are quasi-isomorphisms in the relevant directions because
\(H^i(L)=0\) for \(i\ne0\). The arrow to \(Y[n]\) descends in
degree \(-n\) because it kills the incoming boundaries. Write
\(U\) for this truncated complex, \(s^0:U^0\to X\) for the
augmentation, and \(f:U^{-n}\to Y\) for the other arrow.
Then \(U^{-n}\to U^{-n+1}\) is monic, \(s^0\) is epic,
and the augmented row is exact. Push out the first monomorphism
along \(f\). Pushouts of monomorphisms in an abelian category
preserve the short exact row with the same cokernel, so the
resulting row is an \(n\)-extension by \(Y\). The natural map
from \(U\) to its complex is \(f\) at the left, the pushout map
next, and the identity in the remaining terms. It commutes with
both arrows of the roof and induces the identity on \(X\).
This proves surjectivity.

If two extension roofs agree, use the common refinement of
\ref{res:D02} and apply the same two truncations to it.
The maps to the extension complexes descend through the lower
truncation because their targets vanish below degree \(-n\).
The two maps to \(X[0]\) become identical: a homotopy has no contribution
in degree zero once the source is supported in degrees at most
zero. The two maps to \(Y[n]\) then differ in degree \(-n\) by a
map \(h d_U^{-n}\) for some \(h:U^{-n+1}\to Y\), by the chain
homotopy formula. The two pushouts along \(f\) and \(f+h d_U^{-n}\)
are isomorphic by the biproduct map
\((y,v)\mapsto(y-h(v),v)\). Indeed, it maps the defining relations
\((f(a),-d_U(a))\) to \((f(a)+h d_U(a),-d_U(a))\), and fixes
the endpoint maps. Each truncated refinement induces a map
from its pushout extension to the corresponding original
extension, fixing \(X\) and \(Y\). This gives the required
equivalence and proves injectivity.

Endpoint maps act on the roofs by precomposition and postcomposition:
pullback changes the augmentation by its pullback square, and
pushout changes the projection by its pushout square. The induced
maps of extension complexes give those two identities of roofs.
Taking direct sums and then the diagonal and codiagonal proves
additivity. Multiplying by the chosen sign preserves these
properties. In degree one the factor is \(1\), and \(p_Es_E^{-1}\)
is the positive projection of \ref{res:R13}, namely the negative
of the triangle boundary of \ref{res:D04}.

For the product, take \(n,m>0\). Let \(C\) be the complex of the
spliced extension of \(\alpha\) on the left of \(\beta\).
It maps to \(C_\beta\) by the identity on the \(\beta\)-terms,
by the augmentation of the rightmost \(\alpha\)-term onto \(Y\)
in degree \(-n\), and by zero to the left. Call this map \(a\).
It preserves the augmentation to \(X\). There is also a chain
map \(b:C\to C_\alpha[n]\) which is zero to the right of
degree \(-n\) and is multiplication by \((-1)^{nj}\) on the
\(\alpha\)-term in degree \(-n-j\) for \(0\le j\le m\).
The differential in \(C_\alpha[n]\) has the factor \((-1)^n\),
so these scalars commute with the differential. At degree \(-n\)
the scalar is \(1\), whence \(s_\alpha[n]b=p_\beta a\).
Since \(s_\beta a\) is the augmentation of the spliced extension,
the composite of the unscaled roofs is represented on \(C\) by
\(p_\alpha[n]b\). This equals \((-1)^{nm}\) times the positive
projection of \(C\) onto \(Z[n+m]\).

Finally,
\[
 \frac{(n+m)(n+m-1)}{2}
       =\frac{n(n-1)}{2}+\frac{m(m-1)}{2}+nm .
\]
It follows that the normalized classes give exactly shifted
composition. The cases with a degree-zero factor follow from
naturality. Thus the sign comparison uses the cochain shift
and the degree-one projection throughout.
For the extension classification see
\cite[Tags 06XP, 06XT and 06XU]{stacks-derived}.
The calculation above specifies the normalization for our
cochain shifts and extension differentials.
\end{proof}

\chapter[Cartan--Eilenberg Resolutions and Grothendieck Spectral Sequences]%
  {Cartan--Eilenberg Resolutions\\and Grothendieck Spectral Sequences}
\label{ch:grothendieck}

An injective model computes a total right derived functor. To compute its
cohomology in stages, we resolve each term of a complex while also resolving
its cycles, boundaries, and cohomology. The two filtrations of the resulting
double complex then give distinct first pages. Throughout the right-derived
construction, \(p\) is horizontal cochain degree and \(q\) is vertical
resolution degree. The unsigned horizontal and vertical differentials
commute, and the total differential is \(d_h+(-1)^p d_v\).

\section{Simultaneous resolutions}

For a cochain complex \(K^\bullet\) in an abelian category, write
\[
 B^pK=\operatorname{im}(K^{p-1}\to K^p),\qquad
 Z^pK=\ker(K^p\to K^{p+1}),\qquad
 H^pK=Z^pK/B^pK .
\]
A Cartan--Eilenberg injective resolution is a double complex \(I^{p,q}\)
with \(q\ge0\) and \(p\) bounded below, augmented by a map
\(K^\bullet\to I^{\bullet,0}\) killed by \(d_v\). Its columns
\(I^{p,\bullet}\) and the columns \(B_h^p I\), \(Z_h^p I\), \(H_h^p I\)
obtained by taking horizontal boundaries, cycles, and cohomology are
injective resolutions of \(K^p,B^pK,Z^pK,H^pK\), respectively. Each
augmentation is the one induced by the given map from \(K^\bullet\).
The signed totalization uses the convention of Proposition~\ref{res:B01}.
This definition is \cite[Tag 015H]{stacks-derived}.

\begin{lemma}[Exactness for horizontal invariants]
\label{res:G07}
Suppose \(i:U^\bullet\to V^\bullet\) is injective in each degree and that
\[
      U^p/B^pU\longrightarrow V^p/B^pV
\]
is injective for all \(p\). Put \(W^\bullet=\operatorname{coker}(i)\).
Taking boundaries, cycles, cohomology, or quotients by boundaries of
\(0\to U^\bullet\to V^\bullet\to W^\bullet\to0\) gives a short exact
sequence in each degree. In this chapter an injection satisfying
these two tests is called admissible.
\end{lemma}

\begin{proof}
Identify \(U^\bullet\) with its image and write \(q:V^\bullet\to W^\bullet\)
for the quotient. Generalized elements and lifts along epimorphisms
are used as in the paragraph on subobjects in Chapter~\ref{ch:filtrations}.
A lift can be taken after epic pullback. Image membership tested after
such a pullback implies image membership before it.

The image of \(B^pV\) under \(q\) is \(B^pW\): a representative of
a boundary in \(W^p\) is a differential of a locally lifted
predecessor in \(W^{p-1}\), and that predecessor in turn lifts to
\(V^{p-1}\). The kernel on \(B^pV\) is \(U^p\cap B^pV\). This
intersection equals \(B^pU\) by the injectivity hypothesis on
quotients. Thus boundaries form a short exact sequence.

For the quotient sequence, the first map is monic by hypothesis.
The last map is epic since the composite \(V^p\to W^p\to W^p/B^pW\)
is epic. If a class in \(V^p/B^pV\) maps to zero, lift it to a
representative \(v\) in \(V^p\). Its image \(qv\) belongs to \(B^pW\).
Lift \(qv\) to a boundary \(b\) in \(B^pV\); then \(v-b\) belongs
to \(U^p\) and represents the original class. The composite from
\(U^p/B^pU\) is zero, so this gives precisely the kernel.

For cycles, a lift \(v\in V^p\) of a cycle in \(W^p\) has
\(d_Vv\in U^{p+1}\cap B^{p+1}V=B^{p+1}U\). Lift that boundary to
\(u\in U^p\) with \(d_Uu=d_Vv\). Then \(v-u\) is a cycle lifting
the given one. The kernel on cycles consists of cycles in the
subobject \(U^p\), and they are precisely \(Z^pU\) because the
inclusion in degree \(p+1\) is monic. The first map on cycles is
monic because the inclusion in degree \(p\) is monic.

Finally \(H^pU\to H^pV\) is monic by the already proved quotient
injection, and \(H^pV\to H^pW\) is epic by the cycle lift. If a
cohomology class maps to zero, lift it to a cycle \(v\) in \(V^p\).
Then \(qv\) is a boundary and lifts to a boundary \(b\) in \(V^p\).
The difference \(v-b\) belongs to \(Z^pU\) and gives the original
class. Conversely classes coming from \(U^\bullet\) map to zero.
All maps used are restrictions or induced quotient maps of the
given complex maps, so the short exact sequences have the asserted
natural maps.
\end{proof}

\begin{theorem}[Existence of a Cartan--Eilenberg resolution]
\label{res:G01}
In an abelian category with enough injectives, every bounded-below
cochain complex admits a Cartan--Eilenberg injective resolution.
If \(K^p=0\) for \(p<L\), the resolution can be chosen with
\(I^{p,q}=0\) for \(p<L\).
\end{theorem}

\begin{proof}
Choose injective objects \(E_B^p,E_H^p\) and monomorphisms
\(j_B^p:B^pK\to E_B^p\) and \(j_H^p:H^pK\to E_H^p\).
Choose the target to be zero when the source is zero. Form the complex
\[
   J^p=E_B^p\oplus E_H^p\oplus E_B^{p+1},
   \qquad d_J^p(b,h,b')=(b',0,0).
\]
The first summand is the boundary object and the first two summands
form the cycle object in degree \(p\). By injectivity extend
\(j_B^p\) from \(B^pK\) to a map \(u^p:K^p\to E_B^p\). Also extend
the composite \(Z^pK\twoheadrightarrow H^pK\xrightarrow{j_H^p}E_H^p\)
to a map \(v^p:K^p\to E_H^p\). The third component \(w^p\) is
\(K^p\twoheadrightarrow B^{p+1}K\xrightarrow{j_B^{p+1}}E_B^{p+1}\).
The map \((u^p,v^p,w^p)\) is a chain map. On the image of \(d_K^p\),
the first component in degree \(p+1\) is \(j_B^{p+1}\), while the
second and third vanish. On applying \(d_J^p\), the third component
of the input becomes exactly this first component.

On \(K^p/B^pK\), the last two components give a map to
\(E_H^p\oplus E_B^{p+1}\). On the subobject \(H^pK\) the first of
these components is \(j_H^p\), and on the quotient \(B^{p+1}K\) the
second is \(j_B^{p+1}\). Consequently the map on \(K^p/B^pK\) is
injective by the short exact sequences of these subobjects and
quotients. The map on \(K^p\) is injective by the same argument with
subobject \(B^pK\). We have constructed an admissible injection
\(K^\bullet\to J^\bullet\). It respects both short exact sequences
\[
 0\to B^pK\to Z^pK\to H^pK\to0,\qquad
 0\to Z^pK\to K^p\to B^{p+1}K\to0 .
\]
In \(J\) these become split sequences. This is the simultaneous
injective horseshoe construction for the two sequences, with the
same embedded boundary used on both sides of each differential.

Set \(K_0=K\). Iterate the construction with \(K_q\to J_q\)
and \(K_{q+1}=\operatorname{coker}(K_q\to J_q)\). The vertical
maps of \(I^{p,q}=J_q^p\) are induced by the quotient and embedding
\(J_q\twoheadrightarrow K_{q+1}\hookrightarrow J_{q+1}\). These
are chain maps for \(d_h\), and successive vertical maps compose
to zero. The construction gives short exact sequences with terms
\(K_q,J_q,K_{q+1}\) on horizontal terms and on all three horizontal
invariants by Lemma~\ref{res:G07}. Thus the augmented
vertical complexes are exact, and their terms in nonnegative
vertical degrees are injective finite sums or direct summands of
injectives. This proves all the simultaneous resolution conditions.
For \(p<L\) the first two possible summands of \(J^p\) vanish.
The third vanishes too because \(B^{p+1}K=0\) for \(p+1\le L\).
The same lower bound holds for the quotient at every step.
Compare \cite[Tag 015I]{stacks-derived} and the projective
construction in \cite[Section 5.7]{weibel}.
\end{proof}

\section{Comparison and totalization}

\begin{lemma}[Extension into a split row]
\label{res:G08}
A horizontal row of a Cartan--Eilenberg injective resolution has the
extension property along admissible injections of complexes.
Moreover, in the augmented vertical resolution, the inclusion of
each vertical kernel into its row is admissible.
\end{lemma}

\begin{proof}
In a fixed row, the short exact sequences for cycles and boundaries
split because the subobjects \(B_h^p I\) and \(Z_h^p I\) are injective.
With a choice of splitting the row has the form used in
Theorem~\ref{res:G01}. It is the degreewise product of stalks on
injective cohomology objects and two-term complexes on injective
boundary objects, with identity as the only nonzero differential in
each two-term complex. Only three factors can contribute in a given
degree, so the product exists in the category of complexes using
only finite products in the abelian category. A map to a stalk of
an injective \(E\) in degree \(p\) is a map \(U^p/B^pU\to E\). A
map to the two-term identity complex in degrees \(p-1,p\) is
determined by its degree \(p\) component \(U^p\to E\), with its
degree \(p-1\) component obtained by precomposing with \(d_U^{p-1}\).
Both components extend along the corresponding monomorphisms into
\(V\). Taking the product of these extensions proves the claim.

For a vertical degree \(q\ge0\), put
\(T_q=\ker(d_v:I^{\bullet,q}\to I^{\bullet,q+1})\), so \(T_q\)
is also the previous vertical image. At \(q=0\) this uses the
augmentation from \(K\). Exactness of the augmented boundary
resolutions shows that
\[
  T_q^p\cap B_h^p(I^{\bullet,q})
    = \operatorname{im}\bigl(d_v:B_h^p(I^{\bullet,q-1})
                           \to B_h^p(I^{\bullet,q})\bigr)
    = B^p T_q ,
\]
with \(B^pK\) and its augmentation in place of the previous row at
\(q=0\). For the last equality, take the horizontal boundary of
the image of the previous vertical map and use that its
horizontal differential is induced from the previous row.
This equality of subobjects says exactly that
\(T_q^p/B^pT_q\to I^{p,q}/B_h^p(I^{\bullet,q})\) is injective.
The inclusion is thus admissible.
\end{proof}

\begin{theorem}[Comparison and injective totalization]
\label{res:G02}
Let \(I\) and \(J\) be Cartan--Eilenberg resolutions of bounded-below
complexes \(K\) and \(L\). Every map \(f:K\to L\) lifts to a map
of double complexes \(I\to J\). Any two lifts of \(f\) are
vertically homotopic through maps of horizontal complexes.
The comparison of total complexes is compatible with the
augmentations, which are quasi-isomorphisms. Each total complex is
bounded below and termwise injective. For an additive functor \(U\)
between abelian categories, there is a natural chain isomorphism
\[
        U(\operatorname{Tot}I)\simeq\operatorname{Tot}(U I)
\]
preserving the two split filtrations from the diagonal sums.
\end{theorem}

\begin{proof}
First extend the composite \(K\to L\to J^{\bullet,0}\)
along \(K\hookrightarrow I^{\bullet,0}\) using
Lemma~\ref{res:G08}. Given the comparison
on row \(q\) preserving the preceding vertical image, it induces
a map on the quotient of that row by its vertical kernel. The
quotient is the next vertical kernel. Compose its induced map
with the inclusion into the next row of \(J\) and extend along
the admissible inclusion into the next row of \(I\). This constructs
the lift with all vertical squares commuting. Each extension is
a map of horizontal complexes.

Let \(a,b:I\to J\) be two lifts. We construct horizontal chain
maps \(s^q:I^{\bullet,q}\to J^{\bullet,q-1}\) with \(s^0=0\).
In degree zero, \(a^0-b^0\) vanishes on \(K\), so it factors
through the first vertical image of \(I\). Extend this factor to
the whole first row with target \(J^{\bullet,0}\) to obtain \(s^1\).
At the next step consider
\[
      a^q-b^q-d_v^J s^q .
\]
Its composite with \(d_v^I\) from row \(q-1\) is zero. Indeed,
commutation with \(d_v\) and the preceding step make this
composite \(d_v^Jd_v^J s^{q-1}\). The residual map therefore
factors through the next vertical image, and extends to a map
\(s^{q+1}\) by the extension property. This gives
\(a-b=d_v^J s+s\,d_v^I\). On total complexes the homotopy sends
the component of bidegree \((p,q)\) to \((-1)^p s^q\).
In \(D_J((-1)^p s)+((-1)^p s)D_I\) the horizontal terms cancel
since \(s\) commutes with \(d_h\) and a horizontal step changes
the sign. The two vertical terms have the product of two
factors \((-1)^p\), leaving \(d_v^J s+s\,d_v^I\).

Every vertical resolution has cohomology only in degree zero
with the given augmentation from \(K^p\) or \(L^p\). Its total
augmentation is a quasi-isomorphism by Theorem~\ref{res:B04}.
In degree \(n\) the total complex has summands only for
\(L_0\le p\le n\) and \(q=n-p\), for a fixed lower bound \(L_0\).
Thus finite direct sums of injectives give degreewise injectivity.
For two resolutions of the same complex, lifts of the identity
in both directions have composites vertically homotopic to the
identities by the preceding uniqueness proof. Finally an additive
functor preserves finite biproducts and signs. The canonical
biproduct isomorphism commutes with both component differentials
and the inclusions defining the row and column filtrations. This
proves the assertion about \(U\). See \cite[Tags 015J and 015K]{stacks-derived}.
\end{proof}

\section{The two hypercohomology spectral sequences}

\begin{theorem}[Two filtrations for a total right derived functor]
\label{res:G03}
Let \(F:\mathcal A\to\mathcal B\) be left exact between abelian
categories, with \(\mathcal A\) having enough injectives. For a bounded-below
complex \(K^\bullet\), there are spectral sequences with
\[
 {}_{\mathrm c}E_1^{p,q}=R^qF(K^p),
       \qquad
 {}_{\mathrm r}E_2^{p,q}=R^pF(H^qK).
\]
Both converge to \(H^{p+q}(RF(K))\) with finite filtrations in
each target degree. On the first page of the column sequence,
\(d_1\) is induced by \(d_K\). The column sequence is natural in
chain maps from \(E_1\) onward. The row sequence is natural in
chain maps from \(E_2\) onward and descends from that page,
together with its filtered target, to \(D^+(\mathcal A)\).
\end{theorem}

\begin{proof}
Choose a Cartan--Eilenberg resolution \(K\to I\). The
quasi-isomorphism \(K\to\operatorname{Tot}I\) is a bounded-below
injective replacement by Theorem~\ref{res:G02}. Thus
\(\operatorname{Tot}(FI)\simeq F(\operatorname{Tot}I)\) computes
\(RF(K)\) by Theorem~\ref{res:R06}. The double complex has finite
diagonals and the same lower bound. Apply Theorems~\ref{res:S01}
and~\ref{res:B04}, and transport the filtrations on total
cohomology via this comparison.

In the column filtration, Proposition~\ref{res:B03} identifies
the first page in bidegree \((p,q)\) with the cohomology of
\(F(I^{p,\bullet})\) in degree \(q\). The augmented column is an
injective resolution of \(K^p\), so this is \(R^qF(K^p)\) by
Theorem~\ref{res:R09}. The unsigned horizontal differential
is a chain map of these resolutions lifting \(d_K^p\), so the
induced first differential is \(R^qF(d_K^p)\).

For the row filtration, in every fixed vertical degree the
horizontal row is split into stalks and identity pairs as in
Lemma~\ref{res:G08}. Additivity of \(F\)
preserves the two splitting sequences and their differentials.
Therefore
\[
 {}_{\mathrm r}E_1^{p,q}
    \simeq H_h^q(FI^{\bullet,p})
    \simeq F(H_h^q I^{\bullet,p}).
\]
The second identification can be defined by the maps on the
split kernel, image, and quotient. It is natural in horizontal
chain maps and independent of a splitting. In the signed
interchange coordinates of Proposition~\ref{res:B03}, the
first differential is \(F\) of the induced vertical map on
\(H_h^q I^{\bullet,p}\). In the original coordinates it has the
additional factor \((-1)^q\). The vertical cohomology of
\(F(H_h^q I^{\bullet,*})\) computes \(R^pF(H^qK)\) by the
simultaneous resolution condition. This gives the row second page.

For a chain map, take a lift between the double resolutions.
It preserves each filtration because it preserves both indices.
Two lifts differ by a vertical homotopy commuting with the
horizontal differentials. In each column this is a homotopy
after applying \(F\), so the induced column map on \(E_1\) is
independent of the lift. On the horizontal cohomology of each
row it gives a homotopy in the vertical direction and hence the
same row map on \(E_2\). The maps on later pages are induced
iteratively by homology. The total homotopy gives equality on
underlying cohomology. Both lifted maps preserve the filtrations,
so their restrictions to the filtration steps coincide as well.
Lifting identities and composites then proves naturality.
In particular, different choices of the double resolution
are compared by lifts of the identity, inverse on the pages
in question and on both filtered targets.

For the row sequence a quasi-isomorphism of bounded-below complexes
induces isomorphisms on \(H^q K\) and therefore on \(E_2\).
Finite convergence and Corollary~\ref{res:S03} give an isomorphism
on the filtered target. Thus the row construction inverts
quasi-isomorphisms on the stated pages and admits a unique
descended functor by the localization universal property.
For compatibility with homotopy classes, use the cylinder
\(K\oplus\operatorname{Cone}(\mathrm{id}_K)\), whose projection to
\(K\) is a quasi-isomorphism by Proposition~\ref{res:K03}. With cone
coordinates \((a,b)\) and differential \((da+b,-db)\), its two
inclusions send \(x\) to \((x,(0,0))\) and \((x,(x,0))\).
They induce the same row maps because each is a section of
the projection. If \(g-f=dh+hd\), the map
\((x,(a,b))\mapsto f(x)+(g-f)(a)+h(b)\) is a chain map to the
common target of \(f,g\). The homotopy identity checks the
chain map property on the second cone coordinate. Its
restrictions are \(f\) and \(g\), proving compatibility.
The column first page uses the terms of \(K\) and its
functoriality statement is for complexes. These are the two computations in
\cite[Tags 015J and 015K]{stacks-derived}.
\end{proof}

\section{A composition of right derived functors}

\begin{theorem}[Grothendieck spectral sequence]
\label{res:G04}
Let \(F:\mathcal A\to\mathcal B\) and \(G:\mathcal B\to\mathcal C\)
be left exact between abelian categories. Suppose \(\mathcal A\) and
\(\mathcal B\) have enough injectives and \(F(I)\) is \(G\)-acyclic
for every injective \(I\) of \(\mathcal A\). For \(X\in D^+(\mathcal A)\)
there is a natural cohomological spectral sequence
\[
      E_2^{p,q}=R^pG\bigl(H^q(RF(X))\bigr)
      \quad\Longrightarrow\quad
      H^{p+q}\bigl(R(G\circ F)(X)\bigr)
\]
with differential of bidegree \((r,1-r)\) on page \(r\).
The target filtration is finite in each degree.
\end{theorem}

\begin{proof}
Choose a bounded-below injective model \(I^\bullet\) of \(X\).
Then \(F(I^\bullet)\) is a bounded-below representative of \(RF(X)\).
The acyclicity hypothesis means that every term of this representative
computes its own derived \(G\): its higher \(R^iG\) vanish, and
the degree-zero comparison is the isomorphism given by left exactness.
The acyclic-complex comparison of Theorem~\ref{res:R08} identifies
\(G(F(I^\bullet))\) with \(RG(F(I^\bullet))\) in the derived category.
This is the component of the canonical composition map of
Theorem~\ref{res:R12},
\[
           R(G\circ F)(X)\longrightarrow RG(RF(X)).
\]
It is an isomorphism natural in \(X\).

Apply the row sequence of Theorem~\ref{res:G03} to the functor \(G\)
and the bounded-below complex \(F(I^\bullet)\). Its second page is
the one displayed above, and it converges to \(H^*(RG(RF(X)))\).
Pull back its finite filtration along the canonical composition
isomorphism. Since the row construction from \(E_2\) onward is
functorial on the derived category, the choice of \(I^\bullet\)
or of a resolution for \(F(I^\bullet)\) gives the same comparisons
from the stated page onward. The filtration is finite because
the resolution degree is nonnegative and the horizontal degree
has a fixed lower bound. See \cite[Tags 015M and 015N]{stacks-derived}
and \cite[Section 5.8]{weibel}.
\end{proof}

\begin{corollary}[Object case, edge maps, and low degrees]
\label{res:G05}
Under the hypotheses of Theorem~\ref{res:G04}, an object \(A\) of
\(\mathcal A\), placed in degree zero, gives a first-quadrant sequence
\[
      E_2^{p,q}=R^pG(R^qF(A))
       \quad\Longrightarrow\quad R^{p+q}(G\circ F)(A).
\]
The natural edge maps in total degree \(n\ge0\) have the form
\[
 R^nG(F(A))\ \longrightarrow\ R^n(G\circ F)(A)\
                    \longrightarrow\ G(R^nF(A)).
\]
In degrees one and two they occur in the exact sequence
\[
\begin{split}
 0\longrightarrow R^1G(F(A))
 &\longrightarrow R^1(G\circ F)(A)
 \longrightarrow G(R^1F(A))\\
 &\xrightarrow{d_2^{0,1}} R^2G(F(A))
 \longrightarrow R^2(G\circ F)(A).
\end{split}
\]
\end{corollary}

\begin{proof}
Write \(Y=RF(A[0])\). It has cohomology in nonnegative degrees,
with \(H^0Y=F(A)\) and \(H^qY=R^qF(A)\) by
Theorem~\ref{res:R09}. This gives the quadrant assertion.
The first edge map is induced by the canonical map
\[
       F(A)[0]\simeq\tau_{\le0}Y\longrightarrow Y
\]
after applying \(RG\), taking \(H^n\), and identifying the
target via Theorem~\ref{res:G04}. The row sequence on \(F(A)[0]\)
has only its \(q=0\) row at \(E_2\). Its map to the sequence for
\(Y\) is the identity on the second-page term in bidegree
\((n,0)\). On cohomology it factors through the highest possibly
nonzero filtration step \(F^nH^n\), identifying the map as
\(E_2^{n,0}\twoheadrightarrow E_\infty^{n,0}
  \simeq F^nH^n\hookrightarrow H^n\).

For the second edge map use
\[
      Y\ \longleftarrow\ \tau_{\le n}Y\
               \longrightarrow H^n(Y)[-n].
\]
The left arrow induces an isomorphism on \(H^nRG\): its
truncation triangle has third term \(\tau_{\ge n+1}Y\), and
choosing a bounded-below injective model in degrees at least
\(n+1\) by Theorem~\ref{res:R01} shows its \(RG\) has zero cohomology
in degrees \(n-1,n\). Take the inverse isomorphism on
\(H^nRG\) and then the map to
\(H^nRG(H^n(Y)[-n])=G(H^n(Y))\).
By row functoriality the two arrows are the identity on the
\(E_2^{0,n}\) term, and the sequence for \(H^n(Y)[-n]\) has
only its \(q=n\) row at \(E_2\). In total degree \(n\) its
only term has filtration index zero. Thus the resulting map
is the edge map
\(H^n\twoheadrightarrow E_\infty^{0,n}\hookrightarrow E_2^{0,n}\).
This description also proves naturality of both edges.

In total degree one the graded pieces are
\(E_\infty^{1,0}=E_2^{1,0}\) and
\(E_\infty^{0,1}=\ker(d_2^{0,1})\). The only differential
that can enter \((2,0)\) from page two onward is \(d_2^{0,1}\);
none can leave \((2,0)\). Hence \(E_\infty^{2,0}\) is the
cokernel of this differential and injects into total degree two.
These identifications give exactness at all intervening terms
and no terminal surjectivity assertion. This is the filtration
argument of Theorem~\ref{res:S05} in the present sequence.
See \cite[Tag 015N]{stacks-derived} and
\cite[Theorem 5.8.3]{weibel}.
\end{proof}

\section{The homological dual}

\begin{theorem}[Grothendieck spectral sequence for left derivation]
\label{res:G06}
Let \(F:\mathcal A\to\mathcal B\) and \(G:\mathcal B\to\mathcal C\)
be right exact between abelian categories. Suppose \(\mathcal A\) and
\(\mathcal B\) have enough projectives and \(L_iG(F(P))=0\) for
\(i>0\) and projective \(P\) of \(\mathcal A\). For an object \(A\)
there is a natural first-quadrant homological spectral sequence
\[
      E^2_{p,q}=L_pG(L_qF(A))
       \quad\Longrightarrow\quad L_{p+q}(G\circ F)(A).
\]
More generally, for a derived object represented by a bounded-above
cochain complex \(X\), use homological grading \(H_j=H^{-j}\).
The second page is \(L_pG(H_q(LF(X)))\) and the target is
\(H_{p+q}(L(G\circ F)(X))\). The differential has bidegree
\((-r,r-1)\) and the target filtration is increasing and finite
in each degree.
\end{theorem}

\begin{proof}
Pass to the opposite abelian categories. The opposite functors
are left exact and the opposite source and intermediate
categories have enough injectives. Start with the chain indices
\(K_j=X^{-j}\), with differential \(\partial_j:K_j\to K_{j-1}\).
Viewing \(K_j\) in the opposite category gives the degree \(j\)
term of a cochain complex with differential \(\partial_{j+1}^{\mathrm{op}}\).
The chain lower bound becomes a cochain lower bound there.
Reversing arrows interchanges kernels with cokernels and
carries the degree \(j\) chain homology to degree \(j\)
cohomology in the opposite category. Projective resolutions
become injective resolutions. After reversing back, these
are the bounded-above projective models of
Theorem~\ref{res:R07}. The right acyclicity hypothesis of
Theorem~\ref{res:G04} is exactly the displayed vanishing,
with the degree-zero comparison supplied by right exactness.

Under arrow reversal the canonical comparison becomes
\[
              LG(LF(X))\longrightarrow L(G\circ F)(X).
\]
On a bounded-above projective model it is the canonical
derived comparison for applying \(G\) to a termwise
\(G\)-acyclic complex, hence an isomorphism also by
Theorem~\ref{res:R08}. Applying Theorem~\ref{res:G04}
in the opposite categories and reversing the arrows gives
the specified pages and target. A cohomological differential
from bidegree \((p,q)\) to \((p+r,q-r+1)\) becomes a homological
differential from the latter bidegree to the former. Equivalently,
the resulting page differential has bidegree \((-r,r-1)\).
The signed total in the opposite categories dualizes termwise
to \(d_h+(-1)^p d_v\) with both indices lowered. This agrees with
the chain convention of Chapter~\ref{ch:double-complexes}.

For the filtration, a decreasing subobject \(F^p\) of a target
in the opposite category becomes a quotient of the target in
the original category. Define the increasing step of index \(p-1\)
to be the kernel of that quotient. Reversing the short exact
sequence of successive decreasing steps identifies the
quotient of the increasing steps of indices \(p,p-1\) with the
reversed \(p\)-th graded piece. Finiteness follows from
finiteness in the opposite category. The object \(A[0]\) has
the nonnegative resolution degrees in both axes, giving the
first quadrant and the claimed object formula.

In particular, for \(A[0]\) the right-derived edge maps of
Corollary~\ref{res:G05} reverse to
\[
 G(L_nF(A))\longrightarrow L_n(G\circ F)(A)
                     \longrightarrow L_nG(F(A)).
\]
The five-term sequence reverses to
\[
\begin{split}
 L_2(G\circ F)(A)\longrightarrow L_2G(F(A))
 &\xrightarrow{d_2}G(L_1F(A))
 \longrightarrow L_1(G\circ F)(A)\\
 &\longrightarrow L_1G(F(A))\longrightarrow0 .
\end{split}
\]
Here \(d_2\) goes from \((2,0)\) to \((0,1)\). This is the
homological case of \cite[Corollary 5.8.4]{weibel}.
\end{proof}

\chapter[Group (Co)homology and Lie Cohomology]
  {Group (Co)homology and\\Lie Algebra Cohomology}
\label{ch:group-and-lie}

Throughout this chapter \(k\) is a commutative ring. Group modules are left
modules over the group algebra, and all groups under consideration are
discrete. For a Lie algebra \(\mathfrak g\) over \(k\), a coefficient
module is a left module over \(U(\mathfrak g)\), or equivalently a
\(k\)-module with a compatible \(\mathfrak g\)-action. The trivial module
\(k\) is defined by the augmentation of the respective algebra. We
apply the resolution arguments of Chapter~\ref{ch:derived-functors} and
the composition spectral sequence to these two augmented algebras.
The final section constructs coefficient group homology from a right
bar resolution and identifies it with derived coinvariants.

\section{Invariants and the bar resolution}

\begin{proposition}[Invariants and group cohomology]\label{res:A01}
For a group \(G\) and a left \(k[G]\)-module \(M\), evaluation at \(1\)
identifies \(\operatorname{Hom}_{k[G]}(k,M)\) with \(M^G\), naturally in
\(M\). Its right derived functors, written \(H^n(G,M)\), are
\(\operatorname{Ext}_{k[G]}^n(k,M)\), computed either from an injective
resolution of \(M\) or a projective resolution of \(k\).
\end{proposition}

\begin{proof}
The image of \(1\) under a module map is fixed by each \(g\), and a
fixed element \(m\) defines the map \(a\mapsto am\). This is a natural
identification of left exact functors to \(k\)-modules. The category
of modules over \(k[G]\) has enough injectives and projectives.
Theorems~\ref{res:R09} and~\ref{res:R11} give the derived
identifications and their independence of resolution.
See \cite[Section 6.1]{weibel}.
\end{proof}

\begin{theorem}[The bar resolution and nonhomogeneous cochains]
\label{res:A02}
Let \(B_n(G)=k[G^{n+1}]\), the free \(k\)-module on \((n+1)\)-tuples,
with the diagonal left \(G\)-action. With the alternating deletion
boundary and the augmentation \(B_0(G)\to k\) sending every vertex
to \(1\), this is a free \(k[G]\)-resolution of \(k\). Put
\(C^n(G,M)=\operatorname{Map}(G^n,M)\), with \(C^0(G,M)=M\).
The Hom complex from this resolution identifies with these cochains.
Its differential in degree \(n\) is
\begin{align*}
 (df)(g_1,\ldots,g_{n+1})
  &=g_1 f(g_2,\ldots,g_{n+1})\\
  &\quad+\sum_{i=1}^{n}(-1)^i
          f(g_1,\ldots,g_i g_{i+1},\ldots,g_{n+1})\\
  &\quad+(-1)^{n+1} f(g_1,\ldots,g_n).
\end{align*}
In particular \(H^n(G,M)\) is its \(n\)-th cohomology.
\end{theorem}

\begin{proof}
Two deleted positions in the boundary of a boundary occur in two
orders with opposite signs. The degree-one boundary has equal
augmentation on its two terms, so the augmented differential squares to
zero everywhere. Each tuple has a unique representative with its first
entry \(1\) after translation, which makes \(B_n(G)\) free as a left
\(k[G]\)-module. On the underlying augmented \(k\)-complex set
\[
 s_{-1}(1)=(1),\qquad
 s_n(g_0,\ldots,g_n)=(1,g_0,\ldots,g_n).
\]
Deleting the first entry of \(s_n\) gives the original tuple.
Deleting any other entry occurs with sign opposite to the corresponding
term of \(s_{n-1}\partial\). This proves \(\partial s+s\partial=1\),
including degrees zero and \(-1\). The homotopy need only be
\(k\)-linear: the resulting exactness of the underlying sequence is
exactness of the sequence of \(k[G]\)-modules.

The notation \([g_1|\cdots|g_n]\) denotes the tuple
\((1,g_1,g_1g_2,\ldots,g_1\cdots g_n)\) in a basis over \(k[G]\).
Deleting the first vertex in such a tuple gives
\(g_1[g_2|\cdots|g_n]\). Deleting the \(i\)-th internal vertex
multiplies its two adjacent increments, and deleting the last
vertex drops \(g_n\). Thus precomposition by the boundary gives
the displayed cochain formula. This uses the unsigned
precomposition convention for the Hom from a chain resolution,
as in Theorem~\ref{res:R11}. Equivalently, a homogeneous
cochain is a map \(\varphi:G^{n+1}\to M\) satisfying
\(\varphi(a\boldsymbol g)=a\varphi(\boldsymbol g)\).
Its values on the displayed basis determine it uniquely. The
identification commutes with all coefficient maps, and the free
resolution computes cohomology by Proposition~\ref{res:A01}.
See \cite[Sections 6.5.1--6.5.6]{weibel}.
\end{proof}

\begin{corollary}[Low degrees and the coefficient sequence]
\label{res:A03}
There are natural identifications
\[
 H^0(G,M)=M^G,\qquad
 H^1(G,M)=
 \frac{\{f:G\to M\mid f(gh)=f(g)+g f(h)\}}
      {\{g\mapsto g m-m\mid m\in M\}} .
\]
A short exact sequence of coefficients yields a natural long exact
sequence with connecting map given by the differential of a lifted cochain.
\end{corollary}

\begin{proof}
At degrees zero and one the differential in
Theorem~\ref{res:A02} gives precisely the indicated conditions.
The functors \(\operatorname{Hom}_{k[G]}(B_n(G),-)\) are exact because
the \(B_n(G)\) are free. The coefficient sequence thus gives a short
exact sequence of cochain complexes. If a cocycle in the quotient
complex is lifted to the middle complex, its differential belongs to
the kernel complex. A second lift changes this differential by a
kernel boundary, and a change of cocycle by a boundary permits a
change of lift by a differential. These are the long exact sequence's
connecting maps and they commute with maps of coefficient
sequences. When the coefficient module is injective the groups in
positive degree vanish, by Proposition~\ref{res:A01}. Thus the
cochain long exact sequences have the universality proved by
dimension shifting in Theorem~\ref{res:R09}. The identity on
invariants extends uniquely to their comparison with the
right-derived long exact sequences. See
\cite[Sections 6.1 and 6.4]{weibel}.
\end{proof}

\section{Restriction, Shapiro, and group extensions}

\begin{theorem}[Shapiro's lemma]\label{res:A04}
If \(H\leq G\) and \(L\) is a left \(k[H]\)-module, put
\[
 \operatorname{Coind}_H^G L=\operatorname{Hom}_{k[H]}(k[G],L),
 \qquad (g\varphi)(x)=\varphi(xg).
\]
For every \(n\) there is a natural isomorphism
\(H^n(G,\operatorname{Coind}_H^GL)\simeq H^n(H,L)\). On
homogeneous group cochains the comparison is restriction to tuples
in \(H\), followed by evaluation at \(1\) on the coefficient.
\end{theorem}

\begin{proof}
For any \(G\)-module \(X\), a \(G\)-map from \(X\) to the coinduced
module has evaluation at \(1\) an \(H\)-map \(X|_H\to L\). Conversely
an \(H\)-map \(u\) gives \(x\mapsto(a\mapsto u(ax))\), first for
\(a\in G\) and then by \(k\)-linearity. This map has the required
\(H\)-linearity in \(a\) and \(G\)-linearity in \(x\) by the given
action. The two formulas are inverse and commute with both module
variables. The restriction of \(B_\bullet(G)\) to \(H\) is still a
free resolution since \(k[G]\) is free as a left \(k[H]\)-module
on representatives of \(H\backslash G\). Hence the adjunction applied
in every degree gives the computation on that resolution.
The inclusion \(B_\bullet(H)\to B_\bullet(G)|_H\) lifts the
identity of \(k\), so comparison of resolutions
\cite[Chapter 10]{p103} identifies it with the canonical comparison
on cohomology. Following evaluation at \(1\) by this comparison
restricts homogeneous tuples as claimed. The same maps of complexes
for all \(L\) prove naturality. See \cite[Theorem 6.3.2]{weibel}.
\end{proof}

\begin{theorem}[The Lyndon--Hochschild--Serre spectral sequence]
\label{res:A05}
Let \(N\) be normal in \(G\), \(Q=G/N\), and let \(M\) be a left
\(k[G]\)-module. There is a natural first-quadrant spectral sequence
\[
 E_2^{p,q}=H^p(Q,H^q(N,M))\ \Longrightarrow\ H^{p+q}(G,M)
\]
with finite target filtration in each degree. The \(Q\)-action on
the coefficient is induced from conjugation on \(N\) and the action on
\(M\). Its two edges are inflation and restriction, the latter
taking values in the invariants. In particular it gives
\begin{align*}
 0\longrightarrow H^1(Q,M^N)
  &\xrightarrow{\mathrm{inf}}H^1(G,M)
  \xrightarrow{\mathrm{res}} H^1(N,M)^Q\\
  &\xrightarrow{d_2}H^2(Q,M^N)
  \xrightarrow{\mathrm{inf}}H^2(G,M).
\end{align*}
\end{theorem}

\begin{proof}
The functor \(M\mapsto M^N\) takes values in \(Q\)-modules. An
element \(g\) sends an \(N\)-invariant element to an \(N\)-invariant
element by normality, and \(N\) acts trivially there. A \(G\)-map
from an inflated \(Q\)-module factors uniquely through \(M^N\),
which proves that this functor is right adjoint to inflation.
Inflation is exact on underlying modules. Its right adjoint
sends injectives to injectives: lifting along a monomorphism
before applying the adjunction is exactly lifting along its
inflation. In addition, restriction from \(G\) to \(N\) preserves
injectives because the induction functor \(k[G]\otimes_{k[N]}-\)
is an exact left adjoint, by freeness as a right \(k[N]\)-module.
Thus if \(M\to I^\bullet\) is a \(G\)-injective resolution, the
cohomology of \(C^\bullet=(I^\bullet)^N\) is \(H^q(N,M)\) as a
\(k\)-module, with a canonical \(Q\)-action.

For the identification of the action, on homogeneous \(N\)-cochains
use
\[
 (g\varphi)(n_0,\ldots,n_j)
 =g\,\varphi(g^{-1}n_0g,\ldots,g^{-1}n_jg).
\]
This respects \(N\)-equivariance and commutes with the deletion
differential. On degree-zero invariants it acts through the action
on the coefficient. For \(a\in N\), equivariance identifies this
cochain operation with pullback by simultaneous right translation
by \(a\) on homogeneous tuples. A homotopy on the resolution from
the identity to this translation is
\[
 (n_0,\ldots,n_j)\longmapsto
 \sum_{i=0}^{j}(-1)^i
 (n_0,\ldots,n_i,n_i a,\ldots,n_j a).
\]
Deletions before and after the repeated vertex cancel against the
corresponding prisms of the boundary faces. The two deletions at
each transition cancel between consecutive terms, except for the
two end faces, giving right translation minus identity. This homotopy
is \(N\)-linear. The action on cohomology thus factors through \(Q\).
It agrees with the injective-resolution action: use the
first-quadrant array of homogeneous cochains \(C^j(N,I^r)\).
Taking its cohomology first in the resolution direction produces
\(C^*(N,M)\), since the bar terms are free. Taking it first in the
bar direction produces \((I^\bullet)^N\), since the restricted
terms are injective. Each of the two natural augmentations is a
quasi-isomorphism by Theorem~\ref{res:B04}, and both are \(G\)-linear.

Invariants satisfy \((M^N)^Q=M^G\). The hypotheses of
Theorem~\ref{res:G04} therefore hold for \(M\mapsto M^N\) followed
by \(Q\)-invariants, and its object case gives the asserted page
and the finite abutment filtration. To identify the row edge,
apply derived \(Q\)-invariants to the map
\((M^N)[0]\to (I^\bullet)^N\). On bar cochains it is compared
with the inclusion of coefficients and pullback of tuples along
\(G\to Q\). In detail, for each \(I^r\) pullback gives
\[
 C^j(Q,(I^r)^N)\longrightarrow C^j(G,I^r).
\]
It commutes with both differentials. Both horizontal computations
collapse onto degree zero by injectivity of \(I^r\) and \((I^r)^N\).
There the map is the identity of \((I^r)^G\). The augmentations
from the \(M\)-cochains are compatible with these maps, which
identifies the row edge with the usual inflation map. The column
edge on the models is induced by the inclusion
\((I^\bullet)^G\to (I^\bullet)^N\), with cohomology in the fixed
part. This is the usual restriction: restriction of homogeneous
tuples gives a map \(C^j(G,I^r)\to C^j(N,I^r)\), which restricts
to precisely that inclusion on degree-zero invariants; the two
augmentations of the arrays again commute. The five terms then
follow from Corollary~\ref{res:G05}. See \cite[Section 6.8]{weibel}.
\end{proof}

\section{The Chevalley--Eilenberg resolution}

Let \(\mathfrak g\) be a Lie algebra over \(k\); its bracket is
alternating and satisfies Jacobi. Write \(U=U(\mathfrak g)\).
It is the quotient of the tensor algebra by the relations
\(xy-yx=[x,y]\). The augmentation kills \(\mathfrak g\); a left
\(U\)-module determines its underlying Lie module by restriction.
In this section the underlying \(k\)-module of \(\mathfrak g\) is
free, with no restriction on the rank. Write
\(\widehat{\boldsymbol x}_i\) or \(\widehat{\boldsymbol x}_{ij}\)
for the wedge formed by omitting the indicated entries, preserving
the order of those that remain.

\begin{proposition}[The Chevalley--Eilenberg complex]\label{res:A06}
On \(V_n=U\otimes_k \bigwedge^n\mathfrak g\) put
\begin{align*}
 \partial(u\otimes x_1\wedge\cdots\wedge x_n)
 &=\sum_{i=1}^n (-1)^{i-1}u x_i\otimes
                       \widehat{\boldsymbol x}_i\\
 &\quad+\sum_{1\leq i<j\leq n}(-1)^{i+j}u\otimes
                   [x_i,x_j]\wedge \widehat{\boldsymbol x}_{ij}.
\end{align*}
These maps form an augmented \(U\)-linear chain complex \(V_\bullet\to k\).
Under evaluation on \(1\otimes-\), the complex
\(\operatorname{Hom}_U(V_\bullet,M)\) has degree \(n\) term
\(C^n(\mathfrak g,M)=
 \operatorname{Hom}_k(\bigwedge^n\mathfrak g,M)\).
The coboundary is
\begin{align*}
 (df)(x_1,\ldots,x_{n+1})
 &=\sum_{i=1}^{n+1}(-1)^{i-1}x_i f(\widehat{\boldsymbol x}_i)\\
 &\quad+\sum_{1\leq i<j\leq n+1}(-1)^{i+j}
                 f([x_i,x_j],\widehat{\boldsymbol x}_{ij}).
\end{align*}
\end{proposition}

\begin{proof}
The two formulas for the boundary are multilinear and alternating:
the sign of a removal is the exterior sign of moving the removed
factor to the front, with an additional minus sign for removal of
a pair. If two entries are equal, terms omitting neither retain
them and vanish, and terms omitting just one cancel in pairs.
The bracket of the two equal entries is zero. Thus the formula
descends to exterior powers without a restriction on the
characteristic. Right multiplication of the \(U\)-factor in a
summand commutes with its left \(U\)-action, which gives
\(U\)-linearity. The augmentation kills the image of \(\partial:V_1\to V_0\).

For the square of the differential, collect terms by the original
indices involved. Move them to the front preserving the order of
every other factor. This multiplies both contributions of a given
type by the same exterior sign. For two removed entries \(x,y\),
the two multiplication operations give \(u(xy-yx)\); multiplication
of the newly created bracket gives \(-u[x,y]\). These cancel by
the enveloping relation. For a removed \(x\) and a disjoint
bracket \([y,z]\), the products obtained in the two orders are
\(-ux\otimes[y,z]\) and \(+ux\otimes[y,z]\), with the remaining
factors understood. Two disjoint brackets yield opposite terms
because reversing their orders reverses the order of the two
bracket factors in the wedge. For a triple \(x,y,z\) used in a
nested bracket, the contribution is
\[
 u\otimes\bigl([[x,y],z]-[[x,z],y]+[[y,z],x]\bigr),
\]
wedged with the unused factors. This vanishes by Jacobi. These
cases exhaust the two applications of the differential, proving
\(\partial^2=0\). Finally a \(U\)-linear map \(V_n\to M\) is
determined by its value on \(1\otimes-\), with the inverse given
by \(u\otimes w\mapsto u f(w)\). Precomposition by \(\partial\)
now gives the displayed coboundary, naturally in \(M\).
See \cite[Section 7.7]{weibel}.
\end{proof}

\paragraph{The ordered monomials in the enveloping algebra.}
Choose a totally ordered \(k\)-basis of \(\mathfrak g\). In a word
in basis elements, replace a descending adjacent pair \(yx\) by
\(xy+[y,x]\), expressing the bracket in the basis. This procedure
terminates: in each resulting monomial either the word length
decreases or the length is unchanged and the inversion count
decreases. A word contains only finitely many letters and each
bracket is a finite linear combination, irrespective of the size
of the basis.

Disjoint replacements commute. The only overlap of two different
applicable replacements is a strictly descending triple \(zyx\).
Starting with either pair and commuting the remaining descending
pairs gives the same length-three term \(xyz\). The difference
of the two expressions of smaller length is
\[
 [y,x]\,z-z\,[y,x]+y\,[z,x]-[z,x]\,y+[z,y]\,x-x\,[z,y].
\]
Each difference of concatenations of two basis-linear factors
reduces to their Lie bracket. The resulting sum is
\([[y,x],z]+[y,[z,x]]+[[z,y],x]=0\) by Jacobi. These checks apply
with any unused prefix and suffix. Inducting on length and inversion
count now gives unique linear normal forms. Indeed, to compare the
first reductions, commute them in the disjoint case and use the
computed overlap difference in the overlapping case, and normalize
the resulting terms using the induction hypothesis. The normal form
map kills every instance of the defining relation multiplied by
a prefix and suffix, and thus kills the ideal they generate.
Conversely every replacement is the defining relation in the quotient.
Ordered words therefore form a \(k\)-basis of \(U\), and the
associated graded algebra for the word-length filtration is
\(\operatorname{Sym}_k(\mathfrak g)\): the generators commute in
the top degree, and the sorted monomials on each side are the
same basis. This is the free-module case of the
Poincar\'e--Birkhoff--Witt theorem of \cite[Theorem 7.3.7]{weibel}.

\begin{theorem}[Chevalley--Eilenberg computation of Ext]\label{res:A07}
If \(\mathfrak g\) is free as a \(k\)-module, then the augmentation
of Proposition~\ref{res:A06} is a free \(U(\mathfrak g)\)-resolution
of \(k\). In particular the \(n\)-th cohomology of its cochain
complex is naturally \(\operatorname{Ext}^n_{U(\mathfrak g)}(k,M)\).
\end{theorem}

\begin{proof}
Filter \(V_n\) increasingly by total weight: weight at most \(a\)
means that \(u\) belongs to the word-length step \(F_{a-n}U\),
interpreted as zero for \(a-n<0\). The multiplication part of the
boundary preserves weight and the bracket part lowers it by one.
The sorted-word description above identifies the associated graded
complex with the Koszul chain complex on the variables of
\(\operatorname{Sym}_k(\mathfrak g)\); the differential is the sum
of signed deletions with multiplication by the deleted variable.
The augmentation on its degree-zero terms is the quotient map
by the variables.

For a finite subset of the basis, take the Koszul complex on that
subset over the whole polynomial ring. Starting with no variables,
adjoining one variable adds the two-term factor whose boundary
is multiplication by that variable. Equivalently the larger complex
is the chain mapping cone of multiplication by the new variable
on the smaller complex, with the tensor differential sign of
Chapter~\ref{ch:double-complexes}. The long exact sequence of that
cone computes its homology: by induction the smaller complex has
only degree-zero homology, the quotient by its variables, and
the new variable acts injectively on this polynomial quotient.
The cone thus has degree-zero homology the quotient by all variables
in the enlarged subset and zero homology in positive degrees.
For an infinite basis every wedge chain uses only a finite subset.
A positive-degree cycle is therefore a boundary in some such
finite-subset complex. The degree-zero homology of the union is
the quotient by all variables, namely \(k\). This argument proves
exactness of the augmented associated graded complex, in every
weight since the Koszul boundary preserves that weight.

Now let \(z\in V_n\) be a cycle for \(n>0\), or let \(z\in V_0\)
lie in the kernel of the augmentation. It has finite filtration
weight. Its leading class is a cycle of the augmented graded
complex. By graded exactness lift a graded antecedent to a
chain of one higher degree in the same filtration step. Subtracting
its boundary from \(z\) lowers its filtration weight. Repeating
terminates because the filtration in degree \(n\) is zero below
weight \(n\). The sum of the chosen antecedents shows that \(z\)
is a boundary. The augmentation is onto, and the modules \(V_n\)
are free over \(U\), since ordered wedges of basis elements form
a \(k\)-basis of \(\bigwedge^n\mathfrak g\).
The resolution computes Ext by Theorem~\ref{res:R11}, with the
identification of cochains in Proposition~\ref{res:A06}.
See \cite[Theorem 7.7.2 and Corollary 7.7.3]{weibel}.
\end{proof}

\begin{corollary}[Low-degree Lie algebra cohomology]\label{res:A08}
For the coefficients of Proposition~\ref{res:A06},
\[
 H^0(\mathfrak g,M)=M^{\mathfrak g},\qquad
 H^1(\mathfrak g,M)=
 \frac{\{f:\mathfrak g\to M\text{ \(k\)-linear}\mid
 f([x,y])=x f(y)-y f(x)\}}
 {\{x\mapsto x m\mid m\in M\}} .
\]
The corresponding differentials are \(dm(x)=xm\) and
\(df(x,y)=xf(y)-yf(x)-f([x,y])\).
\end{corollary}

\begin{proof}
These are the formulas of Proposition~\ref{res:A06} with zero
and one input in the cochain, respectively. For a degree-zero
vector \(m\), applying the second formula to \(x\mapsto xm\)
gives \(x(ym)-y(xm)-[x,y]m=0\), using the defining module
relation. Thus quotienting cocycles by these boundaries gives
exactly derivations modulo inner derivations. All identifications
are natural in \(M\), and Theorem~\ref{res:A07} identifies them
with the corresponding Ext groups. See
\cite[Definition 7.4.3 and Theorem 7.4.7]{weibel}.
\end{proof}

\section{The Lie algebra Hochschild--Serre comparison}

\begin{theorem}[An ideal and its quotient]\label{res:A09}
Let \(k\) be a field, \(\mathfrak g\) a Lie algebra over \(k\),
\(\mathfrak h\) an ideal, and \(\mathfrak q=\mathfrak g/\mathfrak h\).
For any \(\mathfrak g\)-module \(M\) there is a natural first-quadrant
spectral sequence
\[
 E_2^{p,q}=H^p(\mathfrak q,H^q(\mathfrak h,M))
           \ \Longrightarrow\ H^{p+q}(\mathfrak g,M)
\]
with finite target filtration in each degree. Its row and column
edges are induced by inflation and restriction of Chevalley--Eilenberg
cochains, and the restriction edge lands in the \(\mathfrak q\)-invariants.
The five terms in degree one and two are
\begin{align*}
 0\longrightarrow H^1(\mathfrak q,M^{\mathfrak h})
 &\xrightarrow{\mathrm{inf}}H^1(\mathfrak g,M)
 \xrightarrow{\mathrm{res}}H^1(\mathfrak h,M)^{\mathfrak q}\\
 &\xrightarrow{d_2}H^2(\mathfrak q,M^{\mathfrak h})
 \xrightarrow{\mathrm{inf}}H^2(\mathfrak g,M).
\end{align*}
\end{theorem}

\begin{proof}
Choose a basis of \(\mathfrak h\) and extend it to a basis of
\(\mathfrak g\). Order the complementary letters before the
\(\mathfrak h\)-letters. The ordered-word computation preceding
Theorem~\ref{res:A07} shows that the enveloping algebra of
\(\mathfrak h\) embeds in \(U(\mathfrak g)\), and that
\(U(\mathfrak g)\) is free as a right \(U(\mathfrak h)\)-module
on the ordered complementary words. Hence restriction preserves
injectives, being right adjoint to exact induction by tensoring
with that free module. The \(\mathfrak h\)-invariant part of a
\(\mathfrak g\)-module is a \(\mathfrak q\)-module:
if \(x\in\mathfrak h\), \(y\in\mathfrak g\), and \(xm=0\) for all
such \(x\), then \(x(ym)=y(xm)+[x,y]m=0\).
The invariant functor is right adjoint to exact inflation
because a map from an inflated module has image annihilated
by \(\mathfrak h\). It thus sends injectives to injectives.
The composite with \(\mathfrak q\)-invariants is precisely
\(\mathfrak g\)-invariants. All three categories of modules
have enough injectives. Theorems~\ref{res:A07} and~\ref{res:R09}
identify the derived invariants with the cohomology groups in
the statement, and Theorem~\ref{res:G04} applies.

For clarity, the quotient action may also be computed on the
cochains of Proposition~\ref{res:A06}. For \(y\in\mathfrak g\)
define
\[
 (L_y f)(x_1,\ldots,x_j)
 =y f(x_1,\ldots,x_j)-
   \sum_{i=1}^{j} f(x_1,\ldots,[y,x_i],\ldots,x_j).
\]
The bracket terms are in \(\mathfrak h\). The action property
follows from the module relation and Jacobi, which identifies the
commutator of the adjoint operators with the adjoint of the bracket.
It commutes with the cochain differential: on action terms this
uses \(y(xm)-x(ym)=[y,x]m\), and on bracket terms it uses that
\([y,-]\) is a derivation of the bracket. For \(a\in\mathfrak h\)
put \((\iota_a f)(x_1,\ldots,x_{j-1})=f(a,x_1,\ldots,x_{j-1})\),
and set \(\iota_a=0\) in degree zero. In
\(d\iota_a f+\iota_a df\) the action and bracket terms
on the \(x_i\) alone cancel in pairs after moving \(a\) to the
first argument. The terms left are the action by \(a\) and
one term \(-f(x_1,\ldots,[a,x_i],\ldots,x_j)\) per index.
Thus \(d\iota_a+\iota_a d=L_a\), so the action descends to
\(\mathfrak q\) on cohomology.

To compare it with the action used by derived invariants, choose
a \(\mathfrak g\)-injective resolution \(M\to I^\bullet\).
The double complex \(C^j(\mathfrak h,I^r)\) has the displayed
\(\mathfrak g\)-action by cochain maps. The same two augmentations
as in the group proof apply: exactness of
\(\operatorname{Hom}_k(\bigwedge^j\mathfrak h,-)\) computes the
resolution-direction cohomology, while restriction preserving
injectives makes the Chevalley--Eilenberg direction vanish in
positive degrees for each \(I^r\). By Theorem~\ref{res:B04},
the total cohomology comparisons are quasi-isomorphisms and
identify the two actions. This proves the needed identification on
the second page.

Finally, the quotient map \(\mathfrak g\to\mathfrak q\) pulls
back alternating cochains with values in \(M^{\mathfrak h}\),
and its pullback commutes with the differential because it is
a Lie map compatible with the action on these coefficients.
Restriction is likewise a cochain map using the ideal inclusion.
For each \(I^r\), pullback compares
\(C^*(\mathfrak q,(I^r)^{\mathfrak h})\) with
\(C^*(\mathfrak g,I^r)\). In degree-zero invariants it is the
identity on \((I^r)^{\mathfrak g}\); the higher cohomology on
both sides vanishes by injectivity. Thus it identifies the row
edge induced by \((M^{\mathfrak h})[0]\to (I^\bullet)^{\mathfrak h}\)
with inflation, because the comparison commutes with the coefficient
augmentations induced by \(M\to I^\bullet\). In the same way,
restriction of cochains
computes the inclusion of invariant complexes
\((I^\bullet)^{\mathfrak g}\to(I^\bullet)^{\mathfrak h}\)
giving the column edge. Apply Corollary~\ref{res:G05} to get
the five terms. See \cite[Section 7.5]{weibel}.
\end{proof}

\Needspace{24\baselineskip}
\section{Coinvariants and coefficient group homology}

\begin{theorem}[Group homology and its coefficient sequence]
\label{res:A10}
Let \(G\) be a discrete group, put \(R=k[G]\), and let \(M\) be
a left \(R\)-module. Its coinvariants are
\[
 M_G=M/\langle gm-m:g\in G,\ m\in M\rangle_k
       \simeq k\otimes_R M,
\]
where \(k\) is the trivial right \(R\)-module. For \(n\ge0\) define
\[
 H_n(G,M)=\operatorname{Tor}^{R}_n(k,M).
\]
These are naturally the left derived functors of coinvariants,
with \(H_0(G,M)=M_G\). They can be computed from the complex
\[
 C_n(G,M)=\bigoplus_{(g_1,\ldots,g_n)\in G^n}M .
\]
Write \(m[g_1|\cdots|g_n]\) for a vector in the displayed summand
and \(C_0(G,M)=M\). Its differential for \(n\ge1\) is
\begin{align*}
 \partial\bigl(m[g_1|\cdots|g_n]\bigr)
 &=g_1^{-1}m[g_2|\cdots|g_n]\\
 &\quad+\sum_{i=1}^{n-1}(-1)^i
       m[g_1|\cdots|g_i g_{i+1}|\cdots|g_n]\\
 &\quad+(-1)^n m[g_1|\cdots|g_{n-1}].
\end{align*}
Every short exact sequence \(0\to M'\to M\to M''\to0\) of
left \(R\)-modules gives a natural long exact sequence
\begin{align*}
 \cdots &\longrightarrow H_n(G,M')\longrightarrow H_n(G,M)
       \longrightarrow H_n(G,M'')\\
 &\xrightarrow{\delta}H_{n-1}(G,M')\longrightarrow\cdots
       \longrightarrow M_G\longrightarrow M''_G\longrightarrow0 .
\end{align*}
\end{theorem}

\begin{proof}
In \(k\otimes_R M\), the tensor relation identifies
\(1\otimes gm\) with \(1g\otimes m=1\otimes m\). Thus
\(m\mapsto1\otimes m\) factors through \(M_G\). Its inverse sends
\(a\otimes m\) to \(a[m]\) for \(a\in k\), and the augmentation
formula for the trivial right action makes this inverse balanced.
This proves the zeroth identification and shows that coinvariants
are right exact.

Take the free left bar resolution \(B_\bullet(G)\to k\) of
\ref{res:A02}. On its underlying \(k\)-modules put
\(b\cdot g=g^{-1}b\), extending \(k\)-linearly to a right
\(R\)-action. Inversion reverses products, so
\((b\cdot g)\cdot h=b\cdot(gh)\). The left-equivariant
differential is also right-equivariant for this action.
Call this augmented complex \(B_\bullet^{\mathrm r}(G)\to k\).
The inverse twist of a free left \(R\)-module is a free right
\(R\)-module: if its left basis is \((e_\alpha)\), then
the elements \(e_\alpha\cdot g=g^{-1}e_\alpha\) are the same
underlying \(k\)-basis as \(g e_\alpha\).
Its exactness is the exactness already proved for the underlying
bar complex, so this is a free right resolution of the trivial
right module.

Use the same free basis
\([g_1|\cdots|g_n]=(1,g_1,g_1g_2,\ldots,g_1\cdots g_n)\)
as in \ref{res:A02}, now as a right-module basis.
Then \(B_n^{\mathrm r}(G)\otimes_R M\) identifies with the
displayed direct sum by
\([g_1|\cdots|g_n]\otimes m\mapsto m[g_1|\cdots|g_n]\).
In the deletion boundary, the first face of this tuple is
the left translate \(g_1[g_2|\cdots|g_n]\).
In the twisted module it is
\([g_2|\cdots|g_n]\cdot g_1^{-1}\); balancing therefore puts
\(g_1^{-1}\) on the coefficient \(m\). The interior and final
faces are unchanged. This proves the differential formula,
and \(\partial^2=0\) also follows from the deletion calculation.
The exponent \(-1\) in its first term records precisely this
choice of right-module basis. At degree one the image is
spanned by \(g^{-1}m-m\), the same submodule as that spanned by
\(gm-m\), in agreement with the formula for \(H_0\).

To compare with the derived functor in the coefficient variable,
choose a projective left resolution \(P_\bullet\to M\).
The first-quadrant double complex
\(B_p^{\mathrm r}(G)\otimes_R P_q\) has finite diagonals.
Each \(B_p^{\mathrm r}(G)\) is free, so resolving \(M\) in the
\(q\)-direction identifies the total homology with the homology
of \(B_\bullet^{\mathrm r}(G)\otimes_R M\).
Each \(P_q\) is flat, so resolving \(k\) in the \(p\)-direction
identifies that same homology with the homology of
\(k\otimes_R P_\bullet\). These are the two resolution
comparisons of \ref{res:B05}. The latter complex computes the
left derived coinvariant functor by \ref{res:R07}; equivalently
place \(P_q\) in cochain degree \(-q\) as in \ref{res:R14}.
Thus the bar model, the two computations of Tor, and the derived
coinvariants agree naturally. Comparison maps between projective
resolutions are unique up to homotopy, so the comparisons are
independent of the choices.

Finally, the complex \(C_n(G,-)\) is a direct sum of copies of
the coefficient module, hence is exact in every degree.
A coefficient short exact sequence therefore gives a short
exact sequence of chain complexes. For a cycle \(z''\) in the
quotient complex, lift it to \(z\) in the middle complex.
The vector \(\partial z\) belongs to the image of the first
complex. Its unique preimage is a cycle by injectivity and
\(\partial^2=0\), and its homology class defines
\(\delta[z'']\). Changing the lift changes that preimage by a
boundary; changing the cycle representative by a boundary has
the same effect after choosing a lifted boundary. These maps
commute with maps of coefficient sequences and give the
displayed long exact sequence. To compare the connecting maps
with the coefficient-derived construction, take a horseshoe of
projective coefficient resolutions for the three modules.
The two augmentations of each double complex used above give
maps of short exact sequences of complexes: the horseshoe is
degreewise split, and the right bar terms are free. Both maps
induce the displayed homology isomorphisms. Naturality of the
lift-and-boundary construction transports the same connecting
map across both comparisons.
See \cite[Section 6.1 and Application 6.5.4]{weibel}.
\end{proof}

\chapter{Sheaf Cohomology and Leray}
\label{ch:sheaves}

We work with sheaves of modules on a ringed space \((X,\mathcal O_X)\).
An open subspace carries the restricted structure sheaf. For abelian
sheaves one can use the constant sheaf of rings \(\underline{\mathbb Z}\).
The calculations retain the module structures over the section rings.
In particular a map of ringed spaces induces restriction of scalars
on global sections even though its pushforward is already a functor of
categories of sheaves of modules.

\section{Cohomology and acyclicity on open sets}

The stalk functor at \(x\) is exact and has a right adjoint sending an
\(\mathcal O_{X,x}\)-module to its skyscraper sheaf at \(x\), with
sections equal to that module on opens containing \(x\) and zero
otherwise. The adjunction follows by sending a morphism into the
skyscraper to its value on germs at \(x\). Choose an embedding of each
stalk of \(\mathcal F\) into an injective \(\mathcal O_{X,x}\)-module.
The induced map into the product of the skyscrapers is monic, since
its projection at each \(x\) induces the chosen injection on the
\(x\)-stalk. Every factor is injective by the adjunction with an exact
functor, and the product is injective by extending one component at
a time. Thus the category of sheaves of modules has enough injectives.

For an open inclusion \(j:V\hookrightarrow X\), extension by zero
\(j_!\) is left adjoint to restriction and is exact. Its stalks are the
original stalks on \(V\) and zero outside \(V\), which proves exactness.
In particular the restriction of an injective sheaf to \(V\) is
injective. The adjunction identifies
\(\operatorname{Hom}(j_!\mathcal O_V,\mathcal F)\) with \(\mathcal F(V)\).
If \(V\subseteq W\) are open, the evident map of extensions by zero
of the unit modules is monic, as is seen on stalks. An injective
sheaf therefore has surjective restrictions between arbitrary opens.
A sheaf with this property is called flasque.

\begin{theorem}[Sections, flasque models, and change to abelian sheaves]
\label{res:H01}
For an open \(U\subseteq X\) define
\[
  H^i(U,\mathcal F)=H^i\bigl(\Gamma(U,I^\bullet)\bigr)
\]
from an injective resolution of \(\mathcal F\) on \(X\). This is also
the \(i\)-th right derived functor of sections of \(\mathcal F|_U\)
on \(U\). Its degree zero is \(\mathcal F(U)\), and a short exact
sequence of sheaves gives the natural cohomology sequence in
\(\mathcal O_X(U)\)-modules. Flasque sheaves are acyclic for sections
over every open. Taking underlying abelian sheaves identifies the
underlying groups of these cohomology objects with ordinary
cohomology of abelian sheaves.
\end{theorem}
\begin{proof}
The adjunction for \(j:U\hookrightarrow X\) shows that the restricted
injective resolution is an injective resolution on \(U\).
The degree-zero and long exact sequence claims follow from left
exactness of sections and Theorem~\ref{res:R09}.

To prove acyclicity, in \(0\to\mathcal F\to\mathcal G\to\mathcal H\to0\)
with \(\mathcal F\) flasque, every section of \(\mathcal H\) over an open
\(U\) lifts over \(U\). Partially lift it over open subsets of \(U\).
A chain of partial lifts glues over the union. If a maximal partial
lift omits a point of \(U\), the epimorphism of sheaves supplies a lift
near that point. On the overlap, the difference lies in \(\mathcal F\);
extend this difference to the new open and subtract it from the new lift.
The two lifts then glue, contradicting maximality. This proves
surjectivity on sections over every open. If \(\mathcal G\) is also
flasque, a section of \(\mathcal H\) lifts over its entire open
and its lift extends to any larger open. Hence \(\mathcal H\) is flasque.
Applied to an injective resolution of a flasque \(\mathcal F\), this
proves by induction that all successive cokernels are flasque and that
the resolution is exact after taking sections in positive degrees.

Forgetting scalars is exact on stalks and does not change restrictions.
An injective resolution of \(\mathcal F\) as an \(\mathcal O_X\)-module
becomes a flasque resolution of the underlying abelian sheaf, so it
computes the cohomology there by Theorem~\ref{res:R08}. The comparison
is natural because it is the canonical comparison through acyclic
models. On abelian sheaves one can also express sections on \(U\) as
\(\operatorname{Hom}(j_!\underline{\mathbb Z},-)\), so
Theorem~\ref{res:R11} gives the comparison with Ext from this free
abelian sheaf. See \cite[Tags 0712, 01E1, 01E3, 09SY]{stacks-sheaf}.
\end{proof}

\section{\v{C}ech cochains and the comparison spectral sequence}

\begin{definition}[\v{C}ech cochains for an indexed cover]\label{res:H02}
Let \(\mathcal U=(U_i)_{i\in I}\) be an open cover of \(U\).
Use all ordered tuples, including repetitions, and put
\[
  \check C^p(\mathcal U,\mathcal F)=
    \prod_{(i_0,\ldots,i_p)\in I^{p+1}}
       \mathcal F(U_{i_0}\cap\cdots\cap U_{i_p}).
\]
Its differential on the component \((i_0,\ldots,i_{p+1})\) is
the sum over \(a\) of \((-1)^a\) times the restriction of the
component with \(i_a\) omitted. The same formula defines the
complex for a presheaf of modules.
\end{definition}

Deleting positions \(a<b\) in opposite orders gives the same restriction
with signs \((-1)^{a+b-1}\) and \((-1)^{a+b}\), so \(d^2=0\).
The degree-zero augmentation from sections on \(U\) lands in the
kernel of the first differential. For a sheaf it is an isomorphism
onto that kernel, by the gluing axiom. All products and restrictions
are in \(\mathcal O_X(U)\)-modules via restriction of scalars.
See \cite[Sections 9--11]{stacks-sheaf}.

We need the following acyclicity calculation for the full, unnormalized
cochains. Write \(j_J:U_{i_0}\cap\cdots\cap U_{i_p}\to U\) for the open
inclusion indexed by a tuple \(J\). On \(U\) consider the augmented
chain complex with \(p\)-th term
\[
     \bigoplus_{J\in I^{p+1}}(j_J)_!\mathcal O_{U_J}
\]
and augmentation to \(\mathcal O_U\). The face maps on each summand
are the extension-by-zero inclusions, with the alternating signs.
At a point \(x\) the complex is the free chain module on tuples
of indices containing \(x\), over \(\mathcal O_{U,x}\). Choose one
index containing \(x\). Prepending it (also in degree \(-1\))
is a contraction of the augmented complex: the first deletion is
the identity and the other deletions cancel the prepending of the
old boundary. Stalks commute with these direct sums, so this proves
exactness of the augmented chain of sheaves. Apply
\(\operatorname{Hom}_{\mathcal O_U}(-,I')\) for an injective \(I'\).
Exactness is preserved contravariantly, and the adjunction for each
\(j_J\) gives exactly the product cochains above. Thus the
augmented \v{C}ech cochain complex is exact for each injective \(I'\).

\begin{theorem}[\v{C}ech-to-cohomology spectral sequence]\label{res:H03}
For a cover \(\mathcal U\) of \(U\) and a sheaf of modules \(\mathcal F\),
there is a first-quadrant spectral sequence
\[
 E_2^{p,q}=
 \check H^p\bigl(\mathcal U,\;V\longmapsto H^q(V,\mathcal F)\bigr)
 \quad\Longrightarrow\quad H^{p+q}(U,\mathcal F).
\]
It is natural in \(\mathcal F\) and induces a finite filtration of each
abutment group. The coefficient in the \v{C}ech expression is the
presheaf of local cohomology groups.
\end{theorem}
\begin{proof}
Let \(\mathcal F\to I^\bullet\) be an injective resolution. Form
\(D^{p,q}=\check C^p(\mathcal U,I^q)\), with commuting differentials
and the totalization of Proposition~\ref{res:B01}. The horizontal
augmentation by \(\Gamma(U,I^q)\) is exact for every \(q\), by the
preceding computation and exact open restriction. Therefore
\(\Gamma(U,I^\bullet)\to\operatorname{Tot}D\) is a quasi-isomorphism
by the augmented comparison in Theorem~\ref{res:B04}.

In the column filtration, vertical cohomology is the product of the
groups \(H^q(U_J,\mathcal F)\). Here products of modules are exact,
and \(I^\bullet|_{U_J}\) is an injective resolution. The horizontal
induced differential is the alternating restriction of these groups.
Consequently the second page is the displayed \v{C}ech term by
Proposition~\ref{res:B03}. First-quadrant support gives convergence and a
finite filtration by Theorems~\ref{res:B04} and~\ref{res:S01}.
Maps of coefficients lift to maps of injective resolutions and any
two lifts are homotopic by Theorem~\ref{res:R04}. On each column
the induced homotopy gives equality on vertical cohomology. The
total homotopy, multiplied by \((-1)^p\) in column \(p\), preserves
the column filtration and gives equality of the abutment maps.
Page functoriality now proves the stated naturality. See
\cite[Tags 01EP, 01ES]{stacks-sheaf}.
\end{proof}

\begin{corollary}[\v{C}ech comparison on acyclic intersections]
\label{res:H04}
If \(H^q(U_{i_0}\cap\cdots\cap U_{i_p},\mathcal F)=0\) for every
finite tuple and every \(q>0\), the natural comparison
\[
      \check H^p(\mathcal U,\mathcal F)\longrightarrow H^p(U,\mathcal F)
\]
is an isomorphism for each \(p\).
\end{corollary}
\begin{proof}
In Theorem~\ref{res:H03} only the \(q=0\) row survives. Its
identification with the target comes from the edge map, since the
target filtration has one nonzero graded layer in each total degree.
Concretely the comparison is induced by
\(\check C^\bullet(\mathcal U,\mathcal F)\to\operatorname{Tot}D\),
using the augmentation in vertical degree zero, followed on cohomology
by the inverse of the horizontal augmentation isomorphism.
It is therefore the usual comparison and natural in coefficients.
See \cite[Tag 01ET]{stacks-sheaf}.
\end{proof}

\begin{theorem}[Mayer--Vietoris for two opens]\label{res:H05}
For \(U=V\cup W\) there is a natural long exact sequence
\[
 \cdots\longrightarrow H^n(U,\mathcal F)
 \longrightarrow H^n(V,\mathcal F)\oplus H^n(W,\mathcal F)
 \xrightarrow{\operatorname{res}_V-\operatorname{res}_W}
 H^n(V\cap W,\mathcal F)
 \longrightarrow H^{n+1}(U,\mathcal F)\longrightarrow\cdots .
\]
\end{theorem}
\begin{proof}
For each injective \(I^q\) the sequence of sections is short exact.
Gluing identifies the kernel of the difference map. Surjectivity
follows by extending a section on \(V\cap W\) to \(V\), since \(I^q\)
is flasque. These short exact sequences form a short exact sequence
of complexes. Its ordinary cohomology sequence, with connecting map
given by lifting and taking the differential, is the displayed
sequence by Theorem~\ref{res:H01}. This convention agrees with the
triangle of Theorem~\ref{res:D04}; using the opposite order of \(V,W\)
reverses the sign of the difference map together with its boundary
identification. See \cite[Tag 01EB]{stacks-sheaf}.
\end{proof}

\section{Higher direct images and composition}

\begin{theorem}[Pushforward, local cohomology, and global sections]
\label{res:H06}
Let \(f:X\to Y\) be a map of ringed spaces. The sheaf \(R^q f_*\mathcal F\)
is the sheafification of the presheaf on \(Y\) sending \(V\) to
\(H^q(f^{-1}V,\mathcal F)\), with its natural \(\mathcal O_Y(V)\)-action.
If \(I\) is injective on \(X\), \(f_*I\) is acyclic for sections over
every open of \(Y\). There is a natural isomorphism
\[
  \operatorname{Res}\circ R\Gamma(X,-)
     \ \simeq\ R\Gamma(Y,-)\circ Rf_*
   \quad\text{on }D^+(\mathcal O_X\text{-}\mathrm{Mod}).
\]
Here \(\operatorname{Res}\) restricts scalars along
\(\mathcal O_Y(Y)\to\mathcal O_X(X)\).
\end{theorem}
\begin{proof}
The pushforward is additive and left exact. Its bounded right derived
functor is constructed in Theorem~\ref{res:R06}. For an injective
resolution \(I^\bullet\), cohomology sheaves of \(f_*I^\bullet\)
are obtained by sheafifying the presheaf cohomology, since the
comparison is an isomorphism on stalks by exactness of filtered
colimits of modules. Sections of \(f_*I^\bullet\) on \(V\) are sections
of \(I^\bullet\) on \(f^{-1}V\), giving the asserted local description
by Theorem~\ref{res:H01}. Pushforward preserves flasqueness since
the relevant restriction maps are those on inverse-image opens.
Thus \(f_*I\) is acyclic on every open by Theorem~\ref{res:H01}.
Apply the right composition comparison of Theorem~\ref{res:R12}
to \(f_*\) and \(\Gamma(Y,-)\). Their underived composite is
\(\operatorname{Res}\circ\Gamma(X,-)\). Restriction of scalars is
exact, and termwise restriction of the complex of sections on
an injective model computes the right derived composite.
This gives the displayed isomorphism, with the indicated coefficient
ring. See \cite[Tags 0713, 01E4, 01EX, 01EZ]{stacks-sheaf}.
\end{proof}

\begin{theorem}[Leray spectral sequence and two degenerations]
\label{res:H07}
For \(\mathcal F\) an \(\mathcal O_X\)-module, there is a first-quadrant
spectral sequence of \(\mathcal O_Y(Y)\)-modules
\[
 H^p(Y,R^q f_*\mathcal F)\ \Longrightarrow\
     \operatorname{Res} H^{p+q}(X,\mathcal F).
\]
For \(K\in D^+(\mathcal O_X\text{-}\mathrm{Mod})\), use \(H^q(Rf_*K)\)
on the second page and hypercohomology at the target. Here \(p\ge0\),
and \(q\in\mathbb Z\) is bounded below for each \(K\). Each abutment has a
finite filtration. For a sheaf \(\mathcal F\) the edge comparison gives
\[
\begin{array}{ll}
 H^n(Y,f_*\mathcal F)\ \simeq\ \operatorname{Res}H^n(X,\mathcal F)
    &\text{if all }R^q f_*\mathcal F=0\text{ for }q>0,\\[2pt]
 \operatorname{Res}H^n(X,\mathcal F)\ \simeq\
    \Gamma(Y,R^n f_*\mathcal F)
    &\text{if all }H^p(Y,R^q f_*\mathcal F)=0\text{ for }p>0 .
\end{array}
\]
For another map of ringed spaces \(g:Y\to Z\) the relative version has
second page \(R^p g_*(H^q(Rf_*K))\) and target
\(H^{p+q}(R(gf)_*K)\), with the same indexing convention.
\end{theorem}
\begin{proof}
Take \(F=f_*\) and \(G=\Gamma(Y,-)\) in
Theorem~\ref{res:G04}. The required acyclicity and the identification
of the composite are those of Theorem~\ref{res:H06}.
Corollary~\ref{res:G05} gives the first-quadrant case for a sheaf.
The two vanishing hypotheses leave only the row \(q=0\), or only the
column \(p=0\), respectively. The single surviving layer in each
total degree identifies the corresponding edge map with an isomorphism.

For the relative version, \(f_*I\) is flasque and is acyclic for
sections over every open of \(Y\). By the local description in
Theorem~\ref{res:H06} applied to \(g\), this implies
\(R^p g_*(f_*I)=0\) for \(p>0\). Thus the same hypothesis of the
derived composition theorem holds for \(g_*\).
Use \(g_*f_*=(gf)_*\) and Theorem~\ref{res:G04} again. The spectral
sequences and their edge comparisons are natural by the corresponding
constructions of Chapter~\ref{ch:grothendieck}.
See \cite[Tags 01F2, 01F4, 01F6]{stacks-sheaf}.
\end{proof}


\backmatter
\begin{thebibliography}{10}
\addcontentsline{toc}{chapter}{Bibliography}

\bibitem{weibel}
Charles A.~Weibel.
\emph{An Introduction to Homological Algebra}.
Cambridge Studies in Advanced Mathematics 38.
Cambridge University Press, 1994, first paperback printing 1995.

\bibitem{mccleary}
John McCleary.
\emph{A User's Guide to Spectral Sequences}, second edition.
Cambridge Studies in Advanced Mathematics 58.
Cambridge University Press, 2001.

\bibitem{yekutieli}
Amnon Yekutieli.
\emph{Derived Categories}.
Cambridge Studies in Advanced Mathematics 183.
Cambridge University Press, 2020.

\bibitem{rotman}
Joseph J.~Rotman.
\emph{An Introduction to Homological Algebra}, second edition.
Universitext. Springer, 2009.

\bibitem{stacks-homology}
The Stacks Project Authors.
\emph{The Stacks Project}, Homology, especially
\href{https://stacks.math.columbia.edu/tag/012K}{Filtered complexes}
and \href{https://stacks.math.columbia.edu/tag/012X}{Double complexes}.
Individual results are identified by section or tag.

\bibitem{stacks-derived}
The Stacks Project Authors.
\emph{The Stacks Project}, \href{https://stacks.math.columbia.edu/tag/05QI}{Derived Categories}.
Individual results are identified by section or tag.
Archived reference snapshot \texttt{ed88ff78}.

\bibitem{mathlib-triangulated}
The Mathlib Community.
\href{https://github.com/leanprover-community/mathlib4/blob/56a005855ea17e0b0958f73ec801ae140b0a1679/Mathlib/CategoryTheory/Localization/Triangulated.lean}{Triangulated localization},
\texttt{Localization/Triangulated.lean}.
Fixed MathGenome dependency revision \texttt{56a005855ea1}.

\bibitem{stacks-sheaf}
The Stacks Project Authors.
\emph{The Stacks Project}, \href{https://stacks.math.columbia.edu/tag/01DW}{Cohomology of Sheaves}.
Individual results are identified by section or tag.
Archived reference snapshot \texttt{ed88ff78}.

\bibitem{p103}
MathGenome.
\emph{Algebra Core} (P1-03), accepted notes and formalization.
MathGenome v0.21.0, Chapters 10--11.

\bibitem{p202}
MathGenome.
\emph{Algebraic Topology and Riemannian Geometry} (P2-02),
accepted notes and formalization. MathGenome v0.21.0.
\end{thebibliography}
\end{document}
